% Quelques exemples d’endémies et d’épidémies — SVTEEHB 3ème — Cours signé OnBuch+
% Programme officiel MINESEC. Compiler avec : tectonic N11-quelques-exemples-endemies-epidemies.tex
\newcommand{\DOCMATIERE}{SVTEEHB}
\newcommand{\DOCNIVEAU}{3ème}
\newcommand{\DOCMODULE}{SVTEEHB – classe de 3ème}
\newcommand{\DOCLECON}{N11}
\newcommand{\DOCTITRE}{Quelques exemples d’endémies et d’épidémies}
% OnBuch+ — préambule « cours signés » ENRICHI (compilé avec tectonic / XeLaTeX)
% Système visuel « Pop » d'OnBuch+ : fond crème (jamais blanc), bordures encre
% épaisses, ombres « dures » décalées, titres Archivo Black en capitales.
% Version MATHÉMATIQUES : graphes (pgfplots/tikz), boîtes pédagogiques
% (définition, propriété, méthode, attention, le savais-tu, exercice résolu,
% exercice, TP, bilan), page de garde et signature OnBuch+ ancrée en bas.
%
% Macros attendues dans main.tex AVANT \input{preamble} :
%   \DOCMATIERE  \DOCNIVEAU  \DOCMODULE  \DOCLECON (numéro)  \DOCTITRE
\documentclass[11pt]{article}

\usepackage{fontspec}
\usepackage[french]{babel}
\usepackage[a4paper,margin=2cm,top=3.4cm,bottom=2.8cm,headheight=1.4cm,footskip=1.3cm]{geometry}
\usepackage{amsmath,amssymb,amsfonts}
\usepackage{mathtools} % \xmapsto, \coloneqq... souvent utilisés par les modèles
\usepackage{cancel} % simplifications barrées (bilans de forces)
% \abs{z} (module, valeur absolue) et \norm{u} (norme), utilisés par les modèles
\providecommand{\abs}{}\let\abs\relax\DeclarePairedDelimiter{\abs}{\lvert}{\rvert}
\providecommand{\norm}{}\let\norm\relax\DeclarePairedDelimiter{\norm}{\lVert}{\rVert}
\usepackage[table]{xcolor} % [table] : fournit aussi \rowcolors (alternance de lignes)
\usepackage{pagecolor}
\usepackage{tikz}
\usetikzlibrary{arrows.meta,positioning,calc,shapes.geometric,shapes.misc,shapes.arrows,decorations.pathmorphing,decorations.pathreplacing,decorations.markings,patterns,shadows,mindmap,trees,fit,backgrounds,matrix,intersections,angles,quotes}
\usepackage{pgfplots}
\usepackage[version=4]{mhchem} % équations nucléaires \ce{^{235}_{92}U}
\usepackage[european,siunitx]{circuitikz} % schémas électriques
\pgfplotsset{compat=1.18}
\usepgfplotslibrary{fillbetween,groupplots} % aires entre courbes (calcul intégral)
\usepackage{enumitem}
\usepackage{fancyhdr}
% [most] charge toutes les bibliothèques tcolorbox (listings, minted, raster,
% documentation...) et gonfle inutilement la mémoire TeX sur un document long
% (~300 boîtes) : on ne charge que ce qui est réellement utilisé.
\usepackage[skins,breakable]{tcolorbox}
\usepackage{graphicx}
\usepackage{colortbl}
\usepackage{booktabs}
\usepackage{tabularx}
\usepackage{multirow}
\usepackage{diagbox} % case d'en-tête diagonale des tableaux à double entrée
% Colonnes extensibles centrée / gauche / droite (C, L, R), souvent utilisées par les modèles
\newcolumntype{C}{>{\centering\arraybackslash}X}
\newcolumntype{L}{>{\raggedright\arraybackslash}X}
\newcolumntype{R}{>{\raggedleft\arraybackslash}X}
\usepackage{array}
\usepackage{multicol}
\usepackage{siunitx}
% parse-numbers=false : accepte aussi \SI{2,5 \times 10^{-3}}{...}
\sisetup{locale=FR,output-decimal-marker={,},group-minimum-digits=4,parse-numbers=false}
\DeclareSIUnit\ohm{\ensuremath{\symup{\Omega}}} % U+2126 absent de la police
\DeclareSIUnit\division{div} % réglages d'oscilloscope (ms/div, V/div)
\DeclareSIUnit\tr{tr} % tours (vitesse de rotation, ex. \SI{1500}{\tr\per\minute})
\DeclareSIUnit\molar{M} % concentration molaire (ex. \SI{3}{\molar}, \SI{1}{\micro\molar})
\DeclareSIUnit\an{an} % année (le modèle confond parfois avec \year, un compteur TeX) — ex. \SI{5}{\kilo\meter\per\an}
\DeclareSIUnit\FCFA{FCFA} % franc CFA (ex. \SI{500}{\FCFA\per\kilo\gram})
\DeclareSIUnit\degreeBrix{\textdegree Brix} % teneur en sucre d'un sirop/jus (réfractométrie)
\DeclareSIUnit\month{mois} % le modèle utilise parfois \month, un compteur TeX, comme unité de durée
\DeclareSIUnit\microliter{\micro\liter} % le modèle écrit parfois \microliter en un seul token
\DeclareSIUnit\million{millions} % dénombrement (ex. \SI{15}{\million\per\mL} spermatozoïdes)
\usepackage{qrcode}
% \degree : commande fréquemment utilisée par les modèles pour un angle ou une
% température en dehors de \SI{}{} (non fournie par les paquets chargés).
\providecommand{\degree}{^{\circ}}
% \qed (fin de démonstration), utilisé par les modèles sans amsthm
\providecommand{\qed}{\hfill$\square$}
% \barre : double barre des valeurs interdites dans un tableau de variations (tkz-tab)
\providecommand{\barre}{\ensuremath{\|}}
% environnement proof (amsthm non chargé) utilisé par les modèles
\ifdefined\proof\else\newenvironment{proof}[1][Démonstration]{\par\noindent\textit{#1.}\ }{\qed\par}\fi
\usepackage{lastpage}
\usepackage{hyperref}
\hypersetup{hidelinks}

% ============================================================
% Palette « Pop » OnBuch+ — mêmes hex que la classe P (plus_ui.dart)
% ============================================================
\definecolor{poporange}{HTML}{F59321}
\definecolor{popdark}{HTML}{E07A0C}
\definecolor{popcream}{HTML}{FFF6E8}
\definecolor{popink}{HTML}{1C1714}
\definecolor{poppurple}{HTML}{7A5AE0}
\definecolor{popgreen}{HTML}{2E9E6A}
\definecolor{poppink}{HTML}{E0457B}
\definecolor{popblue}{HTML}{2F7DE1}
\definecolor{popgold}{HTML}{C98A12}
\definecolor{popmuted}{HTML}{8A7F76}
\definecolor{popline}{HTML}{EADFD2}
\definecolor{popgray}{HTML}{8A7F76}
\colorlet{popgrey}{popgray}
% Alias tolérants (le modèle nomme parfois les couleurs différemment)
\colorlet{popwhite}{white}
\colorlet{popblack}{popink}
\colorlet{popred}{poppink}
\colorlet{popyellow}{popgold}
% Teintes claires (fonds de tableaux/encadrés) — le modèle nomme parfois
% les couleurs avec un suffixe L (ex. popblueL) sans les avoir définies.
\colorlet{poporangeL}{poporange!20}
\colorlet{popdarkL}{popdark!20}
\colorlet{poppurpleL}{poppurple!20}
\colorlet{popgreenL}{popgreen!20}
\colorlet{poppinkL}{poppink!20}
\colorlet{popblueL}{popblue!20}
\colorlet{popgoldL}{popgold!20}
\colorlet{popredL}{popred!20}
\colorlet{popgrayL}{popgray!20}
% Teintes claires pour fonds de tableaux / zones
\colorlet{popblueL}{popblue!10!white}
\colorlet{poporangeL}{poporange!14!white}
\colorlet{popgreenL}{popgreen!12!white}
\colorlet{poppurpleL}{poppurple!12!white}
\colorlet{poppinkL}{poppink!12!white}
\colorlet{popgoldL}{popgold!14!white}

\pagecolor{popcream}

% ============================================================
% Typographie
% ============================================================
\defaultfontfeatures{Ligatures=TeX}
\setmainfont{PlusJakartaSans-Medium.ttf}[Path=../../../fonts/,BoldFont=PlusJakartaSans-Bold.ttf]
% Maths en sans-serif (Fira Math) avec chiffres et lettres droites de Plus
% Jakarta, pour que les formules se fondent dans le texte.
\usepackage[math-style=ISO,bold-style=ISO]{unicode-math}
\setmathfont{FiraMath-Regular.otf}[Path=../../../fonts/]
\setmathfont{PlusJakartaSans-Medium.ttf}[Path=../../../fonts/,range={up/num,up/{Latin,latin},\mathup}]
\usepackage{icomma} % virgule décimale sans espace en mode math
% Caractères mathématiques Unicode absents des polices de texte (Plus Jakarta,
% Archivo Black) : rendus par leur équivalent mathématique, en texte comme en
% formule — sinon ils s'affichent en carré vide (ex. « Arithmétique dans ℤ »).
\usepackage{newunicodechar}
\newunicodechar{ℝ}{\ensuremath{\mathbb{R}}}
\newunicodechar{ℤ}{\ensuremath{\mathbb{Z}}}
\newunicodechar{ℂ}{\ensuremath{\mathbb{C}}}
\newunicodechar{ℕ}{\ensuremath{\mathbb{N}}}
\newunicodechar{ℚ}{\ensuremath{\mathbb{Q}}}
\newunicodechar{𝔻}{\ensuremath{\mathbb{D}}}
\newunicodechar{α}{\ensuremath{\alpha}}
\newunicodechar{β}{\ensuremath{\beta}}
\newunicodechar{γ}{\ensuremath{\gamma}}
\newunicodechar{δ}{\ensuremath{\delta}}
\newunicodechar{ε}{\ensuremath{\varepsilon}}
\newunicodechar{θ}{\ensuremath{\theta}}
\newunicodechar{λ}{\ensuremath{\lambda}}
\newunicodechar{μ}{\ensuremath{\mu}}
\newunicodechar{σ}{\ensuremath{\sigma}}
\newunicodechar{τ}{\ensuremath{\tau}}
\newunicodechar{φ}{\ensuremath{\varphi}}
\newunicodechar{ψ}{\ensuremath{\psi}}
\newunicodechar{ω}{\ensuremath{\omega}}
\newunicodechar{Σ}{\ensuremath{\Sigma}}
% majuscules produites par \MakeUppercase dans les titres (α-aminés -> Α)
\newunicodechar{Α}{\ensuremath{\alpha}}
\newunicodechar{Β}{\ensuremath{\beta}}
\newunicodechar{Γ}{\ensuremath{\Gamma}}
\newunicodechar{Δ}{\ensuremath{\Delta}}
\newunicodechar{Ε}{\ensuremath{\varepsilon}}
\newunicodechar{Θ}{\ensuremath{\Theta}}
\newunicodechar{Λ}{\ensuremath{\Lambda}}
\newunicodechar{Μ}{\ensuremath{\mu}}
\newunicodechar{Τ}{\ensuremath{\tau}}
\newunicodechar{Φ}{\ensuremath{\Phi}}
\newunicodechar{Ψ}{\ensuremath{\Psi}}
\newunicodechar{Ω}{\ensuremath{\Omega}}
\newunicodechar{↦}{\ensuremath{\mapsto}}
\newunicodechar{∈}{\ensuremath{\in}}
\newunicodechar{∉}{\ensuremath{\notin}}
\newunicodechar{∧}{\ensuremath{\wedge}}
\newunicodechar{⇔}{\ensuremath{\Leftrightarrow}}
\newunicodechar{⟺}{\ensuremath{\Longleftrightarrow}}
\newunicodechar{⇒}{\ensuremath{\Rightarrow}}
\newunicodechar{⟹}{\ensuremath{\Longrightarrow}}
\newunicodechar{∘}{\ensuremath{\circ}}
\newunicodechar{∪}{\ensuremath{\cup}}
\newunicodechar{∩}{\ensuremath{\cap}}
\newunicodechar{≡}{\ensuremath{\equiv}}
\newunicodechar{⊂}{\ensuremath{\subset}}
\newunicodechar{⊥}{\ensuremath{\perp}}
\newunicodechar{∃}{\ensuremath{\exists}}
\newunicodechar{∀}{\ensuremath{\forall}}
\newunicodechar{∛}{\ensuremath{\sqrt[3]{\,}}}
\newunicodechar{∜}{\ensuremath{\sqrt[4]{\,}}}
\newunicodechar{⃗}{\ensuremath{{}^{\to}}}
\newunicodechar{ⁿ}{\ifmmode^{n}\else\textsuperscript{\ensuremath{n}}\fi}
\newunicodechar{ˣ}{\ifmmode^{x}\else\textsuperscript{\ensuremath{x}}\fi}
\newunicodechar{ᵇ}{\ifmmode^{b}\else\textsuperscript{\ensuremath{b}}\fi}
\newunicodechar{ʳ}{\ifmmode^{r}\else\textsuperscript{\ensuremath{r}}\fi}
\newunicodechar{ᵘ}{\ifmmode^{u}\else\textsuperscript{\ensuremath{u}}\fi}
\newunicodechar{ʸ}{\ifmmode^{y}\else\textsuperscript{\ensuremath{y}}\fi}
\newunicodechar{ᵃ}{\ifmmode^{a}\else\textsuperscript{\ensuremath{a}}\fi}
\newunicodechar{ᵉ}{\ifmmode^{e}\else\textsuperscript{\ensuremath{e}}\fi}
\newunicodechar{ᵏ}{\ifmmode^{k}\else\textsuperscript{\ensuremath{k}}\fi}
\newunicodechar{ᵖ}{\ifmmode^{p}\else\textsuperscript{\ensuremath{p}}\fi}
\newunicodechar{ᴺ}{\ifmmode^{N}\else\textsuperscript{\ensuremath{N}}\fi}
\newunicodechar{ᶜ}{\ifmmode^{c}\else\textsuperscript{\ensuremath{c}}\fi}
\newunicodechar{ʰ}{\ifmmode^{h}\else\textsuperscript{\ensuremath{h}}\fi}
\newunicodechar{ᵠ}{\ifmmode^{q}\else\textsuperscript{\ensuremath{q}}\fi}
\newunicodechar{ₙ}{\ifmmode_{n}\else\textsubscript{\ensuremath{n}}\fi}
\newunicodechar{ₐ}{\ifmmode_{a}\else\textsubscript{\ensuremath{a}}\fi}
\newunicodechar{ᵢ}{\ifmmode_{i}\else\textsubscript{\ensuremath{i}}\fi}
\newunicodechar{ᵣ}{\ifmmode_{r}\else\textsubscript{\ensuremath{r}}\fi}
\newunicodechar{ₖ}{\ifmmode_{k}\else\textsubscript{\ensuremath{k}}\fi}
\newunicodechar{ᵦ}{\ifmmode_{\beta}\else\textsubscript{\ensuremath{\beta}}\fi}
\newunicodechar{ₓ}{\ifmmode_{x}\else\textsubscript{\ensuremath{x}}\fi}
\newfontfamily\popdisplay{ArchivoBlack-Regular.ttf}[Path=../../../fonts/]
\newfontfamily\poplogo{SpaceGrotesk-Bold.ttf}[Path=../../../fonts/]
\newfontfamily\popsemi{PlusJakartaSans-SemiBold.ttf}[Path=../../../fonts/]
\newfontfamily\popxbold{PlusJakartaSans-ExtraBold.ttf}[Path=../../../fonts/]
\newfontfamily\popxboldls{PlusJakartaSans-ExtraBold.ttf}[Path=../../../fonts/,LetterSpace=3.0]

\setlist[enumerate]{leftmargin=1.7em,itemsep=5pt,topsep=4pt}
\setlist[itemize]{leftmargin=1.7em,itemsep=4pt,topsep=4pt,label={\color{poporange}\textbullet}}
\setlength{\parindent}{0pt}
\setlength{\parskip}{5pt}
\renewcommand{\arraystretch}{1.25}

% pgfplots : style Pop par défaut
\pgfplotsset{
  every axis/.append style={axis line style={popink,line width=1pt},tick style={popink},
    grid style={popline,line width=0.5pt},label style={font=\small},tick label style={font=\footnotesize}},
  popcurve/.style={poporange,line width=1.8pt,no markers,smooth},
}
% Même style disponible dans un \draw TikZ simple (hors pgfplots)
\tikzset{popcurve/.style={poporange,line width=1.8pt}}
\tikzset{popgrid/.style={popline,very thin}}
\colorlet{popcurve}{poporange} % le nom du style est parfois utilisé comme couleur

% ============================================================
% Pill / squiggle
% ============================================================
\newcommand{\poppill}[3]{%
  \tikz[baseline=-0.55ex]\node[draw=popink,line width=1.2pt,rounded corners=2.6pt,%
    fill=#2,inner xsep=6.5pt,inner ysep=3pt]{%
    \popxboldls\fontsize{7}{7}\selectfont\color{#3}\MakeUppercase{#1}};%
}
\newcommand{\popsquiggle}[1]{%
  \tikz[baseline=-0.4ex]{
    \draw[#1,line width=2.6pt,line cap=round]
      plot [smooth] coordinates {(0,0.09) (0.34,0.01) (0.68,0.21) (1.02,0.01) (1.36,0.10)};
  }%
}
% Mot-clé surligné (marqueur orange sous le texte)
\newcommand{\cle}[1]{{\popxbold\color{popdark}#1}}

% ============================================================
% En-tête / pied
% ============================================================
\pagestyle{fancy}
\fancyhf{}
\renewcommand{\headrulewidth}{0pt}
\renewcommand{\footrulewidth}{0pt}
\lhead{\raisebox{-0.15ex}{\poplogo\fontsize{13}{13}\selectfont\color{popink}OnBuch\color{poporange}+}}
\rhead{\poppill{\DOCMATIERE\ \textbullet\ \DOCNIVEAU\ \textbullet\ Leçon \DOCLECON}{white}{popink}}
\lfoot{\popxbold\fontsize{7.6}{7.6}\selectfont\color{popmuted}\MakeUppercase{Cours signé OnBuch+}\ \textbullet\ {\popsemi\DOCTITRE}}
\rfoot{\tikz[baseline=-0.6ex]\node[rounded corners=8.5pt,fill=popink,minimum height=17pt,minimum width=17pt,inner xsep=5pt,inner sep=0]{\popxbold\fontsize{8}{8}\selectfont\color{popcream}\thepage\,/\,\pageref*{LastPage}};}

% ============================================================
% Titres
% ============================================================
\newcommand{\chaptitre}[2]{%
  \begin{tcolorbox}[enhanced,colback=poporange,colframe=popink,boxrule=2.6pt,arc=11pt,
      drop shadow=popink,left=15pt,right=15pt,top=12pt,bottom=11pt,before skip=3pt,after skip=18pt]
    \poppill{#2}{popink}{popcream}\par\vspace{9pt}
    {\raggedright\hyphenpenalty=10000\exhyphenpenalty=10000\popdisplay\fontsize{22}{24}\selectfont\color{popcream}\MakeUppercase{#1}\par}
  \end{tcolorbox}
}
% Section numérotée : pastille orange avec le numéro + titre Archivo
\newcounter{popsec}
\newcommand{\coursec}[1]{%
  \stepcounter{popsec}\setcounter{popsub}{0}%
  \par\vspace{16pt}\noindent
  \begin{minipage}{\linewidth}
  \tikz[baseline=-0.7ex]\node[draw=popink,line width=1.6pt,fill=poporange,rounded corners=4pt,minimum size=22pt,inner sep=0]{\popdisplay\fontsize{12}{12}\selectfont\color{popcream}\thepopsec};%
  \hspace{8pt}{\popdisplay\fontsize{15}{17}\selectfont\color{popink}\MakeUppercase{#1}}\par
  \vspace{4pt}\noindent{\color{poporange}\rule{36pt}{3.6pt}}
  \end{minipage}\par\nopagebreak\vspace{8pt}
}
\newcounter{popsub}
\newcommand{\courssub}[1]{%
  \stepcounter{popsub}%
  \par\vspace{9pt}\noindent{\popxbold\fontsize{11.5}{13}\selectfont\color{popink}{\color{poporange}\thepopsec.\thepopsub}\hspace{6pt}#1}\par\nopagebreak\vspace{4pt}
}

% ============================================================
% Boîtes pédagogiques — même recette : fond blanc, bordure épaisse colorée,
% ombre dure encre, badge plein. Titre optionnel [#1].
% ============================================================
\tcbset{popbase/.style={breakable,enhanced,colback=white,boxrule=2.3pt,arc=9pt,
  shadow={3pt}{-3pt}{0pt}{popink},left=14pt,right=14pt,top=11pt,bottom=11pt,before skip=11pt,after skip=15pt,
  fontupper=\popsemi\fontsize{10.6}{14.5}\selectfont\color{popink},
  fontlower=\popsemi\fontsize{10.6}{14.5}\selectfont\color{popink}}}
% Coche dessinée (le glyphe ✓ n'existe pas dans les polices du document)
\newcommand{\popcheck}[1]{\tikz[baseline=-0.5ex]\draw[#1,line width=1.6pt,line cap=round,line join=round](-0.13,0.0)--(-0.04,-0.09)--(0.13,0.1);}
\providecommand{\coche}{\popcheck{popgreen}}
\providecommand{\resultat}[1]{\par\smallskip\noindent\fcolorbox{popgreen}{popgreenL}{\parbox{\dimexpr\linewidth-2\fboxsep-2\fboxrule\relax}{\textbf{Résultat.} #1}}\par\smallskip}
\newcommand{\popboxhead}[3]{\poppill{#1}{#2}{white}\ifx\relax#3\relax\else\hspace{6pt}{\popxbold\fontsize{10.6}{12}\selectfont\color{popink}#3}\fi\par\vspace{7pt}}

\newtcolorbox{aretenir}[1][]{popbase,colframe=popgreen,before upper={\popboxhead{À retenir}{popgreen}{#1}}}
\newtcolorbox{exemplebox}[1][]{popbase,colframe=poporange,before upper={\popboxhead{Exemple}{poporange}{#1}}}
\newtcolorbox{definition}[1][]{popbase,colframe=popblue,before upper={\popboxhead{Définition}{popblue}{#1}}}
\newtcolorbox{propriete}[1][]{popbase,colframe=poppurple,before upper={\popboxhead{Propriété}{poppurple}{#1}}}
\newtcolorbox{methode}[1][]{popbase,colframe=popgold,before upper={\popboxhead{Méthode}{popgold}{#1}}}
\newtcolorbox{attention}[1][]{popbase,colframe=poppink,before upper={\popboxhead{Attention}{poppink}{#1}}}
\newtcolorbox{experience}[1][]{popbase,colframe=popblue,colback=popblueL,before upper={\popboxhead{Expérience}{popblue}{#1}}}
% « Le savais-tu ? » : carte violette pleine (contexte, culture, Cameroun)
\newtcolorbox{savaistu}[1][]{popbase,colback=poppurple,colframe=popink,
  fontupper=\popsemi\fontsize{10.4}{14.2}\selectfont\color{white},
  before upper={\poppill{Le savais-tu\,?}{white}{poppurple}\ifx\relax#1\relax\else\hspace{6pt}{\popxbold\color{white}#1}\fi\par\vspace{7pt}}}
% Exercice résolu : énoncé en haut, \tcblower puis la solution sur fond crème
\newcounter{popexo}\newcounter{popexr}
\newtcolorbox{exoresolu}[1][]{popbase,colframe=poporange,colbacklower=popcream,
  segmentation style={popink,line width=1.2pt,dashed},
  before upper={\stepcounter{popexr}\popboxhead{Exercice résolu \thepopexr}{poporange}{#1}},
  before lower={\poppill{Solution}{popink}{popcream}\par\vspace{6pt}}}
% Exercice (non résolu en place) — la correction est dans la partie Corrigés
\newtcolorbox{exercice}[1][]{popbase,colframe=popink,
  before upper={\stepcounter{popexo}\popboxhead{Exercice \thepopexo}{popink}{#1}}}
% Corrigé
\newtcolorbox{corrige}[1][]{popbase,colframe=popgreen,colback=popgreenL,
  before upper={\popboxhead{Corrigé}{popgreen}{#1}}}
% Objectifs / prérequis / bilan
\newtcolorbox{objectifs}{popbase,colframe=popink,colback=poporangeL,
  before upper={\popboxhead{Objectifs de la leçon}{popink}{}}}
\newtcolorbox{prerequis}{popbase,colframe=popink,colback=popblueL,
  before upper={\popboxhead{Prérequis}{popblue}{}}}
\newtcolorbox{bilan}{popbase,colframe=popink,colback=white,boxrule=2.8pt,
  before upper={\popboxhead{Fiche bilan}{poporange}{L'essentiel à connaître pour l'examen}}}
% Figure encadrée avec légende
\newtcolorbox{popfigure}[1][]{enhanced,breakable=false,colback=white,colframe=popline,boxrule=1.4pt,arc=8pt,
  left=10pt,right=10pt,top=10pt,bottom=8pt,before skip=10pt,after skip=12pt,halign=center,
  lower separated=false,halign lower=center,
  fontlower=\popsemi\fontsize{9}{11.5}\selectfont\color{popmuted},
  #1}
\newcommand{\legende}[1]{\par\vspace{6pt}{\popsemi\fontsize{9}{11.5}\selectfont\color{popmuted}#1}}

% ============================================================
% Signature OnBuch+
% ============================================================
\newcommand{\poplogoinline}[1]{{\poplogo\fontsize{#1}{#1}\selectfont\color{popink}OnBuch\color{poporange}+}}
\newcommand{\signatureonbuch}{%
  \begin{tcolorbox}[enhanced,colback=popink,colframe=popink,boxrule=2.2pt,arc=10pt,
      drop shadow=poporange,left=14pt,right=14pt,top=10pt,bottom=10pt,before skip=0pt,after skip=0pt]
    \tikz[baseline=-0.7ex]\node[circle,fill=popgreen,draw=popcream,line width=1.3pt,minimum size=22pt,inner sep=0]{\popcheck{white}};%
    \hspace{9pt}\begin{minipage}[c]{0.62\linewidth}
      {\popxbold\fontsize{12}{13}\selectfont\color{popcream}Signé\ \textbullet\ {\poplogo OnBuch\color{poporange}+}}\par\vspace{2pt}
      {\raggedright\popsemi\fontsize{8}{10.5}\selectfont\color{popline}\MakeUppercase{Équipe pédagogique OnBuch+}\\\MakeUppercase{Conforme au programme officiel MINESEC}\par}
    \end{minipage}\hfill
    {\poplogo\fontsize{22}{22}\selectfont\color{popcream}OnBuch\color{poporange}+}
  \end{tcolorbox}%
}

% ============================================================
% Page de garde
% #1 titre · #2 matière · #3 niveau · #4 module · #5 n° leçon · #6 durée ·
% #7 description / objectifs (texte ou itemize)
% ============================================================
\newcommand{\pagegarde}[7]{%
  \thispagestyle{empty}
  \newgeometry{a4paper,margin=1.7cm,top=1.6cm,bottom=1.4cm}%
  \begin{tikzpicture}[remember picture,overlay]
    \fill[poporange!18] ([xshift=-3.2cm,yshift=-3.6cm]current page.north east) circle (4.6cm);
    \fill[poppurple!14] ([xshift=2.4cm,yshift=7.6cm]current page.south west) circle (3.8cm);
  \end{tikzpicture}%
  {\poplogo\fontsize{30}{30}\selectfont\color{popink}OnBuch\color{poporange}+}\hfill
  \raisebox{0.6ex}{\poppill{Cours signé \textbullet\ Premium}{popink}{popcream}}\par
  \vspace{1pt}\popsquiggle{poppurple}\par
  \vspace{8pt}
  \begin{tcolorbox}[enhanced,colback=poporange,colframe=popink,boxrule=2.8pt,arc=13pt,
      drop shadow=popink,left=18pt,right=18pt,top=12pt,bottom=13pt,after skip=10pt]
    \poppill{#2 \textbullet\ #3}{popink}{popcream}\hspace{5pt}\poppill{#4}{white}{popink}\par\vspace{8pt}
    {\popdisplay\fontsize{14}{14}\selectfont\color{popink}LEÇON #5}\par\vspace{4pt}
    {\raggedright\hyphenpenalty=10000\exhyphenpenalty=10000\popdisplay\fontsize{25}{27}\selectfont\color{popcream}\MakeUppercase{#1}\par}
    \vspace{8pt}
    {\popxbold\fontsize{9.5}{9.5}\selectfont\color{popink}\MakeUppercase{Durée conseillée : #6}}
  \end{tcolorbox}
  \begin{tcolorbox}[enhanced,colback=white,colframe=popink,boxrule=2.2pt,arc=10pt,
      drop shadow=popink,left=16pt,right=16pt,top=10pt,bottom=10pt,after skip=8pt]
    \poppill{Dans cette leçon}{poppurple}{white}\par\vspace{5pt}
    \popsemi\fontsize{9.3}{12.6}\selectfont\color{popink}#7
  \end{tcolorbox}
  \noindent\begin{minipage}[t]{0.487\linewidth}
    \begin{tcolorbox}[enhanced,colback=popgreen,colframe=popink,boxrule=2.2pt,arc=10pt,
        drop shadow=popink,left=11pt,right=11pt,top=10pt,bottom=10pt,height=3.1cm,valign=center]
      \tcbox[colback=white,colframe=popink,boxrule=1.3pt,arc=5pt,boxsep=0pt,nobeforeafter,
        left=3pt,right=3pt,top=3pt,bottom=3pt,valign=center]{\qrcode[height=1.9cm]{https://play.google.com/store/apps/details?id=com.luvvix.onbuch&hl=fr}}%
      \hspace{8pt}\begin{minipage}[c]{0.5\linewidth}
        {\popdisplay\fontsize{11}{12.5}\selectfont\color{white}\MakeUppercase{Télécharge OnBuch}}\par\vspace{4pt}
        {\raggedright\popsemi\fontsize{8.4}{11}\selectfont\color{white}Gratuit \textbullet\ Google Play\\Tuteur IA, annales, concours, résultats\par}
      \end{minipage}
    \end{tcolorbox}
  \end{minipage}\hfill
  \begin{minipage}[t]{0.487\linewidth}
    \begin{tcolorbox}[enhanced,colback=poppurple,colframe=popink,boxrule=2.2pt,arc=10pt,
        drop shadow=popink,left=11pt,right=11pt,top=10pt,bottom=10pt,height=3.1cm,valign=center]
      \tikz[baseline=-0.7ex]\node[circle,fill=white,draw=popink,line width=1.3pt,minimum size=1.6cm,inner sep=0]{\popdisplay\fontsize{24}{24}\selectfont\color{poppurple}+};%
      \hspace{8pt}\begin{minipage}[c]{0.6\linewidth}
        {\popdisplay\fontsize{11}{12.5}\selectfont\color{white}\MakeUppercase{Passe à OnBuch+}}\par\vspace{4pt}
        {\raggedright\popsemi\fontsize{8.4}{11}\selectfont\color{white}Tous les cours signés, Tuteur IA illimité, fiches PDF — onglet OnBuch+ de l'app\par}
      \end{minipage}
    \end{tcolorbox}
  \end{minipage}
  \vspace{8pt}
  \popcontenu
  \vfill
  \signatureonbuch
  \restoregeometry
  \newpage
}

% Rangée « Ce cours contient » (page de garde)
\newcommand{\poptile}[3]{% #1 couleur #2 gros texte #3 légende
  \begin{tcolorbox}[enhanced,colback=#1,colframe=popink,boxrule=2pt,arc=9pt,shadow={2.5pt}{-2.5pt}{0pt}{popink},
      left=6pt,right=6pt,top=9pt,bottom=9pt,halign=center,valign=center,height=1.75cm,width=\linewidth,nobeforeafter]
    {\popdisplay\fontsize{12.5}{13}\selectfont\color{white}\MakeUppercase{#2}}\par\vspace{3pt}
    {\centering\popsemi\fontsize{7.8}{9.5}\selectfont\color{white}#3\par}
  \end{tcolorbox}}
\newcommand{\popcontenu}{%
  \par\noindent{\popxboldls\fontsize{7.5}{7.5}\selectfont\color{popmuted}CE COURS SIGNÉ CONTIENT}\par\vspace{7pt}
  \noindent\begin{minipage}[t]{0.235\linewidth}\poptile{popblue}{Cours}{complet, illustré, pas à pas}\end{minipage}\hfill
  \begin{minipage}[t]{0.235\linewidth}\poptile{poporange}{Méthodes}{et exercices résolus}\end{minipage}\hfill
  \begin{minipage}[t]{0.235\linewidth}\poptile{poppink}{Exercices}{type examen + corrigés}\end{minipage}\hfill
  \begin{minipage}[t]{0.235\linewidth}\poptile{popgreen}{Fiche bilan}{l'essentiel à retenir}\end{minipage}\par}

% Sommaire
\newtcolorbox{sommaire}{enhanced,colback=white,colframe=popink,boxrule=2.2pt,arc=10pt,
  drop shadow=popink,left=18pt,right=18pt,top=13pt,bottom=13pt,before skip=6pt,after skip=20pt,
  before upper={\poppill{Sommaire}{poppurple}{white}\par\vspace{9pt}\popsemi\fontsize{10.6}{14.5}\selectfont\color{popink}}}

% Signature de fin de cours — poussée tout en bas de la dernière page
\newcommand{\findecours}{%
  \par\vfill\nopagebreak
  \begin{center}\popsquiggle{poporange}\end{center}\vspace{4pt}
  \signatureonbuch
}
\AtBeginDocument{\def\llbracket{\mathopen{[\mkern-3.5mu[}}\def\rrbracket{\mathclose{]\mkern-3.5mu]}}} % intervalles d'entiers (glyphe ⟦ absent de la police)
% Glyphes absents de Fira Math : redéfinis à partir de glyphes disponibles
\AtBeginDocument{%
  \def\setminus{\mathbin{\backslash}}\let\smallsetminus\setminus
  \def\vdots{\mathord{\vbox{\baselineskip3.2pt\lineskiplimit0pt\kern4pt\hbox{.}\hbox{.}\hbox{.}}}}%
  \def\ddots{\mathinner{\mkern1mu\raise6.5pt\hbox{.}\mkern2mu\raise3.6pt\hbox{.}\mkern2mu\raise0.7pt\hbox{.}\mkern1mu}}%
  \def\triangle{\mathord{\scalebox{0.9}{$\symup{\Delta}$}}}%
  \def\nabla{\mathord{\raisebox{\depth}{\scalebox{1}[-1]{$\symup{\Delta}$}}}}%
  \def\checkmark{\popcheck{popdark}}%
  \def\star{\mathbin{\ast}}\def\bigstar{\mathord{\ast}}%
  \def\diamond{\mathbin{\scalebox{0.6}{$\Diamond$}}}%
  \def\bigcup{\mathop{\mathchoice{\raisebox{-0.3em}{\scalebox{1.7}{$\cup$}}}{\scalebox{1.25}{$\cup$}}{\cup}{\cup}}\displaylimits}%
  \def\bigcap{\mathop{\mathchoice{\raisebox{-0.3em}{\scalebox{1.7}{$\cap$}}}{\scalebox{1.25}{$\cap$}}{\cap}{\cap}}\displaylimits}%
  \def\lesseqqgtr{\mathrel{\lessgtr}}\def\bigoplus{\mathop{\oplus}\displaylimits}%
}
\providecommand{\ssi}{si et seulement si} % abréviation fréquente
\AtBeginDocument{\providecommand{\wideparen}{\overparen}} % arc de cercle

\begin{document}

\pagegarde{Quelques exemples d’endémies et d’épidémies}{SVTEEHB}{3ème}{SVTEEHB – classe de 3ème}{N11}{4 séances (fiche de progression harmonisée)}{Cette leçon présente la distinction entre endémie et épidémie à travers des exemples concrets : le paludisme, endémie majeure au Cameroun, et les épidémies de fièvre Ebola, de COVID-19 et de choléra. Pour chaque maladie, les causes, les manifestations et les moyens de lutte sont établis à partir de faits expérimentaux et de données épidémiologiques.
\vspace{4pt}
\begin{multicols}{2}\small
\begin{itemize}
\item À la fin de cette leçon, je sais définir les notions d'endémie et d'épidémie et distinguer les deux à partir de données chiffrées.
\item Je sais décrire le cycle de transmission du paludisme et identifier son agent causal (Plasmodium) et son vecteur (l'anophèle femelle).
\item Je sais interpréter des résultats expérimentaux et des courbes épidémiologiques pour établir le mode de transmission d'une maladie.
\item Je sais citer les manifestations du paludisme et les moyens de lutte (préventifs et curatifs).
\item Je sais utiliser correctement une moustiquaire imprégnée et expliquer son intérêt dans la lutte contre le paludisme (TP N°9).
\item Je sais décrire les causes et les moyens de lutte contre la fièvre Ebola, la COVID-19 et le choléra.
\end{itemize}\end{multicols}}

\begin{sommaire}
\begin{multicols}{2}
\begin{enumerate}
\item Endémie et épidémie : définitions et distinction
\item Le paludisme : agent causal et cycle de transmission
\item Le paludisme : manifestations et moyens de lutte
\item TP N°9 : utilisation des moustiquaires imprégnées
\item La fièvre Ebola et la COVID-19 : épidémies virales
\item Le choléra : épidémie hydrique et synthèse
\item Activité d'intégration
\item Méthodes et exercices résolus
\item Exercices
\item Corrigés des exercices
\item Fiche bilan
\end{enumerate}
\end{multicols}
\end{sommaire}

\begin{objectifs}
\begin{itemize}
\item À la fin de cette leçon, je sais définir les notions d'endémie et d'épidémie et distinguer les deux à partir de données chiffrées.
\item Je sais décrire le cycle de transmission du paludisme et identifier son agent causal (Plasmodium) et son vecteur (l'anophèle femelle).
\item Je sais interpréter des résultats expérimentaux et des courbes épidémiologiques pour établir le mode de transmission d'une maladie.
\item Je sais citer les manifestations du paludisme et les moyens de lutte (préventifs et curatifs).
\item Je sais utiliser correctement une moustiquaire imprégnée et expliquer son intérêt dans la lutte contre le paludisme (TP N°9).
\item Je sais décrire les causes et les moyens de lutte contre la fièvre Ebola, la COVID-19 et le choléra.
\item Je sais comparer les modes de transmission (vecteur, contact, eau souillée, gouttelettes) de ces maladies.
\item Je sais appliquer ces notions à des situations concrètes de la vie courante et rédiger une conclusion justifiée.
\end{itemize}
\end{objectifs}

\begin{prerequis}
\begin{itemize}
\item Notion de microbe et de maladie infectieuse (classe de 4ème).
\item Notion d'immunité et de vaccination (leçons précédentes de l'unité).
\item Lecture et interprétation de tableaux de données et de courbes.
\item Hygiène du corps et de l'environnement (vie courante).
\end{itemize}
\end{prerequis}

\newpage

\coursec{Endémie et épidémie : définitions et distinction}

\textbf{Situation d'accroche.} À Ngaoundéré comme à Douala, chaque année, des milliers d'élèves manquent l'école à cause du \cle{paludisme}, une maladie qui ne disparaît jamais complètement de nos régions. En 2014, la fièvre \cle{Ebola} a semé la panique en Afrique de l'Ouest, et en 2020, la \cle{COVID-19} a fermé les écoles du monde entier, y compris au Cameroun. Deux situations bien différentes : d'un côté, une maladie toujours présente ; de l'autre, des maladies qui surgissent brutalement, se propagent vite, puis finissent par s'éteindre.

\textbf{Problématique.} Pourquoi certaines maladies restent-elles installées durablement dans une région alors que d'autres apparaissent brutalement puis disparaissent ? Comment l'Homme peut-il se protéger et lutter contre elles ? Pour répondre à ces questions, commençons par apprendre à \emph{classer} les maladies selon leur mode d'apparition dans une population : c'est la distinction entre \cle{endémie} et \cle{épidémie}.

\courssub{Qu'est-ce qu'une endémie ?}

\textbf{Observation.} Interrogeons le registre d'un centre de santé de Yaoundé : chaque mois, année après année, on y soigne des malades du paludisme. Il n'y a pas de « début » ni de « fin » : la maladie est toujours là, avec un nombre de cas qui reste élevé mais relativement stable au fil des ans. Une telle maladie est dite \emph{endémique}.

\begin{definition}[Endémie]
Une \cle{endémie} est une maladie qui sévit de façon \textbf{permanente} dans une région donnée, avec un nombre de cas relativement \textbf{stable} au cours du temps.
\end{definition}

\begin{exemplebox}[Le paludisme, endémie au Cameroun]
Le \cle{paludisme} est une endémie au Cameroun : il est présent toute l'année, dans toutes les régions du pays. Le Cameroun enregistre \textbf{plusieurs millions de cas} de paludisme chaque année. Le nombre de cas peut augmenter pendant la saison des pluies (les moustiques pullulent dans les eaux stagnantes), mais la maladie ne disparaît jamais complètement.
\end{exemplebox}

\begin{attention}[Une endémie n'est pas une « petite » maladie]
Erreur fréquente : croire qu'une endémie est moins grave qu'une épidémie parce qu'elle fait moins parler d'elle. C'est faux ! Une épidémie n'est pas forcément plus mortelle qu'une endémie : le paludisme tue chaque année, dans le monde, \textbf{plus de personnes} que la plupart des épidémies. Le mot « endémie » décrit la \emph{durée} et la \emph{stabilité} de la maladie, pas sa gravité.
\end{attention}

\courssub{Qu'est-ce qu'une épidémie ?}

\textbf{Observation.} En 2021--2022, plusieurs régions du Cameroun, dont le Littoral (Douala), ont connu une flambée de \cle{choléra} : en quelques semaines, le nombre de malades est passé de presque zéro à plusieurs centaines, puis l'épidémie a été maîtrisée et les cas ont chuté. Contrairement au paludisme, le choléra n'est pas présent en permanence : il \emph{surgit}, se \emph{propage rapidement}, puis \emph{disparaît}.

\begin{definition}[Épidémie]
Une \cle{épidémie} est l'apparition \textbf{brutale} et la propagation \textbf{rapide} d'une maladie contagieuse touchant un \textbf{grand nombre de personnes} d'une population, sur une période limitée.
\end{definition}

\begin{exemplebox}[Quelques épidémies célèbres]
\begin{itemize}
\item La \cle{fièvre Ebola} (Afrique de l'Ouest, 2014--2016) : maladie virale très contagieuse par contact direct avec les liquides biologiques (sang, vomissures, selles) des malades.
\item La \cle{COVID-19} (à partir de 2019) : maladie virale respiratoire qui a obligé à fermer les écoles, y compris au Cameroun en 2020.
\item Le \cle{choléra} : maladie bactérienne transmise par l'eau et les aliments souillés ; le Cameroun a connu plusieurs épidémies, notamment en 2021--2022 (plusieurs milliers de cas).
\end{itemize}
\end{exemplebox}

\begin{savaistu}[D'où viennent les mots « épidémie » et « endémie » ?]
Le mot \emph{épidémie} vient du grec ancien : \emph{epi} signifie « sur » et \emph{demos} signifie « peuple ». Une épidémie, c'est littéralement ce qui « tombe sur le peuple » ! De même, \emph{endémie} vient de \emph{en} (« dans ») et \emph{demos} : la maladie qui reste « dans le peuple ».
\end{savaistu}

\textbf{Et la pandémie ?} Lorsqu'une épidémie ne se limite plus à une région ou à un pays, mais s'étend à \textbf{plusieurs continents}, on parle de \cle{pandémie}. C'est le cas de la COVID-19 en 2020, qui a touché presque tous les pays du monde en quelques mois.

\begin{aretenir}[Les trois mots à ne pas confondre]
\begin{itemize}
\item \cle{Endémie} : maladie \textbf{toujours présente} dans une région (ex. : paludisme au Cameroun).
\item \cle{Épidémie} : maladie qui \textbf{apparaît brutalement} et se propage vite dans une population, sur une période limitée (ex. : choléra à Douala).
\item \cle{Pandémie} : épidémie qui s'étend à \textbf{plusieurs continents} (ex. : COVID-19 en 2020).
\end{itemize}
\end{aretenir}

\courssub{Analyse de données : comparer une endémie et une épidémie}

Pour bien visualiser la différence, comparons deux séries de données (valeurs simplifiées mais réalistes) : le nombre de cas de \textbf{paludisme} au Cameroun sur 5 ans, et le nombre de cas de \textbf{choléra} au Cameroun sur la même période, avec l'épidémie de 2021--2022.

\begin{center}
\begin{tabularx}{\linewidth}{|l|X|X|X|X|X|}
\hline
\rowcolor{popblueL} \textbf{Année} & \textbf{2018} & \textbf{2019} & \textbf{2020} & \textbf{2021} & \textbf{2022} \\
\hline
Paludisme (millions de cas) & 2,5 & 2,6 & 2,4 & 2,6 & 2,5 \\
\hline
Choléra (milliers de cas) & 0,1 & 0,2 & 0,1 & 2,5 & 15 \\
\hline
\end{tabularx}
\end{center}

\textbf{Interprétation.} La courbe du paludisme est presque \emph{horizontale} : le nombre de cas reste élevé mais stable, année après année — c'est la signature d'une \cle{endémie}. La courbe du choléra, elle, reste proche de zéro puis forme un \emph{pic brutal} en 2021--2022 avant de redescendre — c'est la signature d'une \cle{épidémie}.

\begin{popfigure}
\begin{center}
\begin{tikzpicture}
\begin{axis}[
width=0.95\linewidth, height=6.5cm,
xlabel={Année}, ylabel={Nombre de cas},
xmin=2018, xmax=2022, ymin=0, ymax=16,
xtick={2018,2019,2020,2021,2022},
xticklabels={2018,2019,2020,2021,2022},
legend style={at={(0.02,0.97)}, anchor=north west, font=\small},
grid=major, grid style={popline!40},
]
\addplot[popcurve, color=popblue, very thick, mark=*] coordinates {
(2018,10) (2019,10.4) (2020,9.6) (2021,10.4) (2022,10)
};
\addlegendentry{Paludisme (unité : 250 000 cas) — endémie}
\addplot[popcurve, color=poporange, very thick, mark=square*] coordinates {
(2018,0.1) (2019,0.2) (2020,0.1) (2021,2.5) (2022,15)
};
\addlegendentry{Choléra (unité : 1 000 cas) — épidémie}
\end{axis}
\end{tikzpicture}
\end{center}
\legende{Comparaison de l'évolution du nombre de cas de paludisme (courbe bleue, stable et élevée : endémie) et de choléra (courbe orange, pic brutal en 2021--2022 : épidémie) au Cameroun sur 5 ans. Pour permettre la comparaison sur le même graphique, les deux courbes utilisent des unités différentes indiquées dans la légende.}
\end{popfigure}

\begin{methode}[Comment distinguer une endémie d'une épidémie à partir de données ?]
\begin{enumerate}
\item Observer l'\textbf{évolution du nombre de cas} au cours du temps (tableau ou courbe).
\item Si le nombre de cas reste \textbf{élevé et stable} pendant des années : c'est une \textbf{endémie}.
\item Si le nombre de cas \textbf{augmente brutalement} (pic) puis diminue jusqu'à presque disparaître : c'est une \textbf{épidémie}.
\item Si l'épidémie s'étend à plusieurs continents : c'est une \textbf{pandémie}.
\end{enumerate}
\end{methode}

\courssub{Tableau récapitulatif : endémie, épidémie, pandémie}

\begin{center}
\begin{tabularx}{\linewidth}{|l|X|X|X|}
\hline
\rowcolor{popblueL} \textbf{Critère} & \textbf{Endémie} & \textbf{Épidémie} & \textbf{Pandémie} \\
\hline
Présence dans le temps & Permanente, stable & Brutale, limitée dans le temps & Brutale, limitée dans le temps \\
\hline
Étendue géographique & Une région donnée & Une population, une région ou un pays & Plusieurs continents \\
\hline
Allure de la courbe des cas & Courbe stable (presque horizontale) & Pic brutal puis disparition & Pics successifs dans le monde entier \\
\hline
Exemples & Paludisme au Cameroun & Choléra (2021--2022), fièvre Ebola (2014--2016) & COVID-19 (2020) \\
\hline
\end{tabularx}
\end{center}

\courssub{Le vocabulaire de la transmission des maladies}

Pour décrire comment une maladie se transmet, les scientifiques utilisent trois mots précis qu'il faut bien connaître :

\begin{definition}[Agent causal, vecteur, hôte]
\begin{itemize}
\item L'\cle{agent causal} est le \textbf{microbe responsable} de la maladie (un micro-organisme : bactérie, virus ou parasite). Exemple : l'agent causal du choléra est une bactérie.
\item Le \cle{vecteur} est l'\textbf{organisme qui transmet} l'agent causal d'un hôte à un autre, sans être lui-même malade. Exemple : le moustique femelle (l'anophèle) est le vecteur du paludisme.
\item L'\cle{hôte} est l'\textbf{organisme infecté} par l'agent causal, chez qui la maladie se développe. Exemple : l'Homme est l'hôte du parasite du paludisme.
\end{itemize}
\end{definition}

\begin{attention}[Ne pas confondre vecteur et agent causal]
Erreur très fréquente au BEPC : dire que « le moustique est le microbe du paludisme ». C'est faux ! Le moustique est le \textbf{vecteur} : il transporte le parasite (l'agent causal) et l'inocule à l'Homme en le piquant. Le moustique lui-même ne souffre pas du paludisme.
\end{attention}

\begin{exoresolu}[Endémie ou épidémie ?]
\textbf{Énoncé.} Voici deux affirmations entendues à la radio :
\begin{enumerate}
\item « Chaque année, la région du Nord enregistre environ le même nombre de cas de paludisme. »
\item « En quelques semaines, plusieurs centaines de cas de choléra sont apparus à Douala, puis l'épidémie a été maîtrisée. »
\end{enumerate}
Pour chaque affirmation, indique s'il s'agit d'une endémie ou d'une épidémie, en justifiant ta réponse.
\tcblower
\textbf{Solution détaillée.}
\begin{enumerate}
\item Il s'agit d'une \textbf{endémie}. Justification : le paludisme est présent \emph{en permanence} dans la région du Nord, avec un nombre de cas \emph{stable} d'une année à l'autre. Or, une endémie est une maladie qui sévit de façon permanente dans une région donnée.
\item Il s'agit d'une \textbf{épidémie}. Justification : les cas de choléra sont apparus \emph{brutalement} et en \emph{grand nombre} sur une période \emph{limitée}, avant de disparaître. Or, une épidémie est l'apparition brutale et la propagation rapide d'une maladie contagieuse dans une population.
\end{enumerate}
\textbf{Conseil de rédaction :} au BEPC, cite toujours la définition (ou l'idée de la définition) pour justifier ta réponse : c'est ce qui rapporte les points.
\end{exoresolu}

\begin{exercice}[À toi de jouer]
En 2020, la COVID-19, apparue d'abord en Chine, s'est propagée en quelques mois à presque tous les pays du monde.
\begin{enumerate}
\item S'agit-il d'une endémie, d'une épidémie ou d'une pandémie ? Justifie.
\item Cite une maladie endémique au Cameroun et précise son vecteur.
\end{enumerate}
\end{exercice}

\begin{corrige}[Exercice : À toi de jouer]
\begin{enumerate}
\item Il s'agit d'une \textbf{pandémie} : la COVID-19 est apparue brutalement et s'est propagée rapidement (caractère épidémique), mais elle s'est étendue à \textbf{plusieurs continents}, ce qui définit une pandémie.
\item Le \textbf{paludisme} est une maladie endémique au Cameroun. Son vecteur est le \textbf{moustique femelle} (l'anophèle), qui transmet le parasite responsable de la maladie en piquant l'Homme.
\end{enumerate}
\end{corrige}

\begin{aretenir}[L'essentiel de la section]
\begin{itemize}
\item Une \cle{endémie} sévit en permanence dans une région, avec un nombre de cas stable (ex. : paludisme au Cameroun, plusieurs millions de cas par an).
\item Une \cle{épidémie} apparaît brutalement, se propage vite, puis s'éteint (ex. : choléra, Ebola, COVID-19).
\item Une \cle{pandémie} est une épidémie étendue à plusieurs continents (ex. : COVID-19 en 2020).
\item Une endémie peut tuer plus qu'une épidémie : ces mots décrivent le mode d'apparition, pas la gravité.
\item Trois acteurs de la transmission : l'\cle{agent causal} (le microbe), le \cle{vecteur} (qui le transporte) et l'\cle{hôte} (qui est infecté).
\end{itemize}
\end{aretenir}

\coursec{Endémie et épidémie : définitions et distinction}

\paragraph{Situation déclenchante.}
Au cours de l'année 2023, le district de santé de \textbf{Bertoua} (Région de l'Est) signale une augmentation brutale des cas de \textbf{fièvre hémorragique à virus Ebola} : 12 cas confirmés en deux semaines, alors qu'aucun cas n'avait été notifié depuis des années. Parallèlement, le même district enregistre habituellement \textbf{300 à 400 cas de paludisme par semaine} tout au long de l'année, avec un pic en saison des pluies. Pourquoi parle-t-on d'\textbf{épidémie} pour Ebola et d'\textbf{endémie} pour le paludisme ?

\begin{definition}[Endémie]
Une \textbf{endémie} est la présence \textbf{constante} et \textbf{habituelle} d'une maladie infectieuse dans une population donnée, dans une zone géographique délimitée, avec une incidence relativement stable au fil du temps (malgré des variations saisonnières). La maladie y est \textbf{en permanence}.
\end{definition}

\begin{definition}[Épidémie]
Une \textbf{épidémie} est la survenue \textbf{brutale} et \textbf{inhabituelle} d'un nombre de cas d'une maladie (souvent infectieuse) \textbf{nettement supérieur} au nombre de cas attendus (seuil d'alerte) dans une population et une zone données, sur une période définie. Si l'épidémie s'étend à plusieurs pays ou continents, on parle de \textbf{pandémie}.
\end{definition}

\begin{popfigure}
\begin{tikzpicture}[scale=0.9, every node/.style={font=\small}]
    \begin{axis}[
        width=\linewidth, height=6cm,
        xlabel={Temps (mois)}, ylabel={Nombre de cas},
        xmin=0, xmax=12, ymin=0, ymax=500,
        xtick={1,2,3,4,5,6,7,8,9,10,11,12},
        xticklabels={J,F,M,A,M,J,J,A,S,O,N,D},
        ytick={0,100,200,300,400,500},
        grid=major, grid style={popline},
        legend pos=north west,
        legend style={fill=popcream, draw=popdark},
        title style={font=\bfseries\small}, title={Évolution mensuelle simulée : Paludisme (Endémie) vs Choléra (Épidémie)}
    ]
    % Paludisme : endémie saisonnière
    \addplot[popblue, thick, mark=*, mark size=2pt] coordinates {
        (1,250) (2,220) (3,280) (4,350) (5,420) (6,450) (7,400) (8,380) (9,410) (10,380) (11,300) (12,260)
    };
    \addlegendentry{Paludisme (Endémie)}
    % Seuil d'épidémie (exemple arbitraire pour le choléra)
    \addplot[popred, dashed, thick, domain=0:12] {100};
    \addlegendentry{Seuil d'alerte épidémique (ex: Choléra)}
    % Choléra : épidémie brutale
    \addplot[poporange, very thick, mark=square*, mark size=2pt] coordinates {
        (1,5) (2,3) (3,2) (4,10) (5,15) (6,20) (7,180) (8,350) (9,280) (10,120) (11,40) (12,15)
    };
    \addlegendentry{Choléra (Épidémie)}
    \end{axis}
\end{tikzpicture}
\legende{Comparaison schématique : le paludisme montre une présence permanente avec variation saisonnière (endémie) ; le choléra reste au niveau de bruit de fond puis explose brutalement au-delà du seuil d'alerte (épidémie).}
\end{popfigure}

\begin{aretenir}[Critères de distinction Endémie vs Épidémie]
\begin{center}
\begin{tabularx}{\linewidth}{|l|X|X|}\hline
\textbf{Critère} & \textbf{Endémie (ex: Paludisme)} & \textbf{Épidémie (ex: Ebola, Choléra, COVID-19)} \\ \hline
\textbf{Présence} & Constante, permanente & Brutale, temporaire, inhabituelle \\ \hline
\textbf{Nombre de cas} & Stable (variation saisonnière) & Explosion > Seuil d'attente/alerte \\ \hline
\textbf{Durée} & Années, décennies (si pas de contrôle) & Semaines à mois \\ \hline
\textbf{Immunité collective} & Souvent élevée (prémunition) & Faible ou nulle (nouvel agent ou réintroduction) \\ \hline
\textbf{Exemples au Cameroun} & Paludisme, Onchocercose, Bilharziose & Choléra, Fièvre Ebola, COVID-19, Rougeole, Méningite \\ \hline
\end{tabularx}
\end{center}
\end{aretenir}

\coursec{Le paludisme : agent causal et cycle de transmission}

Dans la section précédente, nous avons distingué l'endémie (présence constante d'une maladie dans une population donnée) de l'épidémie (survenue brutale d'un nombre de cas supérieur à l'attendu). Le paludisme, encore appelé \textbf{malaria}, constitue l'exemple type d'\textbf{endémie} au Cameroun et dans toute l'Afrique subsaharienne. Selon l'OMS (Rapport 2023, données 2022), la région africaine supportait environ \textbf{94\,\%} des cas mondiaux (\textbf{249 millions}) et \textbf{95\,\%} des décès (\textbf{608\,000}), dont une grande majorité d'enfants de moins de 5 ans. Au Cameroun, le paludisme reste la première cause de consultation et d'hospitalisation dans nos formations sanitaires. Comprendre son agent causal et son mode de transmission est la clé pour mettre en œuvre les moyens de lutte efficaces que nous verrons par la suite.

\courssub{L'agent causal : le Plasmodium, un parasite unicellulaire}

\paragraph{Observation historique (Laveran, 1880).}
C'est en 1880, à l'hôpital militaire de Constantine (Algérie), que le médecin militaire français \textbf{Charles Louis Alphonse Laveran} observe pour la première fois, au microscope optique, des éléments mobiles et pigmentés dans le sang frais d'un soldat fiévreux. Il note la présence de « petits corpuscules mobiles » à l'intérieur des globules rouges, qu'il identifie comme un parasite : \emph{Oscillaria malariae}, rebaptisé plus tard \textbf{Plasmodium}.

\begin{savaistu}[Alphonse Laveran (1845--1922)]
Alphonse Laveran, né à Paris, médecin militaire, reçoit le \textbf{prix Nobel de physiologie ou médecine en 1907} pour sa découverte du parasite du paludisme. Il a consacré la somme du prix à la création d'un laboratoire de médecine tropicale à l'Institut Pasteur de Paris. Au Cameroun, l'Institut Pasteur de Yaoundé et le Centre Pasteur du Cameroun à Garoua perpétuent cette tradition de recherche sur le paludisme.
\end{savaistu}

\begin{definition}[Agent causal du paludisme]
Le paludisme est une maladie parasitaire due à un protozoaire (eucaryote unicellulaire) du genre \textbf{\emph{Plasmodium}}. Cinq espèces infectent l'homme : \emph{Plasmodium falciparum} (la plus fréquente et la plus grave en Afrique), \emph{P. vivax}, \emph{P. ovale}, \emph{P. malariae} et \emph{P. knowlesi} (zoonose d'Asie du Sud-Est).
\end{definition}

\begin{attention}[Espèces et gravité]
Ne confondez pas les espèces. \emph{Plasmodium falciparum} est responsable des formes graves (paludisme cérébral, anémie sévère, détresse respiratoire) car il fait adhérer les globules rouges infectés à l'endothélium vasculaire (séquestration), obstruant la microcirculation. Les autres espèces provoquent généralement des formes bénignes, mais \emph{P. vivax} et \emph{P. ovale} peuvent persister sous forme dormante (\textbf{hypnozoïtes}) dans le foie et causer des rechutes mois ou années plus tard.
\end{attention}

\paragraph{Morphologie microscopique (savoir admis).}
Sur une \textbf{goutte épaisse} (technique de référence pour le diagnostic de routine au laboratoire, concentration du parasite) ou un \textbf{frottis sanguin} (pour l'identification spécifique et le comptage de la parasitémie), on observe les stades érythrocytaires :
\begin{itemize}
    \item \textbf{Formes annulaires (jeunes trophozoïtes)} : anneau de cytoplasme avec un noyau central (aspect « bague »).
    \item \textbf{Trophozoïtes matures} : cytoplasme plus abondant, pigment (hémozoïne) brun-noir issu de la digestion de l'hémoglobine.
    \item \textbf{Schizontes} : stade de multiplication asexuée (mérogonie) donnant 8 à 32 mérozoïtes (rarement vus dans le sang périphérique pour \emph{P. falciparum} car séquestrés dans la microcirculation profonde).
    \item \textbf{Gamétocytes} : formes sexuées (mâle = microgamétocyte, femelle = macrogamétocyte), seul stade infectieux pour le moustique. Chez \emph{P. falciparum}, ils sont en « croissant » ou « banane ».
\end{itemize}

\begin{popfigure}
\begin{tikzpicture}[scale=0.9, every node/.style={font=\small}]
    % Erythrocyte background
    \draw[popred!30, line width=1.5pt] (0,0) ellipse (2.5cm and 1.8cm);
    \draw[popred!50, line width=0.5pt] (0,0) ellipse (2.3cm and 1.6cm);
    \node[popred!80] at (0, -2.2) {Globule rouge (hématie)};
    % Ring stage (young trophozoite)
    \begin{scope}[shift={(-1.5, 0.5)}]
        \draw[popblue, thick] (0,0) circle (0.35);
        \draw[popblue, thick] (0,0) circle (0.15);
        \fill[popblue!20] (0,0) circle (0.35);
        \fill[white] (0,0) circle (0.15);
        \node[popblue] at (0, -0.7) {1. Forme annulaire};
    \end{scope}
    % Mature trophozoite
    \begin{scope}[shift={(1.5, 0.5)}]
        \draw[popblue, thick] (0,0) circle (0.45);
        \fill[popblue!30] (0,0) circle (0.45);
        \fill[poporange!70!popdark] (-0.1,0.1) circle (0.08) (0.1,-0.1) circle (0.08) (-0.15,-0.15) circle (0.06);
        \node[popblue] at (0, -0.7) {2. Trophozoïte mature};
        \node[poporange!70!popdark, font=\tiny] at (0.5, 0.4) {Pigment (hémozoïne)};
    \end{scope}
    % Schizont (rare in peripheral blood for Pf)
    \begin{scope}[shift={(-1.5, -1.2)}]
        \draw[poppurple, thick] (0,0) circle (0.5);
        \fill[poppurple!20] (0,0) circle (0.5);
        \foreach \angle in {0,45,90,135,180,225,270,315} {
            \fill[poppurple] (\angle:0.35) circle (0.07);
        }
        \node[poppurple] at (0, -0.8) {3. Schizonte (mérogonie)};
    \end{scope}
    % Gametocyte (falciparum = crescent)
    \begin{scope}[shift={(1.5, -1.2)}]
        \draw[popgreen!70!black, thick] (-0.5,0) .. controls (-0.5,0.3) and (0.5,0.3) .. (0.5,0) .. controls (0.5,-0.3) and (-0.5,-0.3) .. cycle;
        \fill[popgreen!30!black] (-0.5,0) .. controls (-0.5,0.3) and (0.5,0.3) .. (0.5,0) .. controls (0.5,-0.3) and (-0.5,-0.3) .. cycle;
        \fill[poporange!70!popdark] (0,0) circle (0.06);
        \node[popgreen!70!black] at (0, -0.8) {4. Gamétocyte (falciparum)};
    \end{scope}
\end{tikzpicture}
\legende{Stades érythrocytaires du \emph{Plasmodium falciparum} observables au microscope optique (coloration Giemsa, grossissement $\times$1000). Seuls les gamétocytes (4) sont infectieux pour le moustique vecteur.}
\end{popfigure}

\courssub{Le vecteur : l'anophèle femelle}

\paragraph{Observation.}
Le paludisme ne se transmet pas par contact direct (toux, toucher, objets souillés) entre un malade et une personne saine. La transmission nécessite un intermédiaire biologique obligatoire : un moustique du genre \textbf{\emph{Anopheles}}.

\begin{definition}[Vecteur (en parasitologie)]
Un \textbf{vecteur} est un organisme vivant (souvent un arthropode hématophage : moustique, tique, mouche, puce) qui transmet un agent pathogène (virus, bactérie, parasite) d'un hôte vertébré infecté (réservoir) à un hôte vertébré sain, \textbf{sans être lui-même malade} de cette pathologie. Le parasite effectue une partie de son cycle de développement (multiplication, maturation) à l'intérieur du vecteur.
\end{definition}

\begin{attention}[Seule la femelle pique !]
Chez les moustiques (Culicidés), \textbf{seule la femelle est hématophage} (elle pique pour prélever du sang). Elle a besoin des protéines du sang pour la maturation de ses œufs (vitellogenèse). Le mâle se nourrit exclusivement de nectar et de sucs végétaux. C'est donc \textbf{exclusivement la femelle anophèle} qui peut transmettre le Plasmodium.
\end{attention}

\begin{popfigure}
\begin{tikzpicture}[scale=1.1, every node/.style={font=\small}]
    % Body parts
    % Abdomen (engorged with blood)
    \fill[popred!40] (0,0) ellipse (1.2cm and 0.6cm);
    \draw[popdark, thick] (0,0) ellipse (1.2cm and 0.6cm);
    \node[popdark] at (0, -1.0) {Abdomen (gorgé de sang)};
    % Thorax
    \fill[popgray!30] (-1.3, 0.3) ellipse (0.5cm and 0.4cm);
    \draw[popdark, thick] (-1.3, 0.3) ellipse (0.5cm and 0.4cm);
    \node[popdark] at (-1.3, 1.0) {Thorax};
    % Head with proboscis
    \fill[popgray!30] (-2.1, 0.6) circle (0.3cm);
    \draw[popdark, thick] (-2.1, 0.6) circle (0.3cm);
    % Proboscis
    \draw[popdark, very thick, ->] (-2.4, 0.6) -- (-3.2, 0.6) node[right, pos=0.5] {Proboscis (rostre piqueur)};
    % Wings (resting position: angled)
    \draw[popblue!30, thick, fill=popblue!10] (-1.0, 0.5) .. controls (-0.5, 1.2) and (0.5, 1.0) .. (0.2, 0.3) .. controls (0.0, 0.0) .. cycle;
    \node[popblue] at (-0.5, 1.3) {Ailes tachetées};
    % Legs (characteristic resting angle)
    \draw[popdark, thick] (-1.5, -0.1) -- (-2.2, -0.8);
    \draw[popdark, thick] (-0.8, -0.2) -- (-0.5, -1.0);
    \draw[popdark, thick] (0.2, -0.2) -- (0.8, -0.9);
    \draw[popdark, thick] (-1.8, 0.1) -- (-2.5, 0.6);
    \draw[popdark, thick] (-0.5, 0.2) -- (0.0, 0.8);
    \node[popdark] at (1.0, -0.5) {Pattes (position caractéristique} ;
    \node[popdark] at (1.0, -0.8) {corps incliné, abdomen relevé)};
    % Palps (long in Anopheles female)
    \draw[popdark, thick] (-2.3, 0.4) -- (-2.8, 0.2) node[left] {Palpes maxillaires (longs $\approx$ proboscis)};
    % Antennae
    \draw[popdark, thick] (-2.1, 0.9) -- (-2.3, 1.3);
    \draw[popdark, thick] (-1.9, 0.9) -- (-1.7, 1.3);
    % Scale bar
    \draw[popdark, |-|] (2.5, -1.2) -- (3.5, -1.2) node[midway, below] {1 mm};
\end{tikzpicture}
\legende{Morphologie de la femelle \emph{Anopheles} en position de repos (corps incliné, abdomen relevé) et de piqûre. Critères de distinction avec \emph{Culex} (moustique commun) : palpes maxillaires longs (presque aussi longs que le proboscis), ailes tachetées, position de repos à angle aigu sur la surface. Seule la femelle (représentée ici) pique.}
\end{popfigure}

\courssub{Établissement expérimental du rôle du moustique (Travaux de Ronald Ross, 1897--1898)}

\paragraph{Problématique.}
Après la découverte du parasite dans le sang par Laveran, la question restait ouverte : \emph{comment le parasite passe-t-il d'un homme à un autre ?} Patrick Manson émet l'hypothèse d'une transmission par un moustique. Ronald Ross, médecin britannique en Inde, va le vérifier expérimentalement.

\begin{experience}[Démonstration du rôle vecteur de l'anophèle par Ronald Ross (1897-1898)]
\textbf{Matériel :} Oiseaux infectés par \emph{Plasmodium relictum} (paludisme aviaire), moustiques \emph{Culex} et \emph{Anopheles} élevés en laboratoire (non infectés), microscopes.
\textbf{Protocole :}
\begin{enumerate}
    \item Laisser des moustiques femelles (\emph{Culex} puis \emph{Anopheles}) piquer des oiseaux parasités.
    \item Disséquer les moustiques à intervalles réguliers (1, 2, 3... jours après le repas de sang) et observer le tube digestif (estomac/midgut) au microscope.
    \item Contrôle : disséquer des moustiques élevés en laboratoire n'ayant \textbf{jamais} piqué d'oiseau malade.
    \item Laisser des moustiques ayant piqué des oiseaux malades (et gardés vivants plusieurs jours) piquer des oiseaux sains (non infectés).
\end{enumerate}
\textbf{Observations :}
\begin{itemize}
    \item Dans l'estomac des moustiques ayant piqué des oiseaux malades : apparition de \textbf{kystes (oocystes)} sur la paroi externe de l'intestin moyen, 2 à 3 jours après le repas. Ces kystes grossissent et libèrent des milliers de \textbf{sporozoïtes} qui migrent vers les glandes salivaires.
    \item \textbf{Aucun kyste} chez les moustiques témoins (non nourris sur sang infecté).
    \item Les oiseaux sains piqués par des moustiques « infectés » (porteurs de sporozoïtes aux glandes salivaires) développent le paludisme (parasites dans le sang). Les oiseaux piqués par des moustiques témoins restent sains.
\end{itemize}
\textbf{Interprétation :} Le moustique n'est pas un simple porteur mécanique (piqueur sale). Le Plasmodium \textbf{se développe et se multiplie} dans le tube digestif du moustique (cycle sexué + sporogonie) avant de gagner les glandes salivaires. La transmission se fait par injection de salive contenant des sporozoïtes lors de la piqûre.
\textbf{Conclusion :} \textbf{La piqûre de la femelle anophèle infectée est le mode de transmission exclusif du paludisme.} Ross reçoit le prix Nobel en 1902.
\end{experience}

\begin{propriete}[Cycle biologique du Plasmodium : transmission vectorielle obligatoire]
Le cycle du \emph{Plasmodium} comprend deux phases obligatoires :
\begin{enumerate}
    \item \textbf{Cycle asexué (schizogonie) chez l'Homme (hôte intermédiaire) :} multiplication dans les hépatocytes (foie) puis dans les érythrocytes (sang). Responsable des symptômes.
    \item \textbf{Cycle sexué (gamétogonie) + sporogonie chez l'Anophèle (hôte définitif / vecteur) :} fécondation dans le tube digestif, formation d'oocystes, production de sporozoïtes migrateurs vers les glandes salivaires.
\end{enumerate}
\emph{Il n'y a pas de transmission directe homme-homme (sauf exceptions rares : transfusion sanguine, transmission congénitale mère-enfant, accident d'exposition au sang).}
\end{propriete}

\courssub{Cycle simplifié de transmission Homme $\leftrightarrow$ Anophèle}

Le cycle ci-dessous synthétise les étapes clés. La compréhension de ce cycle est indispensable pour cibler les moyens de lutte (moustiquaires, insecticides, traitement des malades, vaccin).

\begin{popfigure}
\begin{tikzpicture}[scale=0.85, every node/.style={font=\small}, node distance=1.2cm and 1.5cm]
    % Styles
    \tikzset{
        box/.style={draw=popdark, thick, rounded corners, fill=popcream, minimum width=2.8cm, minimum height=1.2cm, align=center},
        arrow/.style={->, >=Stealth, thick, popdark},
        labelarrow/.style={midway, fill=white, inner sep=2pt, font=\footnotesize},
        stage/.style={ellipse, draw=popblue, thick, fill=popblueL, minimum width=2.2cm, minimum height=1cm},
        parasite/.style={rectangle, draw=popred, thick, fill=popred!10, minimum width=2.5cm, minimum height=0.9cm, align=center},
    }
    % HUMAN HOST (Left)
    \node[box, align=center] (human) at (-5, 0){\textbf{HÔTE HUMAIN}\\(Hôte intermédiaire)};
    \node[stage, above left=0.5cm and 0.5cm of human] (liver) {Foie (Hépatocytes)};
    \node[stage, below left=0.5cm and 0.5cm of human] (blood) {Sang (Érythrocytes)};
    % MOSQUITO VECTOR (Right)
    \node[box, align=center] (mosquito) at (5, 0){\textbf{VECTEUR ANOPHÈLE}\\(Hôte définitif / Vecteur)};
    \node[stage, above right=0.5cm and 0.5cm of mosquito] (gut) {Tube digestif (Estomac)};
    \node[stage, below right=0.5cm and 0.5cm of mosquito] (salivary) {Glandes salivaires};
    % PARASITE STAGES (Center flow)
    \node[parasite, align=center] (sporo) at (-1.5, 1.8){SPOROZOÏTES\\(injectés par la salive)};
    \node[parasite, align=center] (liverstage) at (-1.5, 0.5){SCHIZONTES HÉPATIQUES\\(Mérogonie $\to$ Mérozoïtes)};
    \node[parasite, align=center] (bloodstage) at (-1.5, -0.8){CYCLE ÉRYTHROCYTAIRE\\(Anneaux $\to$ Trophozoïtes $\to$ Schizontes $\to$ Mérozoïtes)};
    \node[parasite, align=center] (gametocyte) at (-1.5, -2.1){GAMÉTOCYTES\\(Formes sexuées : $\sigma$ \& $\varphi$)};
    \node[parasite, align=center] (gutstage) at (1.5, 0.5){FÉCONDATION $\to$ ZYGOTE (Ookinète)\\OOCYSTE (Paroi intestinale)};
    \node[parasite, align=center] (sporogony) at (1.5, -0.8){SPOROGONIE\\(Méiose $\to$ Sporozoïtes)};
    \node[parasite, align=center] (salivarystage) at (1.5, -2.1){SPOROZOÏTES\\(Migration vers glandes salivaires)};
    % ARROWS: Human cycle
    \draw[arrow] (sporo.south) -- (liverstage.north) node[labelarrow, left] {1. Piqûre infectante};
    \draw[arrow] (liverstage.south) -- (bloodstage.north) node[labelarrow, left] {2. Rupture hépatocytes};
    \draw[arrow] (bloodstage.south) -- (gametocyte.north) node[labelarrow, left] {3. Différenciation sexuelle};
    % ARROW: Human to Mosquito (Ingestion)
    \draw[arrow, poporange, very thick] (gametocyte.east) -- node[labelarrow, above, align=center]{4. Repas de sang\\ (Ingestion gamétocytes)} (gutstage.west);
    % ARROWS: Mosquito cycle
    \draw[arrow] (gutstage.south) -- (sporogony.north) node[labelarrow, right] {5. Fécondation \& Oocyste};
    \draw[arrow] (sporogony.south) -- (salivarystage.north) node[labelarrow, right] {6. Sporogonie \& Migration};
    % ARROW: Mosquito to Human (Injection)
    \draw[arrow, popgreen!60!black, very thick] (salivarystage.west) -- node[labelarrow, below, align=center]{7. Piqûre\\ (Injection sporozoïtes)} (sporo.east);
    % Legend box
    \node[draw=popdark, fill=popcream, rounded corners, align=left, font=\footnotesize] at (0, -3.5) {
        \textbf{Légende du cycle :}\\
        \textcolor{poporange}{\textbf{\textbullet\ Flèches orange}} = Passage Homme $\to$ Moustique (Ingestion).\\
        \textcolor{popgreen!60!black}{\textbf{\textbullet\ Flèches vertes}} = Passage Moustique $\to$ Homme (Injection).\\
        \textbf{Durée du cycle chez le moustique (sporogonie) :} 10 à 15 jours (dépend de la température).\\
        \textbf{Durée du cycle hépatique :} 5 à 16 jours (selon espèce).\\
        \textbf{Durée du cycle érythrocytaire :} 48h (\emph{P. falciparum, vivax, ovale}) ou 72h (\emph{P. malariae}).
    };
\end{tikzpicture}
\legende{Cycle de transmission du \emph{Plasmodium falciparum} entre l'Homme et l'Anophèle femelle. Le parasite alterne reproduction asexuée (multiplication) dans l'homme et reproduction sexuée + sporogonie dans le moustique. L'interruption du cycle est possible à plusieurs niveaux (voir section lutte).}
\end{popfigure}

\paragraph{Détail des étapes physiologiques.}

\begin{enumerate}
    \item \textbf{Inoculation (Homme $\leftarrow$ Moustique) :} Lors de la piqûre, la femelle anophèle injecte sa salive (anticoagulants, vasodilatateurs) contenant des \textbf{sporozoïtes} (forme allongée, mobile, $\approx$ \SI{10}{\micro\meter}--\SI{15}{\micro\meter}). Quelques dizaines à centaines de sporozoïtes suffisent.
    \item \textbf{Phase hépatique (exo-érythrocytaire) :} Les sporozoïtes atteignent le foie via la circulation sanguine en quelques minutes. Ils envahissent les hépatocytes. Chaque sporozoïte donne un \textbf{schizonte hépatique} qui, par mérogonie (multiplication asexuée), produit \textbf{10\,000 à 30\,000 mérozoïtes} (chez \emph{P. falciparum}). \textbf{Cette phase est asymptomatique.} (Note : pour \emph{P. vivax/ovale}, certains sporozoïtes deviennent des \textbf{hypnozoïtes} dormants $\to$ rechutes).
    \item \textbf{Phase érythrocytaire (symptomatique) :} Les mérozoïtes libérés envahissent les globules rouges (GR). Cycle de \SI{48}{\hour} : Anneau $\to$ Trophozoïte (digestion hémoglobine $\to$ hémozoïne) $\to$ Schizonte $\to$ Éclatement (hémolyse) libérant 16--32 mérozoïtes + toxines + débris cellulaires $\to$ Fièvre paroxysmique. Une infime fraction se différencie en \textbf{gamétocytes} (stade sexuel, non pathogène, infectieux pour le moustique).
    \item \textbf{Ingestion par le moustique (Homme $\to$ Moustique) :} Lors d'un repas de sang sur un porteur de gamétocytes matures, le moustique ingère les gamétocytes.
    \item \textbf{Cycle dans le moustique (Sporogonie) :}
    \begin{itemize}
        \item Dans le lumen de l'estomac (midgut) : \textbf{Gamétogenèse}. Le microgamétocyte (mâle) effectue 3 divisions mitotiques rapides $\to$ 8 microgamètes flagellés (exflagellation). Le macrogamétocyte (femelle) mûrit $\to$ 1 macrogamète.
        \item \textbf{Fécondation} : Un microgamète féconde le macrogamète $\to$ \textbf{Zygote (diploïde)} $\to$ \textbf{Ookinète} (mobile, traverse la paroi intestinale).
        \item L'ookinète se ronde sous la membrane basale $\to$ \textbf{Oocyste}. Division méiotique (réductionnelle) puis mitotiques multiples $\to$ des milliers de \textbf{sporozoïtes} (haploïdes).
        \item \textbf{Rupture de l'oocyste} (\SIrange{10}{15}{\day} post-infection, température-dépendant) $\to$ sporozoïtes dans l'hémocèle $\to$ migration active vers les \textbf{glandes salivaires}. Le moustique devient \textbf{infectant} pour toute sa vie (quelques semaines).
    \end{itemize}
\end{enumerate}

\begin{methode}[Identifier le mode de transmission d'une maladie (Démarche d'investigation)]
Pour déterminer si une maladie est vectorielle, hydrique, aérienne ou de contact, on suit une démarche comparative :
\begin{enumerate}
    \item \textbf{Observation épidémiologique :} Décrire la distribution spatio-temporelle des cas (groupement autour d'un point d'eau ? saison des pluies ? contact familial ?).
    \item \textbf{Formulation d'hypothèses :} « La transmission se fait par l'eau », « par un moustique », « par l'air »...
    \item \textbf{Expérimentation / Enquête cas-témoins :} Comparer un groupe \textbf{exposé} au facteur suspect (ex: dormir sans moustiquaire, boire l'eau du puits X) et un groupe \textbf{non exposé} (moustiquaire imprégnée, eau traitée). Mesurer l'incidence de la maladie dans les deux groupes.
    \item \textbf{Recherche de l'agent :} Identifier l'agent pathogène dans le vecteur suspect (dissection moustique, analyse eau, prélèvement air).
    \item \textbf{Reproduction expérimentale :} Infecter l'hôte sain via le facteur suspect uniquement (expérience de Ross).
    \item \textbf{Conclusion :} Valider ou invalider le mode de transmission.
\end{enumerate}
\end{methode}

\courssub{Exercice résolu : Analyse du cycle et conséquences sur la lutte}

\begin{exoresolu}[Cycle du Plasmodium et cibles thérapeutiques]
\textbf{Énoncé :} Un chercheur teste trois molécules candidates (A, B, C) sur cultures de \emph{P. falciparum}.
\begin{itemize}
    \item Molécule A : bloque l'invasion des hépatocytes par les sporozoïtes.
    \item Molécule B : empêche la maturation des gamétocytes dans le sang.
    \item Molécule C : inhibe la division méiotique dans l'oocyste du moustique.
\end{itemize}
\textbf{Questions :}
\begin{enumerate}
    \item Classer ces molécules selon leur cible : prophylaxie (prévention infection), traitement curatif (arrêt symptômes), ou blocage de la transmission (réduction réservoir).
    \item Laquelle serait la plus utile pour un vaccin « pré-érythrocytaire » type RTS,S/AS01 (Mosquirix\textsuperscript{\textregistered}) ?
    \item Pourquoi le traitement des porteurs sains de gamétocytes (asymptomatiques) est-il crucial pour l'élimination ?
\end{enumerate}

\tcblower
\textbf{Solution détaillée :}
\begin{enumerate}
    \item \textbf{Molécule A (blocage invasion hépatocytes)} : Cible la phase \textbf{pré-érythrocytaire}. Empêche l'installation de l'infection $\to$ \textbf{Prophylaxie / Vaccin pré-érythrocytaire}. Si prise avant exposition, empêche la maladie. Si prise après infection hépatique établie, inefficace sur le cycle sanguin.
    \item \textbf{Molécule B (blocage maturation gamétocytes)} : Cible le stade \textbf{infectieux pour le moustique}. N'agit pas sur les schizontes érythrocytaires (donc pas d'effet direct sur la fièvre du patient \emph{au moment de la prise}), mais empêche le patient de contaminer les moustiques $\to$ \textbf{Blocage de la transmission (transmission-blocking)}. Exemple réel : Primaquine (seule molécule active sur gamétocytes matures de \emph{P. falciparum} et hypnozoïtes de \emph{P. vivax}).
    \item \textbf{Molécule C (blocage méiose dans l'oocyste)} : Cible le développement \textbf{dans le vecteur}. Empêche la formation de sporozoïtes $\to$ \textbf{Blocage de la transmission (côté vecteur)}. Stratégie : « Transmission-blocking vaccine » (TBV) ou modification génétique des moustiques (gene drive).
    \item \textbf{Vaccin RTS,S/AS01 (Mosquirix\textsuperscript{\textregistered})} : Cible la \textbf{protéine CSP (Circumsporozoite Protein)} à la surface du sporozoïte. Il vise à induire des anticorps neutralisant le sporozoïte \textbf{avant qu'il n'envahisse l'hépatocyte}. Cible = \textbf{Molécule A}. Efficacité partielle ($\approx$ 30--40\,\% sur 4 ans chez l'enfant), ne protège pas contre l'infection hépatique établie ni le cycle sanguin.
    \item \textbf{Portage asymptomatique :} Dans les zones d'endémie stable (comme le Cameroun), une grande partie de la population (enfants, adultes semi-immuns) porte des parasites à faible densité \textbf{sans fièvre}. Ces porteurs abritent souvent des \textbf{gamétocytes}. Ils constituent le \textbf{réservoir principal} qui maintient la transmission endémique. Les traiter (dépistage et traitement de masse, ou traitement préventif intermittent) réduit la source d'infection pour les moustiques, brisant le cycle Homme $\to$ Moustique.
\end{enumerate}
\end{exoresolu}

\begin{aretenir}[Synthèse : Agent causal et transmission du paludisme]
\begin{itemize}
    \item \textbf{Agent :} \emph{Plasmodium} spp. (protozoaire), principalement \emph{P. falciparum} en Afrique. Découvert par Laveran (1880, Nobel 1907).
    \item \textbf{Vecteur :} Femelle \emph{Anopheles} (hématophage pour maturation ovaires). Rôle vecteur prouvé par Ross (1897-98, Nobel 1902) : cycle sexué + sporogonie dans le moustique.
    \item \textbf{Transmission :} \textbf{Vectorielle obligatoire} (piqûre). Pas de transmission directe interhumaine (sauf sang/congénital).
    \item \textbf{Cycle :} Sporozoïtes (salive moustique) $\to$ Foie (schizogonie asymptomatique) $\to$ Sang (cycle érythrocytaire 48h = symptômes) $\to$ Gamétocytes $\to$ Estomac moustique (fécondation, oocyste, sporogonie) $\to$ Glandes salivaires (sporozoïtes).
    \item \textbf{Point clé :} L'interruption du cycle est possible au niveau du vecteur (moustiquaires, insecticides), du parasite dans le sang (traitement, vaccin), ou du contact vecteur-homme.
\end{itemize}
\end{aretenir}

\coursec{Le paludisme : manifestations et moyens de lutte}

\courssub{Manifestations cliniques}

\paragraph{Observation clinique.}
Un enfant de 3 ans, résident à Douala, est amené à la consultation pour \textbf{fièvre élevée (39,5\,°C)}, frissons, céphalées, vomissements et abattement survenant par \textbf{accès répétés toutes les 48 heures}. L'examen retrouve une \textbf{splénomégalie} (rate palpable à 3 cm du rebord costal) et une pâleur des conjonctives. Le test de diagnostic rapide (TDR) au sang périphérique est positif pour \emph{P. falciparum} (ligne HRP2).

\begin{definition}[Accès palustre typique (P. falciparum)]
L'accès palustre classique comprend trois stades successifs :
\begin{enumerate}
    \item \textbf{Stade des frissons} (15--60 min) : sensation de froid intense, chair de poule, contraction musculaire, élévation brutale de la température.
    \item \textbf{Stade de la fièvre} (2--6 h) : température \SIrange{39}{41}{\celsius}, peau chaude et sèche, tachycardie, céphalées intenses, délire possible chez l'enfant.
    \item \textbf{Stade de la déférence} : transpiration abondante, chute brutale de la température, asthénie marquée, sommeil.
\end{enumerate}
La périodicité est de \textbf{48 h} (tertiane) pour \emph{P. falciparum, vivax, ovale} et \textbf{72 h} (quarte) pour \emph{P. malariae}.
\end{definition}

\begin{attention}[Paludisme grave : urgence vitale]
Chez l'enfant < 5 ans et la femme enceinte, \emph{P. falciparum} peut évoluer vers des formes graves en quelques heures :
\begin{itemize}
    \item \textbf{Paludisme cérébral :} troubles de la conscience (score de Blantyre < 3), convulsions, coma.
    \item \textbf{Anémie sévère :} Hb < \SI{5}{\gram\per\deci\liter} (ou Ht < 15\,\%), pâleur extrême, essoufflement au repos.
    \item \textbf{Détresse respiratoire / Acidose :} respiration de Kussmaul (ample et profonde).
    \item \textbf{Hypoglycémie :} glycémie < \SI{2.2}{\milli\mole\per\liter} (\SI{40}{\milli\gram\per\deci\liter}), fréquente sous quinine/artésunate.
    \item \textbf{Insuffisance rénale aiguë :} oligurie, hémoglobinurie (urines « Coca-Cola »).
\end{itemize}
\textbf{Tout signe de gravité impose une référence immédiate vers un hôpital pour traitement parentéral (artésunate IV/IM).}
\end{attention}

\paragraph{Paludisme et grossesse.}
La femme enceinte (surtout primigeste) est plus susceptible : le parasite se séquestre dans le placenta (adhésion à l'acide chondroïtine sulfate A). Conséquences : \textbf{anémie maternelle sévère}, \textbf{insuffisance pondérale à la naissance} (IPN), \textbf{avortement}, \textbf{mortinatalité}. La \textbf{Chimioprévention du Paludisme pendant la Grossesse (CPG)} par Sulfadoxine-Pyriméthamine (SP) aux 2ème et 3ème trimestres (doses espacées d'au moins 1 mois) est la stratégie nationale.

\courssub{Moyens de lutte : stratégie intégrée}

La lutte anti-vectorielle (LAV) et la prise en charge des cas (PEC) sont les deux piliers. Au Cameroun, le \textbf{Programme National de Lutte contre le Paludisme (PNLP)} coordonne ces interventions.

\begin{propriete}[Moyens de lutte classés par cible du cycle]
\begin{center}
\begin{tabularx}{\linewidth}{|l|X|X|}\hline
\textbf{Cible du cycle} & \textbf{Moyen de lutte} & \textbf{Principe / Mise en œuvre au Cameroun} \\ \hline
\textbf{Moustique vecteur} (Adulte) & \textbf{Moustiquaire imprégnée d'insecticide à longue durée d'action (MILDA)} & Barrière physique + effet chimique (pyréthrinoïdes : perméthrine, deltaméthrine, ou nouvelles générations PBO, chlorfénapyr). Distribution gratuite massive (campagnes tous les 3 ans) + distribution continue (CPN, vaccination). \textbf{Objectif :} couverture universelle (1 MILDA / 2 personnes). \\ \hline
& \textbf{Pulvérisation intra-domiciliaire à effet rémanent (PID/IRS)} & Pulvérisation d'insecticides (pyréthrinoïdes, organophosphates, néonicotinoïdes) sur les murs intérieurs. Ciblée dans zones à forte transmission/résistance (ex: Nord, Extrême-Nord). \\ \hline
\textbf{Moustique vecteur} (Larve) & \textbf{Lutte anti-larvaire (LAL)} & Destruction des gîtes larvaires (assainissement, drainage, comblement), larvicides biologiques (\emph{Bacillus thuringiensis} var. \emph{israelensis} - Bti) ou chimiques. Adaptée aux zones urbaines (Yaoundé, Douala) où les gîtes sont identifiables. \\ \hline
\textbf{Parasite chez l'Homme} & \textbf{Diagnostic précoce \& Traitement efficace (TPE)} & TDR (HRP2/pLDH) ou Microscopie (référence). \textbf{Traitement de 1ère ligne :} \textbf{Artéméthrine-Luméfantrine (AL)} (ACT) 6 prises sur 3 jours. \textbf{Alternative :} Artesunate-Amodiaquine (ASAQ). Traitement gratuit pour < 5 ans et femmes enceintes dans le public. \\ \hline
& \textbf{Chimioprévention du Paludisme Saisonnier (CPS)} & Administration mensuelle de SP + Amodiaquine (SPAQ) pendant 4 mois (saison des pluies) aux enfants \textbf{3--59 mois} dans le Sahel (Régions Nord, Adamaoua, Extrême-Nord). Réduit la mortalité de $\approx$ 75\,\%. \\ \hline
& \textbf{Chimioprévention Grossesse (CPG)} & SP aux contacts CPN (dose 1 au 2ème trimestre, dose 2 au 3ème trimestre, min 1 mois d'intervalle). \\ \hline
& \textbf{Vaccin RTS,S/AS01 (Mosquirix\textsuperscript{\textregistered})} & Introduit au Cameroun en \textbf{janvier 2024} (Pilotage : Régions Centre, Sud, Est, Littoral). 4 doses : 6, 7, 9 mois + rappel 24 mois. Cible sporozoïtes (protéine CSP). Efficacité $\approx$ 30--40\,\% sur paludisme grave. Complémentaire, ne remplace pas MILDA. \\ \hline
\end{tabularx}
\end{center}
\end{propriete}

\begin{attention}[Résistance aux insecticides et aux antipaludéens]
\begin{itemize}
    \item \textbf{Résistance aux pyréthrinoïdes :} Fréquente chez \emph{An. gambiae} s.l. et \emph{An. funestus} au Cameroun (mutation \emph{kdr}, métabolique). Justifie l'usage de MILDA de nouvelle génération (PBO, chlorfénapyr, pyriproxyfène) et la rotation d'insecticides pour la PID.
    \item \textbf{Résistance à l'artémisinine :} Mutations gène \emph{kelch13} (C580Y etc.) émergentes en Afrique de l'Est (Rwanda, Ouganda, Érythrée). Surveillance moléculaire active au Cameroun (Réseau ReSMAC). Pas de résistance clinique confirmée à ce jour pour l'AL ou l'ASAQ, mais vigilance accrue (suivi de la clairance parasitaire J3).
    \item \textbf{Délétion du gène \emph{pfhrp2/3} :} Rend les TDR basés sur HRP2 faussement négatifs. Surveillance nécessaire ; utilisation de TDR pan-pLDH ou microscopie en cas de doute.
\end{itemize}
\end{attention}

\begin{exemplebox}[Exemple de calcul : Couverture en MILDA]
\textbf{Situation :} Un ménage de 7 personnes (2 adultes, 5 enfants) reçoit 3 MILDA lors d'une campagne de distribution de masse.
\textbf{Question :} Ce ménage atteint-il l'objectif de couverture universelle (1 MILDA / 2 personnes) ?
\textbf{Raisonnement :} Objectif = 7 personnes / 2 = 3,5 $\to$ \textbf{4 MILDA nécessaires}. Le ménage a 3 MILDA $\to$ \textbf{Couverture insuffisante} (déficit de 1 MILDA). Deux personnes dormiront sans protection.
\textbf{Action :} Compléter via les canaux de distribution continue (CPN, vaccination) ou redistribution ciblée.
\end{exemplebox}

\coursec{TP N°9 : Utilisation des moustiquaires imprégnées}

\begin{experience}[TP N°9 : Observation, manipulation et évaluation de l'état d'une MILDA]
\textbf{Objectifs (Savoir-faire) :}
\begin{enumerate}
    \item Identifier les différents types de moustiquaires (non imprégnée, imprégnée conventionnelle, MILDA).
    \item Manipuler correctement une MILDA (déballage, suspension, rentrée, lavage, séchage, réparation).
    \item Évaluer l'état physique d'une MILDA (trous, propreté) selon la grille OMS (proportionate Hole Index - pHI).
    \item Sensibiliser un pair sur les bonnes pratiques d'utilisation.
\end{enumerate}

\textbf{Matériel :}
\begin{itemize}
    \item Échantillons de moustiquaires : MILDA neuve (ex: PermaNet 2.0/3.0, Olyset Net, Royal Guard), moustiquaire usagée (propre, sale, déchirée), moustiquaire non imprégnée (filet simple).
    \item Kit de réparation : aiguille, fil, morceaux de filet (patchs) ou ruban adhésif médical.
    \item Bassin, eau, savon doux (type savon de Marseille), corde, clous/crochets pour suspension.
    \item Grille d'évaluation pHI simplifiée (OMS 2013) : Taille 1 (< 2 cm), Taille 2 (2--10 cm), Taille 3 (> 10 cm). pHI = (Nb T1 $\times$ 1) + (Nb T2 $\times$ 23) + (Nb T3 $\times$ 196). Bon état si pHI $\le$ 64 ; Acceptable si 64 < pHI $\le$ 642 ; À remplacer si pHI > 642.
\end{itemize}

\textbf{Protocole :}
\begin{enumerate}
    \item \textbf{Identification :} Observer l'étiquette (marque, type d'insecticide, année de production, norme WHOPES/PQ). Différencier MILDA (insecticide incorporé dans la fibre, durée 3 ans / 20 lavages) et moustiquaire conventionnelle (imprégnation de surface, 6 mois / 3 lavages).
    \item \textbf{Installation (Simulation) :} Suspendre la MILDA au-dessus d'un lit/matelas simulé. Rentrer soigneusement les bords sous le matelas ou le tapis. Vérifier qu'aucun espace ne permet le passage des moustiques.
    \item \textbf{Évaluation physique (pHI) :} Étaler la moustiquaire usagée à plat. Compter les trous par catégorie de taille (à l'aide d'un gabarit ou règle). Calculer le pHI. Décider : « Bon état », « À réparer », « À remplacer ».
    \item \textbf{Réparation :} Coudre les trous de taille 1 et 2. Poser un patch (morceau de filet cousu) ou du ruban adhésif pour les trous de taille 3 (si réparable). Noter : une réparation ne restaure pas l'insecticide.
    \item \textbf{Lavage (Démonstration) :} Laver une MILDA neuve ou peu usagée : eau froide/tiède (< \SI{30}{\celsius}), savon doux, \textbf{pas de détergent, pas de javel, pas de brosse}. Frotter délicatement à la main. Rincer abondamment. \textbf{Sécher à l'ombre} (pas au soleil direct : UV dégradent l'insecticide et la fibre). Ne pas tordre violemment.
    \item \textbf{Message clé à restituer :} « Chaque nuit, toute l'année, tout le monde dort sous MILDA bien tendue, bien rentrée, lavée correctement, réparée si trouée, remplacée au bout de 3 ans. »
\end{enumerate}

\textbf{Exploitation / Compte-rendu attendu :}
\begin{itemize}
    \item Fiche descriptive des échantillons (Type, État pHI, Décision).
    \item Schéma annoté : Installation correcte (bords rentrés) vs incorrecte (bords lâchés, moustique entrant).
    \item Rédaction d'un message de sensibilisation court (radio/affichage) en français ou langue locale : « La MILDA protège ma famille, je l'utilise bien. »
\end{itemize}

\end{experience}

\begin{aretenir}[Bonnes pratiques MILDA (Règle des 3 « L ») : \textbf{L}aver (savon doux, ombre), \textbf{L}ouer (réparer trous), \textbf{L}ancer (remplacer après 3 ans / 20 lavages).]
\end{aretenir}

\coursec{La fièvre Ebola et la COVID-19 : épidémies virales}

\courssub{La fièvre hémorragique à virus Ebola (FHVE)}

\paragraph{Agent causal et réservoir.}
Maladie virale due à un \textbf{Filovirus} (famille \emph{Filoviridae}), genre \emph{Ebolavirus}. 6 espèces identifiées : \emph{Zaire ebolavirus} (EBOV, la plus létale, épidémies RDC, Guinée 2014, 2021), \emph{Sudan ebolavirus} (SUDV, Ouganda, Soudan), \emph{Bundibugyo ebolavirus} (BDBV), \emph{Taï Forest ebolavirus} (TAFV, Côte d'Ivoire), \emph{Reston ebolavirus} (RESTV, non pathogène pour l'homme), \emph{Bombali ebolavirus} (BOMV, chauves-souris). \textbf{Réservoir naturel :} Chauves-souris frugivores (ex: \emph{Hypsignathus monstrosus}, \emph{Epomops franqueti}, \emph{Myonycteris torquata}). Transmission primaire : contact avec fluides corporels d'animaux infectés (chasse, manipulation viande de brousse : singes, antilopes, chauves-souris).

\paragraph{Transmission interhumaine (Épidémie).}
\textbf{Contact direct} avec le sang, sécrétions, organes ou liquides biologiques (selles, urines, salive, sperme, lait maternel, larmes) de personnes \textbf{malades ou décédées}. \textbf{Pas de transmission aérienne} (aérosols) ni vectorielle (moustique). Risque maximal : phase terminale, funérailles (lavage/embrassement du corps), soignants sans EPI (Équipements de Protection Individuelle), contacts sexuels (persistance du virus dans le sperme > 12 mois).

\paragraph{Manifestations et diagnostic.}
Période d'incubation : \textbf{2 à 21 jours} (moyenne 8--10 j). Début brutal : fièvre, fatigue intense, myalgies, céphalées, pharyngite. Suivi de : vomissements, diarrhées, éruption cutanée, hémorragies (gencives, injection, digestives), insuffisance multi-viscérale. Létalité : \textbf{25--90\,\%} selon souche et prise en charge.
\textbf{Diagnostic :} RT-PCR (sang, écouvillon) en laboratoire de niveau de sécurité P3/P4 (Centre Pasteur Yaoundé, INRB Kinshasa). TDR antigéniques (ReEBOV) pour triage rapide.

\paragraph{Moyens de lutte (Stratégie « Anneau » / Ring Vaccination).}
\begin{enumerate}
    \item \textbf{Surveillance / Détection précoce :} Définition de cas (OMS), alerte immédiate, isolement du cas suspect (CST - Centre de Soins de Traitement).
    \item \textbf{Prise en charge :} Réanimation symptomatique (réhydratation IV/ORale aggressive, correction électrolytes, transfusion, vasopresseurs). Antiviraux : \textbf{Ansuvimab (Mab114)}, \textbf{Regn-EB3 (Inmazeb\textsuperscript{\textregistered})} (anticorps monoclonaux anti-EBOV, efficacité prouvée PALM trial). Antiviraux à large spectre (Remdesivir, Favipiravir) efficacité moindre.
    \item \textbf{Vaccination (Prophylaxie post-exposition / Anneau) :} \textbf{rVSV-ZEBOV (Ervebo\textsuperscript{\textregistered})} : vaccin vivant atténué recombinant (VSV vecteur + glycoprotéine EBOV). Efficacité $\approx$ 100\,\% (essai Guinée 2015). Administré aux contacts (cas contacts, contacts des contacts) et personnel de santé (1 dose IM). Nouveau vaccin 2 doses (Ad26.ZEBOV + MVA-BN-Filo - Zabdeno/Mvabea) pour prophylaxie préventive.
    \item \textbf{Prévention/Contrôle infection (PCI) :} EPI complet (combinaison étanche, double gants, masque FFP2/FFP3, lunettes/visière, bottes) pour soignants/enterreurs. Enterrements sécurisés et dignes (équipes formées). Désinfection (chlore 0,05\,\% / 0,5\,\%).
    \item \textbf{Communication / Engagement communautaire :} Lutte contre la rumeur, stigmatisation, résistance aux équipes de riposte.
\end{enumerate}

\courssub{La maladie à Coronavirus 2019 (COVID-19)}

\paragraph{Agent causal et émergence.}
\textbf{SARS-CoV-2} (Coronavirus 2 du Syndrome Respiratoire Aigu Sévère). Famille \emph{Coronaviridae}, genre \emph{Betacoronavirus}. Origine zoonotique probable (chauve-souris $\to$ hôte intermédiaire ? $\to$ homme). Émergence : Wuhan (Chine) déc. 2019. Pandémie déclarée OMS mars 2020. Au Cameroun : 1er cas mars 2020.

\paragraph{Transmission.}
\textbf{Aérienne (gouttelettes > 5 µm et aérosols < 5 µm)} : toux, éternuement, parole, chant, respiration. \textbf{Contact direct/indirect} : mains contaminées $\to$ muqueuses (yeux, nez, bouche). \textbf{Contagiosité :} Pic 1--2 jours avant symptômes $\to$ 8--10 jours après. Variants (Alpha, Delta, Omicron BA.1/BA.2/BA.5/XBB/JN.1...) : transmissibilité accrue, échappement immunitaire partiel.

\paragraph{Manifestations.}
Asymptomatique ($\approx$ 30--40\,\%) à forme critique. Symptômes fréquents : fièvre, toux sèche, fatigue, anosmie/agueusie (perte odorat/goût - caractéristique variants initiaux), céphalées, myalgies, gorge irritée. Formes graves : pneumopathie bilatérale (verre dépoli scanner), SDRA (détresse respiratoire), tempête cytokinique, thrombose, atteinte multi-organes. Facteurs de risque : âge > 60 ans, obésité, diabète, HTA, immunodépression, grossesse.
\textbf{COVID long :} symptômes persistants > 4 semaines (fatigue, dyspnée, troubles cognitifs « brouillard cérébral »).

\paragraph{Moyens de lutte.}
\begin{enumerate}
    \item \textbf{Mesures barrières (Non pharmaceutiques) :} Port du masque (chirurgical/FFP2) en espace clos/affluence, hygiène des mains (gel hydroalcoolique/lavage), distanciation physique (\SI{1}{\meter}), aération des lieux clos (10 min/h), isolement des cas (7--10 j), recherche des contacts.
    \item \textbf{Diagnostic :} RT-PCR (référence, \emph{gène E, RdRp, N}). TDR Ag (antigène) rapide (15--30 min) pour dépistage massif/autotest. Séquençage (surveillance variants).
    \item \textbf{Prise en charge :} Symptomatique (antipyrétiques, antitussifs, O2, corticoïdes \emph{dexaméthasone} si O2 requis, anticoagulants préventifs, anticorps monoclonaux \emph{selon variant}, antiviraux \emph{nirmatrelvir/ritonavir (Paxlovid\textsuperscript{\textregistered}), molnupiravir} précocement chez sujets à risque).
    \item \textbf{Vaccination (Outil majeur) :} Vaccins ARNm (\emph{BNT162b2 - Pfizer/BioNTech}, \emph{mRNA-1273 - Moderna}), vecteur viral (\emph{ChAdOx1 - AstraZeneca/Covishield}, \emph{Ad26.COV2.S - Janssen}), protéine recombinante (\emph{NVX-CoV2373 - Novavax}), inactivé (\emph{CoronaVac - Sinovac}, \emph{BBIBP-CorV - Sinopharm}). Au Cameroun : AstraZeneca, Sinopharm, Pfizer, Janssen, Pfizer Pédiatrique. \textbf{Calendrier :} Primo-vaccination (1 ou 2 doses selon vaccin) + Rappels (6--12 mois) pour populations cibles (personnes âgées, comorbidités, personnel soignant). Efficacité élevée contre formes graves/hospitalisation (> 90\,\% maintenue), variable contre infection symptomatique.
\end{enumerate}

\begin{aretenir}[Comparatif Ebola vs COVID-19]
\begin{center}
\begin{tabularx}{\linewidth}{|l|X|X|}\hline
\textbf{Caractère} & \textbf{Fièvre Ebola (FHVE)} & \textbf{COVID-19 (SARS-CoV-2)} \\ \hline
\textbf{Agent} & Filovirus (ARN, enveloppé) & Coronavirus (ARN, enveloppé) \\ \hline
\textbf{Réservoir} & Chauves-souris frugivores & Chauves-souris (probable) \\ \hline
\textbf{Transmission} & Contact direct fluides (malade/mort) & Aérienne (aérosols/gouttelettes) + Contact \\ \hline
\textbf{Contagiosité (R0)} & Faible à modérée (1,5--2,5) & Élevée (3--10+ selon variants) \\ \hline
\textbf{Létalité (CFR)} & Élevée (25--90\,\%) & Faible à modérée (0,1--3\,\%, selon variant/âge/vaccin) \\ \hline
\textbf{Vaccin} & rVSV-ZEBOV (Ervebo\textsuperscript{\textregistered}) - Anneau & Multiples plateformes (ARNm, vecteur, etc.) - Masse \\ \hline
\textbf{Traitement spécifique} & Anticorps monoclonaux (Mab114, Regn-EB3) & Antiviraux oraux (Paxlovid\textsuperscript{\textregistered}), Anticorps (selon variant) \\ \hline
\textbf{Mesure clé rupture} & Isolement cas + Enterrements sécurisés + PCI & Masque + Aération + Vaccination + Isolement cas \\ \hline
\end{tabularx}
\end{center}
\end{aretenir}

\coursec{Le choléra : épidémie hydrique et synthèse}

\courssub{Agent causal et transmission}

\paragraph{Agent causal.}
\textbf{\emph{Vibrio cholerae}} : bacille Gram négatif, mobile (flagelle polaire), courbe « en virgule », aérobie, halophile (aime le sel). Deux sérogroupes épidémiques : \textbf{O1} (sous-types Ogawa, Inaba, Hikojima) et \textbf{O139} (Bengale, 1992). Production d'une \textbf{entérotoxine cholérique (Choléra Toxine - CT)} : sous-unité A (activité ADP-ribosyltransférase) + 5 sous-unités B (liaison récepteur GM1).

\paragraph{Mode de transmission : Hydrique et Alimentaire (Fécalo-orale).}
\begin{itemize}
    \item \textbf{Eau contaminée :} Eau de boisson (puits, forages, rivières, réseau de distribution défaillant) souillée par des matières fécales (assainissement inexistant ou défaillant, inondations).
    \item \textbf{Aliments contaminés :} Crudités lavées à l'eau souillée, fruits non pelés, poisson/crustacés crus ou mal cuits (milieu salin favorable), riz/repas réchauffés insuffisamment (milieu de culture).
    \item \textbf{Contact direct :} Mains souillées (déjections, vomissures) $\to$ bouche (hygiène défaillante). Portage asymptomatique fréquent (ratio asymptomatique/symptomatique $\approx$ 10/1 à 100/1).
\end{itemize}
\textbf{Facteurs de risque au Cameroun :} Saison des pluies (inondations, contamination nappes/puits), quartiers précaires (bidonvilles) sans assainissement, déplacements de populations (crises Nord-Ouest/Sud-Ouest, Extrême-Nord), enterrements traditionnels (lavage corps), pénurie d'eau potable.

\courssub{Physiopathologie et manifestations}

\paragraph{Mécanisme de la diarrhée (Savoir admis).}
\emph{V. cholerae} colonise l'intestin grêle (duodénum, jéjunum) $\to$ fixation aux récepteurs GM1 des entérocytes via sous-unités B $\to$ entrée sous-unité A $\to$ activation adénylate cyclase $\to$ augmentation \textbf{AMPc} intracellulaire $\to$ ouverture canaux \textbf{CFTR} (Chlorure) $\to$ sortie massive \ce{Cl-}, \ce{Na+}, \ce{HCO3-}, \ce{K+} $\to$ entrée d'eau par osmose $\to$ \textbf{diarrhée sécrétoire massive, isotone, indolore}.

\begin{definition}[Choléra typique]
\begin{itemize}
    \item \textbf{Diarrhée profuse :} « Eau de riz » (liquide, incolore, éclats de mucus, sans selles formées, sans sang, sans fièvre, indolore). Volume : \textbf{10 à 20 L/jour} (jusqu'à \SI{1}{\liter\per\hour}).
    \item \textbf{Vomissements :} En jet, sans nausées, survenant après la diarrhée.
    \item \textbf{Déshydratation aiguë sévère :} Soif intense, peau « en gants de boxeur » (pli cutané lent), yeux enfoncés, voix éteinte (aphonie), oligurie/anurie, crampes musculaires (perte \ce{K+}), hypotension, tachycardie, troubles de la conscience (choc hypovolémique).
    \item \textbf{Acidose métabolique :} Perte \ce{HCO3-} $\to$ respiration de Kussmaul.
    \item \textbf{Hypokaliémie :} Perte \ce{K+} $\to$ iléus paralytique, troubles cardiaques (arythmies).
\end{itemize}
Sans traitement : létalité \textbf{25--50\,\%} (choc hypovolémique). Avec réhydratation adéquate : \textbf{< 1\,\%}.
\end{definition}

\courssub{Moyens de lutte : La réhydratation, clé de la survie}

\begin{propriete}[Prise en charge du choléra (OMS / PNLC Cameroun)]
\begin{enumerate}
    \item \textbf{Évaluation de la déshydratation (Plan A, B, C) :}
    \begin{itemize}
        \item \textbf{Plan A (Pas de déshydratation) :} Traitement à domicile. SRO (Sel de Réhydratation Orale) après chaque selle (\SI{200}{\milli\liter}--\SI{500}{\milli\liter} enfant, \SI{1000}{\milli\liter} adulte). Zinc (\SI{20}{\milli\gram}/j 10--14 j < 5 ans). Conseil : continuer alimentation, eau traitée, retour si aggravation.
        \item \textbf{Plan B (Déshydratation modérée) :} SRO au centre de santé (CTC/CTU) : \SI{75}{\milli\liter\per\kilo\gram} en 4 h (enfant) / 4 h (adulte). Surveillance toutes les 15--30 min. Si échec $\to$ Plan C.
        \item \textbf{Plan C (Déshydratation sévère / Choc) :} \textbf{URGENCE VITALE}. Voie IV (perfusion) : \textbf{Ringer Lactate (RL)} ou \textbf{Saline Normale (0,9\,\% NaCl)}. \textbf{Bolus initial :} \SI{30}{\milli\liter\per\kilo\gram} en \textbf{30 min} (enfant \SI{1}{\hour}) $\to$ puis \SI{70}{\milli\liter\per\kilo\gram} en \SI{2.5}{\hour} (enfant \SI{5}{\hour}). Réévaluer toutes les 15--30 min. Passer au SRO (Plan B) dès que le patient boit (généralement 3--4 h). \textbf{Potassium :} Ajouter KCl dans les perfusions de relais (RL en contient \SI{4}{\milli\mole\per\liter}, souvent insuffisant) ou SRO (contient \SI{15}{\milli\mole\per\liter} \ce{K+}) dès reprise per os.
    \end{itemize}
    \item \textbf{Antibiothérapie (Adjunctive, réduit durée/volume selles/portage) :} \textbf{Doxycycline} (\SI{300}{\milli\gram} dose unique adulte) ou \textbf{Azithromycine} (\SI{20}{\milli\gram\per\kilo\gram} dose unique enfant/femme enceinte). Résistances fréquentes aux tétracyclines, cotrimoxazole, ciprofloxacine $\to$ adapter selon antibiogramme local.
    \item \textbf{Alimentation :} Reprise immédiate (lait maternel, bouillies épaisses, riz, banane plantain). Pas de jeûne.
\end{enumerate}
\end{propriete}

\paragraph{Prévention et contrôle de l'épidémie (Approche multisectorielle).}
\begin{enumerate}
    \item \textbf{Eau, Hygiène, Assainissement (EHA) :} \textbf{Accès à l'eau potable} (chloration point de puisage/réseau \ce{Cl2} résiduel \SIrange{0.2}{0.5}{\milli\gram\per\liter}, traitement ménager : Aquatabs\textsuperscript{\textregistered}, PuR\textsuperscript{\textregistered}, ébullition, filtration céramique). \textbf{Assainissement :} Latrines améliorées, gestion boues de vidange, drainage eaux usées. \textbf{Hygiène :} Lavage mains au savon (postes de lavage), cuisson complète aliments, pelage fruits, protection aliments (mouches).
    \item \textbf{Surveillance / Riposte :} Déclaration obligatoire (Maladie à Déclaration Obligatoire - MDO). Définition de cas standard. CTC (Centre de Traitement du Choléra) / CTU (Unité) isolés. Ligne de liste (Line list). Enquête de terrain (source, propagation).
    \item \textbf{Vaccination orale anti-cholérique (VOC) :} Vaccins inactivés bivalents O1/O139 (Shanchol\textsuperscript{\textregistered}, Euvichol\textsuperscript{\textregistered}, Euvichol-Plus\textsuperscript{\textregistered}). 2 doses (14 j d'intervalle) $\to$ protection 3--5 ans. 1 dose (urgence) $\to$ protection $\approx$ 6--12 mois. Campagnes réactives (foyers) ou préventives (zones à haut risque : Nord, Extrême-Nord, Littoral). \textbf{Ne remplace pas les mesures EHA.}
    \item \textbf{Communication / Mobilisation sociale :} Messages clairs (Eau traitée, Mains propres, Latrines, Enterrements sécurisés, Dépistage précoce). Lutte contre la stigmatisation des malades/guéris.
\end{enumerate}

\begin{exemplebox}[Préparation de la Solution de Réhydratation Orale (SRO) - Norme OMS 2002 (Faible osmolarité 245~mOsm/L)]
\begin{center}
\begin{tabular}{|c|c|c|}\hline
\textbf{Composant} & \textbf{Concentration (mmol/L)} & \textbf{Poids pour 1 L} \\ \hline
Chlorure de sodium (\ce{NaCl}) & 75 & \SI{2.6}{\gram} \\ \hline
Chlorure de potassium (\ce{KCl}) & 20 & \SI{1.5}{\gram} \\ \hline
Citrate de sodium trisubstitué dihydraté (\ce{Na3C6H5O7.2H2O}) & 10 & \SI{2.9}{\gram} \\ \hline
Glucose (anhydre) & 75 & \SI{13.5}{\gram} \\ \hline
\textbf{Total osmolarité} & \textbf{245} & \\ \hline
\end{tabular}
\end{center}
\textbf{Mode d'emploi :} Dissoudre 1 sachet dans \textbf{1 L d'eau potable} (bouillie et refroidie ou traitée au chlore). Ne pas ajouter sucre, sel, jus. Jeter après 24 h (6 h si T > \SI{25}{\celsius}). Goût : légèrement salé/sucré.
\end{exemplebox}

\courssub{Synthèse de l'unité : Endémies et Épidémies}

\begin{aretenir}[Tableau récapitulatif des maladies du programme (3ème)]
\begin{center}
\small
\begin{tabularx}{\linewidth}{|l|c|X|X|X|}\hline
\textbf{Maladie} & \textbf{Type} & \textbf{Agent} & \textbf{Transmission} & \textbf{Moyens de lutte clés (Programme 3ème)} \\ \hline
\textbf{Paludisme} & Endémie & Parasite (\emph{Plasmodium} spp.) & Vectorielle : Piqûre \emph{Anopheles} $\varphi$ & MILDA, PID, LAL, TPE (ACT), CPS, CPG, Vaccin RTS,S \\ \hline
\textbf{Fièvre Ebola} & Épidémie & Virus (Filovirus \emph{Ebolavirus}) & Contact direct fluides (malade/mort) & Isolement, PCI/EPI, Enterrements sécurisés, Vaccin rVSV-ZEBOV (Anneau), Anticorps monoclonaux \\ \hline
\textbf{COVID-19} & Épidémie/Pandémie & Virus (Coronavirus SARS-CoV-2) & Aérienne (aérosols/gouttelettes) + Contact & Masque, Aération, Distanciation, Hygiène mains, Vaccination (ARNm, vecteur...), Antiviraux précoces \\ \hline
\textbf{Choléra} & Épidémie & Bactérie (\emph{Vibrio cholerae} O1/O139) & Hydrique / Alimentaire (Fécalo-orale) & \textbf{Réhydratation (SRO / RL Plan A/B/C)}, Antibio (Doxy/Azithro), EHA (Eau/Assainissement/Hygiène), VOC \\ \hline
\end{tabularx}
\end{center}
\end{aretenir}

\begin{exercice}[Activités d'intégration - Synthèse Unité 11]
\begin{enumerate}
    \item \textbf{Analyse de situation :} Dans un quartier périphérique de Yaoundé, on observe une augmentation brutale de cas de diarrhée aqueuse « eau de riz » avec vomissements en saison des pluies, après une inondation ayant détruit les latrines et contaminé les puits.
    \begin{enumerate}
        \item Quelle maladie suspectez-vous ? Justifiez par 3 arguments épidémiologiques/cliniques.
        \item Précisez le mode de transmission et le réservoir.
        \item Proposez 3 mesures d'urgence à mettre en œuvre (1 médicale, 1 environnementale, 1 communautaire).
    \end{enumerate}
    \item \textbf{Comparaison :} Complétez le tableau comparatif ci-dessous pour le paludisme et le choléra.
    \begin{center}
    \begin{tabularx}{\linewidth}{|l|X|X|}\hline
    \textbf{Critère} & \textbf{Paludisme} & \textbf{Choléra} \\ \hline
    Type d'épidémiologie & & \\ \hline
    Agent causal (règne) & & \\ \hline
    Mode de transmission & & \\ \hline
    Organe cible principal & & \\ \hline
    Signe clinique majeur & & \\ \hline
    Traitement curatif de 1ère ligne & & \\ \hline
    Mesure préventive individuelle majeure & & \\ \hline
    \end{tabularx}
    \end{center}
    \item \textbf{Vaccination :} Le Cameroun a introduit le vaccin RTS,S en 2024.
    \begin{enumerate}
        \item Contre quelle maladie ? Quel stade du parasite vise-t-il ?
        \item Pourquoi ce vaccin ne suffit-il pas seul pour éliminer le paludisme ? (2 arguments).
    \end{enumerate}
    \item \textbf{Épidémie à virus Ebola :} Un cas confirmé est détecté dans un village.
    \begin{enumerate}
        \item Citez les 3 piliers de la riposte « Anneau ».
        \item Pourquoi l'enterrement sécurisé est-il critique ? (Mécanisme de transmission).
    \end{enumerate}
\end{enumerate}
\end{exercice}

\begin{corrige}[Exercices d'intégration]
\begin{enumerate}
    \item \textbf{Choléra.}
    \begin{enumerate}
        \item Arguments : (1) Saison des pluies + inondation (facteur déclenchant hydrique), (2) Diarrhée « eau de riz » profuse + vomissements en jet + déshydratation rapide sans fièvre (tableau clinique typique), (3) Contexte assainissement défaillant (puits contaminés, latrines détruites).
        \item Transmission : Hydrique (eau de boisson contaminée fèces) et alimentaire (fécalo-orale). Réservoir : Homme (malade + porteur sain) + Environnement aquatique (eau saumâtre, plancton).
        \item Mesures : (Médicale) Installation CTC + Réhydratation SRO/RL (Plan A/B/C) + Doxycycline. (Environnementale) Chloration points d'eau / distribution Aquatabs + Interdiction puits contaminés + Assainissement d'urgence (latrines d'urgence). (Communautaire) Sensibilisation lavage mains / eau traitée / cuisson / enterrements sécurisés + Recherche active cas / VOC réactive.
    \end{enumerate}
    \item
    \begin{center}
    \begin{tabularx}{\linewidth}{|l|X|X|}\hline
    \textbf{Critère} & \textbf{Paludisme} & \textbf{Choléra} \\ \hline
    Type d'épidémiologie & Endémie (stable, saisonnière) & Épidémie (brutale, explosive) \\ \hline
    Agent causal (règne) & Parasite (Protozoaire \emph{Plasmodium}) & Bactérie (\emph{Vibrio cholerae}) \\ \hline
    Mode de transmission & Vectorielle (Piqûre \emph{Anopheles} $\varphi$) & Hydrique / Alimentaire (Fécalo-orale) \\ \hline
    Organe cible principal & Foie (phase asymptomatique) puis Sang (GR) & Intestin grêle (entérocytes) \\ \hline
    Signe clinique majeur & Fièvre paroxysmique (accès 48h) + Splénomégalie & Diarrhée « eau de riz » profuse + Déshydratation aiguë \\ \hline
    Traitement curatif de 1ère ligne & ACT : Artéméthrine-Luméfantrine (AL) & Réhydratation (SRO / Ringer Lactate) + Doxycycline/Azithromycine \\ \hline
    Mesure préventive individuelle majeure & MILDA (Moustiquaire Imprégnée Longue Durée) & Eau potable traitée (chloration/ébullition) + Lavage mains savon \\ \hline
    \end{tabularx}
    \end{center}
    \item
    \begin{enumerate}
        \item Maladie : \textbf{Paludisme} (due à \emph{P. falciparum}). Stade visé : \textbf{Sporozoïte} (phase pré-érythrocytaire) via anticorps anti-protéine CSP.
        \item Arguments : (1) Efficacité partielle seulement ($\approx$ 30--40\,\% sur paludisme grave), ne protège pas tous les vaccinés. (2) Ne cible pas les stades sanguins (symptômes) ni les gamétocytes (transmission) ; nécessite le maintien de la LAV (MILDA) et du TPE.
    \end{enumerate}
    \item
    \begin{enumerate}
        \item (1) \textbf{Détection/Isolation} précoce des cas (CST) + Prise en charge (Anticorps monoclonaux, réanimation). (2) \textbf{Vaccination Anneau} (rVSV-ZEBOV) : contacts + contacts des contacts + personnel soignant. (3) \textbf{PCI/Enterrements sécurisés} + Surveillance contacts (21 j) + Communication communautaire.
        \item Le corps d'un défunt Ebola est \textbf{très fortement virulent} (charge virale maximale dans tous les fluides : sang, selles, sueur, peau). Les rites funéraires traditionnels (lavage, toucher, embrassement, veillée) exposent massivement la famille/assistants $\to$ amplification explosive de l'épidémie. L'enterrement sécurisé (équipe formée, EPI, sac hermétique, dignité respectée) coupe cette chaîne majeure.
    \end{enumerate}
\end{enumerate}
\end{corrige}

\coursec{Le paludisme : manifestations et moyens de lutte}

Dans la section précédente, nous avons identifié l'agent causal du paludisme (un parasite microscopique, le \emph{Plasmodium}, transmis par la piqûre d'un moustique femelle du genre \emph{Anopheles}) et décrit son cycle de transmission entre l'Homme et le moustique vecteur. Nous comprenons maintenant \emph{comment} le parasite nous infecte. Reste à répondre à trois questions concrètes : quels sont les \cle{signes} de la maladie qui doivent alerter ? Comment le personnel de santé \cle{confirme}-t-il le diagnostic ? Et surtout, comment \cle{lutter} efficacement contre cette endémie qui frappe durement nos communautés camerounaises ?

\courssub{Les manifestations du paludisme}

Lorsqu'un moustique \emph{Anopheles} infecté pique un homme, il injecte le parasite dans le sang. Après une première multiplication silencieuse dans le foie, le parasite envahit les \cle{globules rouges}, s'y multiplie, puis les détruit. C'est la destruction des globules rouges qui provoque les signes de la maladie.

\begin{definition}[Symptôme]
Un \textbf{symptôme} est une manifestation d'une maladie, ressentie par le malade (douleur, fatigue) ou observable par le médecin (fièvre, pâleur). L'ensemble des symptômes constitue le \textbf{tableau clinique} de la maladie.
\end{definition}

Le paludisme simple (non compliqué) se manifeste par :
\begin{itemize}
    \item une \textbf{fièvre élevée} (souvent supérieure à $38,5\,^\circ\mathrm{C}$), qui survient par \textbf{accès} : le malade a très chaud, puis la température redescend, puis un nouvel accès survient ;
    \item des \textbf{frissons} intenses (sensation de froid violent) suivis de \textbf{sueurs abondantes} ;
    \item des \textbf{maux de tête}, une \textbf{fatigue} intense, des \textbf{nausées} et des \textbf{vomissements}.
\end{itemize}

\begin{propriete}[Origine de la fièvre]
La fièvre du paludisme résulte de la \textbf{destruction des globules rouges} infectés : lorsque les parasites sortent des globules rouges en grand nombre, des substances toxiques sont libérées dans le sang, ce qui déclenche la fièvre, les frissons puis les sueurs (savoir admis).
\end{propriete}

\begin{experience}[Observation de la fièvre chez un malade]
\textbf{Contexte :} Un enfant consulte dans un centre de santé de Yaoundé pour de la fièvre. L'infirmier relève sa température toutes les 6 heures pendant 24 heures.
\textbf{Résultats :}
\begin{center}
\begin{tabularx}{\linewidth}{|l|X|X|X|X|X|}\hline
Heure & 6 h & 12 h & 18 h & 0 h & 6 h \\ \hline
Température ($^\circ\mathrm{C}$) & 37,2 & 39,8 & 37,5 & 37,0 & 37,1 \\ \hline
\end{tabularx}
\end{center}
\textbf{Interprétation :} On observe un pic de fièvre ($39,8\,^\circ\mathrm{C}$) suivi d'un retour à une température normale. Cette fièvre qui va et vient est typique du paludisme : elle correspond aux moments où les parasites détruisent les globules rouges et se libèrent dans le sang.
\textbf{Conclusion :} Une fièvre élevée par accès, en zone où sévit le paludisme, doit faire penser à cette maladie et conduire à consulter.
\end{experience}

\begin{attention}[Toute fièvre n'est pas un paludisme !]
Au Cameroun, beaucoup de maladies donnent de la fièvre : la typhoïde, la grippe, la COVID-19, les infections respiratoires... \textbf{On ne peut pas être sûr qu'il s'agit d'un paludisme sans examen.} Il faut donc consulter un centre de santé pour faire confirmer le diagnostic avant tout traitement, et éviter l'automédication.
\end{attention}

\courssub{Les formes graves et les personnes les plus exposées}

En l'absence de traitement rapide, le paludisme simple peut évoluer vers un \textbf{paludisme grave}, qui est une urgence pouvant entraîner la mort.

\begin{aretenir}[Signes de gravité à reconnaître]
\begin{itemize}
    \item \textbf{Troubles de la conscience} : le malade ne répond plus, confusion, \textbf{convulsions} répétées (paludisme cérébral).
    \item \textbf{Pâleur importante} (paupières, paumes des mains) : signe d'\textbf{anémie sévère} due à la destruction massive des globules rouges.
    \item \textbf{Difficultés respiratoires}, respiration rapide et profonde.
    \item \textbf{Jaunisse} (yeux et peau jaunes) et urines très foncées.
\end{itemize}
Devant l'un de ces signes, il faut conduire le malade \textbf{en urgence} à l'hôpital.
\end{aretenir}

\begin{savaistu}[Les personnes les plus vulnérables]
Les \textbf{enfants de moins de 5 ans} et les \textbf{femmes enceintes} sont les plus exposés aux formes graves. Chez la femme enceinte, le paludisme peut provoquer une anémie sévère, un accouchement prématuré ou un bébé de faible poids à la naissance. C'est pourquoi, au Cameroun, les femmes enceintes reçoivent gratuitement un traitement préventif lors des consultations prénatales, ainsi qu'une moustiquaire imprégnée.
\end{savaistu}

\courssub{Le diagnostic : confirmer avant de traiter}

Le seul examen des symptômes ne suffit pas, car ils ressemblent à ceux d'autres maladies. Le personnel de santé dispose de deux moyens de confirmation :

\begin{methode}[Confirmer un paludisme au centre de santé]
\begin{enumerate}
    \item \textbf{Le Test de Diagnostic Rapide (TDR) :} on prélève une goutte de sang au bout du doigt ; le test détecte la présence du parasite en 15 à 20 minutes. Il est utilisable partout, même sans laboratoire.
    \item \textbf{La goutte épaisse :} on étale une goutte de sang sur une lame de verre, on la colore, puis on l'observe au \textbf{microscope} pour rechercher directement le parasite dans les globules rouges. C'est l'examen de référence, réalisé au laboratoire.
\end{enumerate}
\textbf{Règle d'or :} on teste d'abord, on traite ensuite.
\end{methode}

\begin{exemplebox}[Lecture d'un Test de Diagnostic Rapide]
\begin{center}
\begin{tikzpicture}[scale=0.75, every node/.style={font=\small}]
% Cas positif
\begin{scope}
\draw[thick, fill=white, rounded corners] (0,0) rectangle (5,2);
\draw[thick] (1.2,0) -- (1.2,2);
\node at (0.6,1.5) {\textbf{C}};
\node at (0.6,0.5) {\textbf{T}};
\draw[fill=popink] (0.45,1.35) rectangle (0.75,1.65);
\draw[fill=popink] (0.45,0.35) rectangle (0.75,0.65);
\node at (3,1) {Fenêtre de lecture};
\node at (2.5,-0.6) {\textbf{Positif} : bandes C et T};
\end{scope}
% Cas négatif
\begin{scope}[xshift=7cm]
\draw[thick, fill=white, rounded corners] (0,0) rectangle (5,2);
\draw[thick] (1.2,0) -- (1.2,2);
\node at (0.6,1.5) {\textbf{C}};
\node at (0.6,0.5) {\textbf{T}};
\draw[fill=popink] (0.45,1.35) rectangle (0.75,1.65);
\node at (3,1) {Fenêtre de lecture};
\node at (2.5,-0.6) {\textbf{Négatif} : bande C seule};
\end{scope}
% Cas invalide
\begin{scope}[xshift=14cm]
\draw[thick, fill=white, rounded corners] (0,0) rectangle (5,2);
\draw[thick] (1.2,0) -- (1.2,2);
\node at (0.6,1.5) {\textbf{C}};
\node at (0.6,0.5) {\textbf{T}};
\node at (3,1) {Fenêtre de lecture};
\node at (2.5,-0.6) {\textbf{Invalide} : pas de bande C};
\end{scope}
\end{tikzpicture}
\end{center}
\legende{Principe de lecture d'un TDR du paludisme. C = bande de contrôle (elle prouve que le test a fonctionné) ; T = bande de test (elle révèle la présence du parasite). Si la bande C n'apparaît pas, le test est invalide et doit être refait.}
\end{exemplebox}

\courssub{La lutte préventive : une approche combinée}

Le paludisme étant transmis par un moustique, la prévention repose sur trois idées simples : \textbf{se protéger des piqûres}, \textbf{détruire le moustique et ses larves}, \textbf{traiter les malades} pour couper la transmission.

\begin{methode}[Les moyens de lutte contre le paludisme]
\begin{enumerate}
    \item \textbf{Dormir sous moustiquaire imprégnée d'insecticide (MII)} : la moustiquaire est une barrière physique, et l'insecticide qu'elle contient tue ou repousse le moustique. C'est le moyen de protection principal, distribué gratuitement lors des campagnes nationales. Elle doit être utilisée \textbf{toutes les nuits, toute l'année}, par toute la famille, en priorité par les enfants et les femmes enceintes.
    \item \textbf{Détruire les gîtes larvaires} : la larve du moustique se développe dans les \textbf{eaux stagnantes} (flaques, marécages, ornières, pneus usagés, boîtes et bouteilles abandonnées, pieds de bananiers creux). Actions : combler les flaques, drainer, vider ou recouvrir les récipients d'eau, ramasser les déchets qui retiennent l'eau de pluie.
    \item \textbf{Assainir l'environnement} : débroussailler autour des maisons, curer les caniveaux, gérer correctement les ordures.
    \item \textbf{Pulvérisations d'insecticides} à l'intérieur des habitations, réalisées par des équipes spécialisées dans certaines zones.
    \item \textbf{Traitement préventif} des femmes enceintes (lors des consultations prénatales) et, dans les régions du Nord à saison des pluies marquée, des enfants pendant la saison de forte transmission.
\end{enumerate}
\end{methode}

\begin{popfigure}
\begin{tikzpicture}[scale=0.9, every node/.style={font=\small}]
% Maison
\draw[fill=poporangeL, thick] (0.5,3) rectangle (3.5,5.5);
\draw[thick, fill=popgoldL] (1.7,3) rectangle (2.5,4); % porte
\draw[fill=popblueL] (0.8,4.3) rectangle (1.4,4.9); % fenêtre
\draw[fill=popblueL] (2.8,4.3) rectangle (3.3,4.9); % fenêtre
\draw[thick, fill=popgold] (0.3,5.5) -- (2,6.4) -- (3.7,5.5) -- cycle; % toit
\node at (2,6.8) {\textbf{Maison}};
% Moustiquaire
\draw[fill=popblueL, pattern=north east lines, pattern color=popblue] (0.7,3.1) rectangle (1.6,4.1);
\node[align=center] at (1.15,2.75) {Moustiquaire\\imprégnée};
% Gîte larvaire
\draw[fill=popblueL, thick] (5.5,1) .. controls (6.5,0.5) and (8.5,0.5) .. (9.5,1) .. controls (10,2) and (9,3) .. (7.5,3.3) .. controls (6,3) and (5,2) .. (5.5,1);
\draw[fill=popgreen] (6.3,1.6) circle (0.18);
\draw[fill=popgreen] (7.3,1.9) circle (0.18);
\draw[fill=popgreen] (8.3,1.6) circle (0.18);
\draw[fill=popgreen] (7.6,2.4) circle (0.18);
\node at (7.5,0.1) {\textbf{Gîte larvaire (eau stagnante)}};
\node[popgreen!60!popdark] at (7.5,1.9) {Larves};
% Actions sur le gîte
\draw[->, very thick, popdark] (10,2.4) -- (11.3,2.9);
\node[align=center, fill=white, rounded corners, draw=popdark] at (12.6,3.1) {Comblage /\\drainage};
\draw[->, very thick, popdark] (10,1.4) -- (11.3,1);
\node[align=center, fill=white, rounded corners, draw=popdark] at (12.6,0.8) {Vidange des\\récipients};
% Déchets
\draw[fill=popmuted!40] (0.8,0.6) rectangle (1.6,1.3);
\draw[fill=popmuted!40] (2,0.6) rectangle (2.7,1.3);
\node at (1.75,0.2) {Déchets (pneus, boîtes)};
\draw[->, very thick, popdark] (1.75,1.35) -- (1.75,2.1);
\node[align=center, fill=white, rounded corners, draw=popdark] at (1.75,2.45) {Ramassage /\\élimination};
% Moustique
\node[align=center] at (4.6,5.6) {Moustique\\(\emph{Anopheles})};
\draw[->, thick, poppink] (4.3,5.2) -- (3.2,4.6);
\draw[->, thick, poppink] (5,5.2) -- (7.2,3.4);
\end{tikzpicture}
\legende{Schéma synthétique de la lutte contre le paludisme : protection de l'homme (moustiquaire imprégnée dans la maison) et destruction des gîtes larvaires (comblage, drainage, vidange des récipients, élimination des déchets qui retiennent l'eau).}
\end{popfigure}

\begin{exemplebox}[L'efficacité des moustiquaires imprégnées]
Selon l'Organisation Mondiale de la Santé (OMS), l'usage régulier et correct des \textbf{moustiquaires imprégnées} peut réduire d'environ \textbf{50\,\% les cas de paludisme} et d'environ \textbf{20\,\% la mortalité} chez les enfants de moins de 5 ans en zone d'endémie. Au Cameroun, les campagnes nationales de distribution gratuite de moustiquaires ont contribué à la baisse de la fréquence du paludisme chez les enfants.
\end{exemplebox}

\courssub{La lutte curative : traiter correctement}

Le traitement vise à éliminer le parasite du sang du malade et à éviter les formes graves.

\begin{propriete}[Le traitement du paludisme simple]
Au Cameroun, le paludisme simple confirmé se traite avec des \textbf{médicaments antipaludiques associés} (combinaisons à base d'artémisinine), prescrits par le personnel de santé et pris pendant \textbf{3 jours complets}. Le paludisme grave est une urgence hospitalière : le malade reçoit un traitement injectable.
\end{propriete}

\begin{attention}[Les erreurs à éviter absolument]
\begin{itemize}
    \item \textbf{Les antibiotiques ne guérissent pas le paludisme} : le \emph{Plasmodium} est un parasite, pas une bactérie.
    \item \textbf{Ne jamais interrompre le traitement} quand la fièvre tombe : il faut prendre la totalité des 3 jours, sinon des parasites survivent, la maladie revient, et le parasite devient \textbf{résistant} au médicament.
    \item \textbf{Éviter l'automédication} et les médicaments achetés dans la rue (souvent de mauvaise qualité ou falsifiés) : consulter un centre de santé.
    \item Les tisanes et remèdes traditionnels peuvent soulager, mais \textbf{ne remplacent pas} le traitement prescrit.
\end{itemize}
\end{attention}

\courssub{Analyse de données : l'impact des moustiquaires}

Pour évaluer l'efficacité d'une campagne de distribution de moustiquaires, on compare la mortalité due au paludisme chez les enfants de moins de 5 ans avant et après la campagne.

\begin{popfigure}
\begin{center}
\begin{tikzpicture}
\begin{axis}[
    ybar,
    bar width=30pt,
    width=0.9\linewidth,
    height=8cm,
    ylabel={Décès par paludisme (pour 1000 enfants de moins de 5 ans)},
    symbolic x coords={Avant (2010), Après (2016)},
    xtick=data,
    nodes near coords,
    ymin=0, ymax=100,
    ymajorgrids=true,
    grid style=dashed,
    fill=popblue!70,
]
\addplot coordinates {(Avant (2010), 85) (Après (2016), 42)};
\end{axis}
\end{tikzpicture}
\end{center}
\legende{Mortalité due au paludisme chez les enfants de moins de 5 ans dans une région camerounaise, avant et après une campagne de distribution de moustiquaires imprégnées (données fictives inspirées des tendances réelles observées au Cameroun).}
\end{popfigure}

\begin{exoresolu}[Interpréter l'impact des moustiquaires]
\textbf{Document :} le graphique ci-dessus montre l'évolution de la mortalité par paludisme chez les enfants de moins de 5 ans, avant (2010) et après (2016) une campagne de distribution de moustiquaires imprégnées.
\textbf{Questions :}
\begin{enumerate}
    \item Calculer le pourcentage de baisse de la mortalité entre 2010 et 2016.
    \item Expliquer biologiquement pourquoi la moustiquaire imprégnée fait baisser la mortalité.
    \item Pourquoi la mortalité n'est-elle pas tombée à zéro ? Donner deux raisons.
\end{enumerate}
\tcblower
\textbf{Solution détaillée :}
\begin{enumerate}
    \item Mortalité initiale : $M_i = 85$ décès pour 1000 enfants ; mortalité finale : $M_f = 42$.
    \[ \text{Baisse} = \frac{M_i - M_f}{M_i} \times 100 = \frac{85 - 42}{85} \times 100 = \frac{43}{85} \times 100 \approx 50,6\,\% \]
    La mortalité a donc baissé d'environ \textbf{51\,\%}.
    \item La moustiquaire est une \textbf{barrière physique} qui empêche le moustique de piquer pendant le sommeil (l'\emph{Anopheles} pique surtout la nuit) ; de plus, l'\textbf{insecticide} dont elle est imprégnée tue ou repousse les moustiques qui se posent dessus. Moins de piqûres signifie moins de contaminations, donc moins de malades et moins de décès.
    \item La mortalité n'est pas nulle car : (a) certaines familles \textbf{n'utilisent pas} la moustiquette toutes les nuits (chaleur, oubli, moustiquaire déchirée) ; (b) le moustique peut piquer \textbf{en début de soirée}, avant le coucher, ou à l'extérieur ; (c) certains moustiques deviennent \textbf{résistants} aux insecticides ; (d) des malades arrivent trop tard au centre de santé.
\end{enumerate}
\end{exoresolu}

\courssub{Perspectives : un vaccin contre le paludisme}

\begin{savaistu}[Le vaccin antipaludique arrive au Cameroun]
Après des décennies de recherche, un premier \textbf{vaccin contre le paludisme} a été recommandé par l'OMS et introduit au Cameroun à partir de 2024, dans les districts de santé les plus touchés, pour les jeunes enfants. \textbf{Message clé :} ce vaccin ne protège pas à 100\,\% et \textbf{ne remplace pas} les moustiquaires, l'assainissement ni les traitements : c'est un outil \textbf{supplémentaire} qui s'ajoute aux moyens de lutte existants.
\end{savaistu}

\begin{exercice}[Prise en charge d'un cas suspect]
Un enfant de 3 ans, vivant à Yaoundé, est amené au centre de santé : fièvre à $39,2\,^\circ\mathrm{C}$ depuis 24 heures, vomissements, refus de manger. Il est conscient, boit correctement et n'a pas convulsé. Le test de diagnostic rapide est \textbf{positif}.
\begin{enumerate}
    \item S'agit-il d'un paludisme simple ou grave ? Justifier.
    \item Que doit faire le personnel de santé, et que doit-on conseiller à la mère concernant le traitement ?
    \item Citer trois conseils de prévention à donner à la famille.
\end{enumerate}
\end{exercice}

\begin{corrige}[Exercice : Prise en charge d'un cas suspect]
\begin{enumerate}
    \item Il s'agit d'un \textbf{paludisme simple} : fièvre + test positif, mais \textbf{absence de signes de gravité} (l'enfant est conscient, ne convulse pas, boit normalement, pas de pâleur signalée).
    \item Le personnel de santé prescrit un \textbf{traitement antipaludique adapté au poids de l'enfant, pendant 3 jours complets}. Conseils à la mère : donner le médicament exactement comme prescrit, \textbf{ne pas arrêter} quand la fièvre disparaît ; si l'enfant vomit juste après une prise, recommencer la dose ; revenir immédiatement si l'état s'aggrave (convulsions, prostration, pâleur, refus total de boire).
    \item Conseils de prévention : (a) faire dormir l'enfant (et toute la famille) \textbf{sous moustiquaire imprégnée toutes les nuits} ; (b) \textbf{vider, recouvrir ou éliminer} les récipients et déchets qui retiennent l'eau autour de la maison (gîtes larvaires) ; (c) débroussailler et curer les caniveaux autour de la concession ; (d) consulter rapidement en cas de fièvre.
\end{enumerate}
\end{corrige}

\begin{aretenir}[Bilan : le paludisme, manifestations et lutte]
\begin{itemize}
    \item \textbf{Manifestations :} fièvre élevée par accès, frissons, sueurs, maux de tête, vomissements ; la fièvre vient de la destruction des globules rouges par le parasite.
    \item \textbf{Gravité :} convulsions, troubles de la conscience, pâleur (anémie), difficultés respiratoires $\rightarrow$ urgence hospitalière. Les plus exposés : enfants de moins de 5 ans et femmes enceintes.
    \item \textbf{Diagnostic :} test de diagnostic rapide (TDR) ou goutte épaisse au microscope. \textbf{On teste avant de traiter.}
    \item \textbf{Prévention :} moustiquaires imprégnées (outil principal), destruction des gîtes larvaires, assainissement, traitement préventif des femmes enceintes ; un vaccin vient s'y ajouter depuis 2024.
    \item \textbf{Traitement :} antipaludiques prescrits, 3 jours complets ; pas d'antibiotiques, pas d'automédication.
\end{itemize}
\end{aretenir}

\coursec{TP N°9 : utilisation des moustiquaires imprégnées}

Après avoir étudié les manifestations cliniques du paludisme et les grands principes de sa lutte (diagnostic précoce, traitement efficace, lutte contre le moustique vecteur), nous allons maintenant manipuler l'outil de prévention individuel le plus efficace et le plus accessible au Cameroun : la \cle{moustiquaire imprégnée de longue durée d'action (MILD)}. Ce travail pratique ne se limite pas à « savoir installer un filet » ; il vise à comprendre \emph{comment} et \emph{pourquoi} cet objet technique protège efficacement, et à identifier les erreurs d'usage qui en annulent le bénéfice. C'est une démarche d'investigation concrète : du déballage à l'installation correcte, en passant par l'entretien, nous allons valider les règles du bon usage.

\courssub{Objectif et matériel}

\begin{aretenir}[Objectif du TP]
Manipuler une moustiquaire imprégnée de longue durée (MILD), en respecter le protocole d'installation et d'entretien, et dégager les règles de bon usage qui garantissent une protection efficace contre les piqûres d'anophèles vecteurs du paludisme.
\end{aretenir}

\begin{experience}[Matériel nécessaire]
\begin{itemize}
    \item Une moustiquaire imprégnée de longue durée (MILD) neuve, encore dans son emballage d'origine (modèle rectangulaire ou conique, taille adaptée au lit : simple, double, familiale).
    \item Des ficelles solides ou des crochets de suspension (souvent fournis avec la moustiquaire).
    \item Un mètre ruban pour vérifier la hauteur de suspension.
    \item La notice d'utilisation du fabricant (indiquant le nombre maximal de lavages, la nature de l'insecticide, les précautions d'emploi).
    \item Un lit avec matelas (ou une table simulant un couchage) pour l'installation pratique.
    \item De l'eau propre, un savon doux (savon neutre), une bassine, un étendoir placé à l'ombre.
\end{itemize}
\end{experience}

\courssub{Protocole expérimental : de l'emballage au lit}

Nous suivons une démarche en cinq étapes rigoureuses. Chaque étape répond à une contrainte physique, chimique ou biologique précise.

\begin{methode}[Protocole d'installation d'une MILD en 5 étapes]
\begin{enumerate}
    \item \textbf{Déballage et aération initiale (étape cruciale).} Sortir la moustiquaire de son sac en plastique. La déplier entièrement à l'air libre, \textbf{à l'ombre}, pendant environ \textbf{24 heures} avant la première utilisation. Ne pas l'exposer au soleil direct.
    \item \textbf{Suspension.} Accrocher les points d'attache (anneaux ou boucles aux quatre coins) aux ficelles ou crochets fixés au plafond ou aux murs. La hauteur idéale : le bas de la moustiquaire doit retomber jusqu'au sol ou jusqu'au bord du matelas une fois tendue. Pour un lit standard, prévoir environ 1,50 m à 1,80 m du sol au point d'attache.
    \item \textbf{Tension de la toile.} Tirer sur les ficelles latérales pour que la toile soit bien tendue, sans plis excessifs qui créeraient des poches où les moustiques pourraient se poser sans contact franc avec l'insecticide. Une moustiquaire conique nécessite un seul point central au plafond ; une rectangulaire en nécessite quatre.
    \item \textbf{Bordage (rentrer les bords).} Une fois le lit fait, \textbf{rentrer soigneusement tout le pourtour de la moustiquaire sous le matelas}, tout autour du lit. C'est la condition indispensable de l'étanchéité physique. Aucun espace ne doit subsister entre le matelas et la toile.
    \item \textbf{Vérification finale.} Inspecter toute la surface : absence de trous, de déchirures, de coutures décousues. Vérifier que l'ouverture d'accès (chevauchement des pans) se referme correctement. S'asseoir sur le lit et vérifier que la bordure ne se soulève pas.
\end{enumerate}
\end{methode}

\begin{popfigure}
\begin{tikzpicture}[scale=0.85, every node/.style={font=\small}]
    % Lit (vue de face, coupe)
    % Sommier
    \draw[popdark, line width=2pt, fill=popgold!30] (0,0) rectangle (10,1);
    % Matelas
    \draw[popgold, line width=1.5pt, fill=popgold!50] (0,1) rectangle (10,3.5);
    % Oreiller
    \draw[fill=white, draw=popmuted] (0.5,3.5) rectangle (3,4.5);
    
    % Plafond (ligne pointillée)
    \draw[dashed, popmuted, thick] (-0.5, 7) -- (10.5, 7);
    \node[popmuted, above] at (5, 7) {Plafond / points d'attache (h $\approx$ 1,80 m)};
    
    % Ficelles de suspension
    \foreach \x in {0.5, 9.5} {
        \draw[thick, popdark] (\x, 7) -- (\x, 4.5);
        \draw[fill=popdark] (\x, 7) circle (3pt);
    }
    
    % Toile de la moustiquaire (forme tombante)
    \draw[popblue, line width=1.5pt, fill=popblueL, opacity=0.6] 
        (0.5, 4.5) -- (9.5, 4.5) 
        -- (10, 1) -- (10, 0) -- (0, 0) -- (0, 1) 
        -- cycle;
    
    % Bordure rentrée sous le matelas (détail)
    \draw[popblue, line width=2.5pt, decorate, decoration={snake, amplitude=1pt, segment length=4pt}] 
        (0, 1) -- (10, 1);
    \node[popblue, align=center, font=\bfseries\small] at (5, 0.4) {Bordure rentr\'ee sous le matelas};
    
    % Moustiques à l'extérieur
    \foreach \i in {1,...,5} {
        \node[popmuted, rotate=90, opacity=0.7] at (-1.2, {0.8 + \i*0.9}) {\tiny $\triangleright$};
        \node[popmuted, rotate=-90, opacity=0.7] at (11.2, {0.8 + \i*0.9}) {\tiny $\triangleleft$};
    }
    
    % Flèches d'action insecticide
    \draw[poporange, very thick, ->, >=Stealth] (10.6, 2) -- (10.1, 2);
    \node[poporange, font=\tiny, align=left, anchor=west] at (10.7, 2) {Contact $\to$ mort};
    \draw[poporange, very thick, ->, >=Stealth] (-0.6, 2) -- (-0.1, 2);
    \node[poporange, font=\tiny, align=right, anchor=east] at (-0.7, 2) {R\'epulsion};
    
    % Personne sous la moustiquaire
    \draw[fill=poppinkL, draw=poppink, line width=1pt, rounded corners=3pt] (3, 1.2) -- (4, 1.2) -- (4.2, 3) -- (3.8, 3.2) -- (2.8, 3.2) -- (3, 3) -- cycle;
    \node[align=center, font=\small] at (3.5, 2.2) {Dormeur\\prot\'eg\'e};
    
    % Légendes numérotées
    \node[draw, circle, fill=popblue, text=white, font=\bfseries\small] at (0.5, 6.5) {1};
    \node[draw, circle, fill=popblue, text=white, font=\bfseries\small] at (9.5, 6.5) {2};
    \node[draw, circle, fill=popgreen, text=white, font=\bfseries\small] at (5, 1.6) {3};
    \node[draw, circle, fill=poporange, text=white, font=\bfseries\small] at (11.6, 3.2) {4};
    
\end{tikzpicture}
\legende{Schéma d'un lit correctement équipé d'une MILD. 1 et 2 : points de suspension au plafond (ficelles tendues). 3 : bordure de la moustiquaire rentrée hermétiquement sous le matelas (barrière physique). 4 : action insecticide au contact de la toile (paralysie puis mort des anophèles). Les moustiques (symboles gris) sont tenus à l'extérieur.}
\end{popfigure}

\courssub{Interprétation : le double mécanisme de protection}

Pourquoi cette installation rigoureuse est-elle indispensable ? L'efficacité de la MILD ne repose pas sur un seul principe, mais sur la \textbf{combinaison de deux effets} que nous allons dégager.

\begin{propriete}[Double effet protecteur de la MILD]
La moustiquaire imprégnée de longue durée (MILD) protège le dormeur par un \textbf{double effet} :
\begin{enumerate}
    \item \textbf{Effet barrière (mécanique) :} la toile en polyester ou en polyéthylène, à mailles très fines, empêche physiquement l'anophèle d'atteindre la peau du dormeur, \emph{à condition} que la moustiquaire soit intacte et correctement bordée sous le matelas.
    \item \textbf{Effet insecticide (chimique) :} les fibres de la toile contiennent un insecticide de la famille des pyréthrinoïdes (souvent la \textbf{perméthrine} ou la \textbf{deltaméthrine}) qui migre lentement vers la surface. Lorsque l'anophèle se pose sur la toile (attiré par la chaleur et l'odeur du dormeur), il entre en contact avec l'insecticide : il est d'abord paralysé (il tombe, c'est l'\emph{effet knock-down}), puis il meurt. L'insecticide a aussi un \emph{effet répulsif} qui fait fuir une partie des moustiques avant la piqûre.
\end{enumerate}
La combinaison des deux effets assure une protection même si la toile présente de micro-trous (l'effet chimique agit encore) ou si l'insecticide s'épuise partiellement (l'effet mécanique subsiste).
\end{propriete}

\begin{attention}[Erreur fréquente : croire que l'insecticide seul suffit]
Ne pas border la moustiquaire sous le matelas en se disant « l'insecticide va les tuer de toute façon » est une \textbf{erreur grave}. Un anophèle qui entre par l'ouverture du bas pique \emph{avant} de se poser sur la toile imprégnée. L'effet barrière est la première ligne de défense ; l'insecticide est la seconde (et il protège aussi les autres habitants en tuant les moustiques présents dans la pièce).
\end{attention}

\begin{attention}[Erreur fréquente : usage saisonnier]
Ne pas utiliser la moustiquaire en saison sèche sous prétexte qu'« il n'y a pas de moustiques » est une \textbf{erreur}. Au Cameroun, les anophèles sont présents \textbf{toute l'année}. Leur densité diminue en saison sèche, mais la transmission du paludisme ne s'arrête jamais complètement. La protection doit être \textbf{quotidienne, 12 mois sur 12}.
\end{attention}

\courssub{Entretien et durée de vie : préserver l'efficacité chimique}

L'insecticide n'est pas éternel. Il se dégrade par les lavages répétés, par l'action des rayons ultraviolets du soleil et par la chaleur. Le protocole d'entretien est donc aussi codifié que l'installation.

\begin{methode}[Entretien d'une MILD : règles d'or]
\begin{enumerate}
    \item \textbf{Lavage :} uniquement quand la moustiquaire est visiblement sale (poussière, taches), au maximum \textbf{1 fois tous les 3 mois} (idéalement moins). Utiliser de l'\textbf{eau froide ou tiède ($< 30^\circ$C) et un savon doux neutre} (pas de détergent en poudre, pas d'eau de Javel). Frotter délicatement à la main, puis bien rincer.
    \item \textbf{Séchage :} \textbf{obligatoirement à l'ombre}, à plat ou suspendue, jamais au soleil direct. Les rayons ultraviolets dégradent très rapidement l'insecticide.
    \item \textbf{Ne pas repasser :} la chaleur du fer abîme les fibres et détruit l'insecticide de surface.
    \item \textbf{Réparation :} tout trou, même petit (de la taille d'un doigt), doit être cousu immédiatement (fil et aiguille) ou couvert d'un patch de réparation (souvent fourni). Un trou = une porte ouverte pour l'anophèle.
    \item \textbf{Remplacement :} respecter la durée de vie indiquée par le fabricant (généralement \textbf{3 ans} ou \textbf{20 lavages} maximum, selon le modèle). Au-delà, l'insecticide n'est plus garanti : la moustiquaire redevient une simple moustiquaire non imprégnée, beaucoup moins efficace.
\end{enumerate}
\end{methode}

\begin{exemplebox}[Exemple chiffré : durée de vie réelle d'une MILD au Cameroun]
Une MILD standard distribuée lors des campagnes nationales porte la mention : « Efficacité garantie 3 ans ou 20 lavages ».
\begin{itemize}
    \item Si une famille lave sa moustiquaire 1 fois par mois (fréquence excessive mais observée), les 20 lavages sont atteints en \textbf{20 mois} (environ 1 an et 8 mois) : l'insecticide est épuisé bien avant les 3 ans.
    \item Si la famille lave 1 fois par trimestre (recommandation), les 20 lavages ne seraient atteints qu'au bout de \textbf{5 ans} : c'est donc la limite des \textbf{3 ans} (vieillissement des fibres et de l'insecticide) qui sera atteinte en premier.
    \item \textbf{Conclusion :} la durée de vie effective est le \textbf{minimum} des deux seuils (temps OU lavages). Au Cameroun, avec la poussière et la chaleur, on observe souvent une perte d'efficacité insecticide vers 24 à 30 mois. D'où l'importance des campagnes de renouvellement organisées tous les 3 ans.
\end{itemize}
\end{exemplebox}

\courssub{Analyse des comportements à risque : enquête de terrain}

Au-delà du protocole idéal, les enquêtes menées dans les ménages camerounais révèlent des usages détournés qui compromettent la lutte contre le paludisme. Le tableau ci-dessous synthétise les bons et les mauvais usages observés.

\begin{popfigure}
\begin{center}
\begin{tabularx}{\linewidth}{|>{\bfseries}l|X|X|}
\hline
\rowcolor{popblueL}
\textbf{Situation / Comportement} & \textbf{Bon usage (protecteur)} & \textbf{Mauvais usage (à risque / inefficace)} \\
\hline
\textbf{Moment d'utilisation} & Toute l'année, chaque nuit (saison sèche et saison des pluies). & Uniquement en saison des pluies ; rangée en saison sèche. \\
\hline
\textbf{Installation} & Suspendue, tendue, bordée sous le matelas \emph{avant} de se coucher. & Posée simplement sur le lit (non suspendue) ; bords pendants hors du lit ; non tendue (plis). \\
\hline
\textbf{État de la toile} & Intacte, sans trous ; trous réparés immédiatement (couture ou patch). & Trouée, déchirée, coutures ouvertes ; « ça ne gêne pas, il y a l'insecticide ». \\
\hline
\textbf{Entretien (lavage)} & Savon doux, eau froide ou tiède, rinçage soigné, séchage \textbf{à l'ombre}. Fréquence $\le$ 1 fois par trimestre. & Détergent puissant, eau de Javel, brosse dure, eau chaude ; séchage \textbf{au soleil direct} ; lavage hebdomadaire. \\
\hline
\textbf{Usage détourné} & Aucun. Réservée au sommeil, sur le lit. & Utilisée comme filet de pêche, clôture de jardin ou de poulailler, rideau de porte, sac de transport. \\
\hline
\textbf{Remplacement} & Remplacée à la campagne de distribution suivante (tous les 3 ans) ou dès qu'elle est usée. & Gardée 5 à 7 ans « parce qu'elle a encore ses mailles » ; insecticide depuis longtemps disparu. \\
\hline
\end{tabularx}
\end{center}
\legende{Tableau comparatif des bons et mauvais usages de la MILD observés dans les ménages camerounais.}
\end{popfigure}

\begin{savaistu}[Impact communautaire : l'effet de masse]
Au-delà de la protection individuelle, l'utilisation massive des MILD (quand plus de 80\,\% de la population visée en dort sous une chaque nuit) crée un \textbf{effet de masse} (ou effet communautaire). Les anophèles qui entrent dans les maisons pour piquer sont tués au contact des moustiquaires, même s'ils ne piquent personne. Cela réduit le nombre de moustiques et leur durée de vie dans le village ou le quartier, protégeant \emph{indirectement} les personnes non couvertes (nourrissons, personnes âgées). C'est la justification des campagnes de distribution gratuite menées par le MINSANTE et le Programme National de Lutte contre le Paludisme (PNLP) au Cameroun tous les 3 ans.
\end{savaistu}

\courssub{Exercice résolu : analyse d'une situation familiale}

\begin{exoresolu}[Analyse de pratiques à risque]
\textbf{Énoncé :} Dans un quartier de Yaoundé, une enquête auprès de 5 ménages utilisant des MILD distribuées lors de la campagne de 2021 révèle les pratiques suivantes au mois de janvier 2024 (saison sèche) :
\begin{itemize}
    \item Ménage A : moustiquaire installée, bordée, lavée 1 fois par mois au savon en poudre, séchée au soleil.
    \item Ménage B : moustiquaire rangée dans son carton « parce qu'il n'y a pas de moustiques en ce moment ».
    \item Ménage C : moustiquaire installée mais avec un trou de 5 cm non réparé ; bordée correctement.
    \item Ménage D : moustiquaire utilisée pour clôturer le potager ; la famille dort sans protection.
    \item Ménage E : moustiquaire installée, bordée, lavée 1 fois en 2 ans (savon doux, séchage à l'ombre), trou de 1 cm réparé.
\end{itemize}
\textbf{Questions :}
\begin{enumerate}
    \item Identifier pour chaque ménage la ou les erreur(s) majeure(s) commise(s) par rapport au protocole.
    \item Préciser la conséquence probable sur l'efficacité insecticide et/ou mécanique.
    \item Quel ménage a la pratique la plus conforme ? Justifier.
    \item Calculer l'âge approximatif de ces moustiquaires en janvier 2024. Sont-elles encore dans leur durée de vie théorique ?
\end{enumerate}

\tcblower
\textbf{Solution détaillée :}
\begin{enumerate}
    \item \textbf{Analyse des erreurs :}
    \begin{itemize}
        \item \textbf{A :} erreurs d'entretien majeures : savon en poudre (détergent agressif), séchage au soleil, fréquence excessive (12 lavages par an). $\to$ Dégradation chimique accélérée de l'insecticide.
        \item \textbf{B :} erreur d'usage majeure : non-utilisation en saison sèche, fondée sur la croyance fausse que les moustiques disparaissent en saison sèche.
        \item \textbf{C :} erreur de maintenance : trou de 5 cm non réparé, c'est une porte ouverte. La barrière physique est rompue : l'insecticide tue le moustique qui se pose sur la toile, mais pas celui qui entre par le trou et pique directement.
        \item \textbf{D :} détournement total de l'usage (filet de jardin). Protection individuelle nulle et gaspillage d'un bien de santé publique.
        \item \textbf{E :} pratique quasi parfaite : lavage rare et doux, séchage à l'ombre, réparation immédiate.
    \end{itemize}
    \item \textbf{Conséquences :}
    \begin{itemize}
        \item A : perte rapide de l'effet insecticide (probablement inefficace chimiquement dès 2022--2023). Ne reste que l'effet barrière, si la toile est intacte.
        \item B : protection nulle (0\,\%) : exposition aux piqûres nocturnes toute l'année.
        \item C : protection mécanique compromise par le trou ; protection chimique intacte sur le reste de la toile.
        \item D : protection nulle pour la famille ; risque de pollution du sol et de l'eau du jardin par l'insecticide.
        \item E : double effet préservé (mécanique + chimique) : protection optimale.
    \end{itemize}
    \item \textbf{Ménage E.} C'est le seul à respecter les 4 piliers : \textbf{usage permanent} + \textbf{installation correcte (bordage)} + \textbf{entretien doux et rare} + \textbf{maintenance (réparation)}.
    \item \textbf{Âge :} campagne de 2021 $\to$ janvier 2024, soit environ \textbf{3 ans} (30 à 36 mois). Ces moustiquaires sont \textbf{à la limite supérieure} de leur durée de vie théorique (3 ans). Même le ménage E doit prévoir le remplacement lors de la prochaine campagne (2024). La moustiquaire du ménage A est probablement « chimiquement morte » depuis environ 1 an.
\end{enumerate}
\end{exoresolu}

\courssub{Synthèse du TP : les 5 commandements de la MILD}

\begin{aretenir}[Règles d'or pour une protection efficace -- À mémoriser]
\begin{enumerate}
    \item \textbf{Aérer 24 h à l'ombre} avant la première utilisation.
    \item \textbf{Suspendre, tendre, border} sous le matelas \textbf{chaque nuit, toute l'année}.
    \item \textbf{Laver peu (au plus 1 fois par trimestre), au savon doux, et sécher à l'ombre} pour préserver l'insecticide.
    \item \textbf{Réparer tout trou immédiatement} (couture ou patch).
    \item \textbf{Remplacer à 3 ans} (ou après 20 lavages) lors des campagnes officielles de distribution.
\end{enumerate}
\end{aretenir}

\begin{exercice}[Application : conception d'un message de sensibilisation]
Vous êtes agent de santé communautaire. À partir des erreurs identifiées dans l'exercice résolu (ménages A, B, C, D), rédigez \textbf{4 messages courts} (style SMS ou affiche, 160 caractères maximum chacun) en français ou dans une langue locale (Ewondo, Bassa, Fulfuldé, Pidgin\ldots) pour corriger chaque comportement à risque. Chaque message doit citer l'erreur, sa conséquence et l'action correcte.
\end{exercice}

\begin{corrige}[Exercice : conception de messages de sensibilisation]
\textbf{Propositions de messages (exemples) :}
\begin{enumerate}
    \item \textit{(Cible : ménage A)} « Laver la moustiquaire au savon en poudre et la sécher au soleil détruit l'insecticide ! Lave rarement (1 fois / 3 mois), savon doux, sèche à l'ombre. »
    \item \textit{(Cible : ménage B)} « Les moustiques piquent même en saison sèche ! Dors sous ta moustiquaire TOUTES les nuits, 12 mois sur 12. Le palu ne prend pas de vacances. »
    \item \textit{(Cible : ménage C)} « Un trou = une porte ouverte pour le moustique. Répare tout trou (couture ou patch) AVANT de dormir. Protège bien ta famille. »
    \item \textit{(Cible : ménage D)} « La moustiquaire est faite pour dormir, pas pour le jardin ! Elle protège ta famille la nuit. Remets-la sur le lit dès ce soir. »
\end{enumerate}
\end{corrige}

Ce TP vous a permis de passer de la théorie (cycle du plasmodium, rôle de l'anophèle femelle) à la pratique concrète de la prévention. Maîtriser l'usage de la MILD, c'est maîtriser le maillon le plus accessible de la chaîne de transmission du paludisme. Dans la section suivante, nous changerons d'échelle pour analyser des épidémies virales (fièvre Ebola, COVID-19) et une épidémie bactérienne transmise par l'eau (le choléra), où la prévention repose sur d'autres leviers : isolement des malades, hygiène des mains, vaccination et assainissement de l'eau de boisson.

\coursec{La fièvre Ebola et la COVID-19 : épidémies virales}

Dans la section précédente, tu as appris à utiliser correctement une moustiquaire imprégnée pour te protéger du paludisme, une \cle{endémie} transmise par un moustique. Mais toutes les maladies infectieuses ne se transmettent pas par un insecte ! Certaines passent directement d'une personne à une autre, et peuvent alors se répandre très vite dans une population : ce sont des \cle{épidémies}. Dans cette section, nous allons étudier deux épidémies virales qui ont marqué l'histoire récente, y compris celle du Cameroun : la \cle{fièvre Ebola} et la \cle{COVID-19}.

\courssub{La fièvre Ebola : une épidémie transmise par contact}

\textbf{Observation.} En 2014, une terrible épidémie se déclare en Guinée, puis s'étend au Liberia et à la Sierra Leone, en Afrique de l'Ouest. Les malades présentent une forte fièvre, des vomissements, des diarrhées et parfois des saignements. Beaucoup en meurent. Les soignants eux-mêmes tombent malades en grand nombre. Comment expliquer une propagation aussi rapide ?

\begin{definition}[La fièvre Ebola]
La \cle{fièvre Ebola} (ou maladie à virus Ebola) est une maladie infectieuse grave causée par un \cle{virus} : le \cle{virus Ebola}. Elle est apparue pour la première fois en 1976, près de la rivière Ebola (en République démocratique du Congo), d'où elle tire son nom. Elle se manifeste par des poussées épidémiques, surtout en Afrique.
\end{definition}

\begin{savaistu}[Un réservoir animal probable]
Le virus Ebola vit normalement chez des animaux sauvages, sans forcément les rendre malades : on dit que ces animaux constituent le \cle{réservoir} du virus. Les chauves-souris sont le réservoir probable du virus Ebola. L'homme se contamine au contact d'animaux infectés (chauves-souris, singes, antilopes), puis le virus passe d'homme à homme.
\end{savaistu}

\textbf{Comment le virus Ebola se transmet-il d'un homme à un autre ?}

L'observation des chaînes de contamination pendant les épidémies a montré que les personnes qui tombaient malades avaient presque toujours \emph{touché} un malade, un mort, ou des objets souillés (draps, vêtements, seringues). En revanche, les personnes qui avaient simplement partagé la même pièce, sans contact, ne tombaient pas malades. On en conclut que :

\begin{propriete}[Transmission de la fièvre Ebola]
Le virus Ebola se transmet par \cle{contact direct} avec :
\begin{itemize}
\item le \cle{sang} ou les \cle{liquides biologiques} (vomissements, diarrhées, sueur, salive, urine) d'une personne malade ou décédée ;
\item les \cle{objets souillés} par ces liquides (draps, vêtements, matériel médical non désinfecté).
\end{itemize}
Le virus pénètre dans l'organisme par une plaie de la peau ou par les muqueuses (yeux, bouche, nez).
\end{propriete}

\begin{attention}[Erreur fréquente : Ebola ne se transmet pas par l'air !]
Contrairement à une idée répandue, le virus Ebola \textbf{ne se transmet pas par l'air} comme la grippe ou la COVID-19. Il faut obligatoirement un \emph{contact} avec les liquides biologiques d'un malade ou d'un mort. C'est pourquoi les enterrements traditionnels, où la famille touche et lave le corps du défunt, ont été de grands moments de contamination pendant les épidémies.
\end{attention}

\textbf{Manifestations de la maladie.} Après une période d'incubation de 2 à 21 jours, la maladie se déclare brutalement :
\begin{itemize}
\item une \cle{fièvre brutale} et intense ;
\item une \cle{fatigue intense}, des maux de tête et des douleurs musculaires ;
\item des \cle{vomissements} et des \cle{diarrhées}, qui provoquent une forte déshydratation ;
\item dans les formes graves, des \cle{hémorragies} (saignements du nez, des gencives, ou internes).
\end{itemize}
La fièvre Ebola est une maladie à \cle{forte létalité} : une grande proportion des malades en meurent (en moyenne environ un malade sur deux, et parfois davantage lors de certaines épidémies).

\begin{exemplebox}[L'épidémie d'Ebola de 2014-2016 en chiffres]
L'épidémie d'Ebola de 2014-2016 en Afrique de l'Ouest (Guinée, Liberia, Sierra Leone) a été la plus grave de l'histoire de cette maladie : elle a infecté plus de \num{28000} personnes et causé plus de \num{11000} morts. C'est cette épidémie qui a conduit à la mise au point du premier vaccin contre Ebola.
\end{exemplebox}

\textbf{Comment lutter contre Ebola ?} Puisqu'il n'existe pas de traitement simple qui guérit la maladie, la lutte repose sur la \emph{rupture des chaînes de transmission} :
\begin{itemize}
\item l'\cle{isolement des malades} dans des centres de traitement spécialisés ;
\item le port d'\cle{équipements de protection} (combinaisons, gants, masques, lunettes) par le personnel soignant ;
\item les \cle{enterrements sécurisés}, réalisés par des équipes formées, sans contact avec le corps ;
\item le \cle{traçage des contacts} : on recherche et on surveille toutes les personnes qui ont touché un malade ;
\item la \cle{vaccination} : un vaccin efficace, le \cle{vaccin rVSV-ZEBOV}, protège les personnes exposées ; il a été utilisé à grande échelle lors des épidémies récentes en République démocratique du Congo.
\end{itemize}

\courssub{La COVID-19 : une épidémie devenue pandémie}

\textbf{Observation.} Fin 2019, une nouvelle maladie apparaît en Chine, dans la ville de Wuhan. En quelques mois, elle se répand dans le monde entier : écoles fermées, déplacements interdits, masques obligatoires... Tu as toi-même vécu cette période ! Cette maladie, c'est la COVID-19.

\begin{definition}[Pandémie]
Une \cle{pandémie} est une épidémie qui s'étend à \cle{plusieurs pays ou continents} en même temps. Exemple : la COVID-19, apparue fin 2019, est devenue une pandémie en 2020, touchant presque tous les pays du monde, dont le Cameroun.
\end{definition}

\begin{definition}[La COVID-19]
La \cle{COVID-19} est une maladie infectieuse causée par un virus de la famille des coronavirus : le \cle{SARS-CoV-2}. Apparue fin 2019, elle s'est propagée dans le monde entier et est devenue une pandémie en 2020.
\end{definition}

\textbf{Comment le SARS-CoV-2 se transmet-il ?}

L'observation des chaînes de contamination a montré que les contaminations avaient lieu surtout dans les lieux fermés et bondés, entre personnes qui se parlaient, toussaient ou chantaient à proximité, \emph{sans se toucher}. On en conclut que le virus se transmet par les voies respiratoires :

\begin{propriete}[Transmission de la COVID-19]
Le SARS-CoV-2 se transmet par les \cle{gouttelettes respiratoires} et les \cle{aérosols} (fines particules en suspension dans l'air) émis par une personne infectée lorsqu'elle \emph{tousse}, \emph{parle} ou simplement \emph{respire}. Ces particules chargées de virus pénètrent dans l'organisme par le nez, la bouche ou les yeux.
\end{propriete}

\begin{propriete}[Deux virus, deux modes de transmission différents]
Les virus Ebola et SARS-CoV-2 se transmettent \emph{différemment} : le virus Ebola se transmet par \cle{contact} avec les liquides biologiques d'un malade ou d'un mort, tandis que le SARS-CoV-2 se transmet par \cle{voie respiratoire} (gouttelettes et aérosols). C'est pourquoi les moyens de lutte ne sont pas exactement les mêmes.
\end{propriete}

\begin{popfigure}
\begin{center}
\begin{tikzpicture}[scale=0.95, every node/.style={font=\small}]
% ---- EBOLA : contact ----
\begin{scope}
\node[font=\bfseries, poppurple] at (0,3.4) {EBOLA : transmission par contact};
% malade
\draw[fill=poppinkL, draw=poppink] (0,1.6) circle (0.35);
\draw[poppink, thick] (0,1.25) -- (0,0.5);
\draw[poppink, thick] (-0.5,0.95) -- (0,1.05) -- (0.5,0.95);
\draw[poppink, thick] (0,0.5) -- (-0.35,-0.1);
\draw[poppink, thick] (0,0.5) -- (0.35,-0.1);
\node[below, poppink] at (0,-0.2) {Malade};
% liquides
\draw[fill=poppink, poppink] (0.75,0.75) circle (0.08);
\draw[fill=poppink, poppink] (0.95,0.55) circle (0.06);
\draw[fill=poppink, poppink] (0.85,0.95) circle (0.05);
\node[popmuted, font=\scriptsize, align=center] at (1.7,0.75) {sang, vomissements,\\diarrhées, sueur};
% main contact
\draw[-{Stealth}, very thick, poppink] (1.1,0.8) -- (2.6,1.0);
\node[above, poppink, font=\scriptsize] at (1.85,1.05) {contact direct};
% personne saine
\draw[fill=popgreenL, draw=popgreen] (3.1,1.6) circle (0.35);
\draw[popgreen, thick] (3.1,1.25) -- (3.1,0.5);
\draw[popgreen, thick] (2.6,0.95) -- (3.1,1.05) -- (3.6,0.95);
\draw[popgreen, thick] (3.1,0.5) -- (2.75,-0.1);
\draw[popgreen, thick] (3.1,0.5) -- (3.45,-0.1);
\node[below, popgreen] at (3.1,-0.2) {Personne saine};
\node[align=center, font=\scriptsize, popmuted] at (1.55,-0.9) {ou objets souillés,\\corps d'un défunt};
\end{scope}
% ---- COVID : gouttelettes ----
\begin{scope}[xshift=8.2cm]
\node[font=\bfseries, popblue] at (0,3.4) {COVID-19 : transmission par l'air};
% malade
\draw[fill=poppinkL, draw=poppink] (-1.1,1.6) circle (0.35);
\draw[poppink, thick] (-1.1,1.25) -- (-1.1,0.5);
\draw[poppink, thick] (-1.6,0.95) -- (-1.1,1.05) -- (-0.6,0.95);
\draw[poppink, thick] (-1.1,0.5) -- (-1.45,-0.1);
\draw[poppink, thick] (-1.1,0.5) -- (-0.75,-0.1);
\node[below, poppink] at (-1.1,-0.2) {Malade};
% gouttelettes
\foreach \x/\y/\r in {-0.55/1.55/0.07, -0.25/1.75/0.05, 0.05/1.5/0.08, 0.35/1.7/0.05, 0.6/1.45/0.06, 0.9/1.65/0.05, 1.15/1.5/0.04}
  \draw[fill=popblueL, draw=popblue] (\x,\y) circle (\r);
\draw[-{Stealth}, very thick, popblue] (-0.6,1.9) .. controls (0.3,2.5) .. (1.3,2.0);
\node[above, popblue, font=\scriptsize] at (0.35,2.45) {gouttelettes, aérosols};
\node[popmuted, font=\scriptsize, align=center] at (0.35,0.7) {toux, parole,\\respiration};
% personne saine
\draw[fill=popgreenL, draw=popgreen] (1.7,1.6) circle (0.35);
\draw[popgreen, thick] (1.7,1.25) -- (1.7,0.5);
\draw[popgreen, thick] (1.2,0.95) -- (1.7,1.05) -- (2.2,0.95);
\draw[popgreen, thick] (1.7,0.5) -- (1.35,-0.1);
\draw[popgreen, thick] (1.7,0.5) -- (2.05,-0.1);
\node[below, popgreen] at (1.7,-0.2) {Personne saine};
\end{scope}
\end{tikzpicture}
\end{center}
\legende{Comparaison des modes de transmission de deux épidémies virales. À gauche : le virus Ebola passe d'une personne à l'autre par \emph{contact direct} avec le sang, les liquides biologiques ou les objets souillés. À droite : le SARS-CoV-2 se transmet par les \emph{gouttelettes respiratoires} et les aérosols émis lors de la toux, de la parole ou de la respiration.}
\end{popfigure}

\textbf{Manifestations de la COVID-19.} Les signes les plus fréquents sont :
\begin{itemize}
\item la \cle{fièvre} ;
\item une \cle{toux sèche} ;
\item une \cle{fatigue} ;
\item la \cle{perte d'odorat} et la \cle{perte du goût} (signes très caractéristiques) ;
\item dans les \cle{formes graves} : des \cle{difficultés respiratoires} pouvant nécessiter une hospitalisation, voire entraîner la mort, surtout chez les personnes âgées ou déjà malades.
\end{itemize}

\textbf{Comment lutter contre la COVID-19 ?} La lutte repose sur :
\begin{itemize}
\item les \cle{gestes barrières} : \emph{lavage régulier des mains} au savon, \emph{port du masque}, \emph{distanciation} (garder ses distances, éviter les rassemblements), aérer les pièces ;
\item le \cle{dépistage} (tests) pour identifier les personnes infectées ;
\item l'\cle{isolement} des malades et des personnes contaminées ;
\item la \cle{vaccination}, qui protège contre les formes graves de la maladie.
\end{itemize}

\begin{savaistu}[La vaccination contre la COVID-19 au Cameroun]
Le Cameroun a organisé des campagnes de vaccination contre la COVID-19 à partir de 2021. Les vaccins ont d'abord été proposés aux personnels de santé, aux personnes âgées et aux personnes fragiles, puis à toute la population. Des centres de vaccination ont été installés dans les hôpitaux et les centres de santé de tout le pays.
\end{savaistu}

\textbf{Étude de document : l'évolution de l'épidémie de COVID-19 au Cameroun en 2020.} Le graphique ci-dessous représente l'évolution du nombre approximatif de nouveaux cas de COVID-19 détectés chaque mois au Cameroun au cours de l'année 2020.

\begin{popfigure}
\begin{center}
\begin{tikzpicture}
\begin{axis}[
  width=0.85\linewidth, height=6.2cm,
  xlabel={Mois de l'année 2020},
  ylabel={Nouveaux cas par mois (approx.)},
  xtick={1,2,3,4,5,6,7,8,9,10,11,12},
  xticklabels={Jan,Fév,Mar,Avr,Mai,Juin,Juil,Août,Sep,Oct,Nov,Déc},
  x tick label style={rotate=45, anchor=east, font=\scriptsize},
  ymin=0, ymax=7500,
  grid=major, grid style={popline!40},
  tick label style={font=\scriptsize},
  label style={font=\small},
]
\addplot[popcurve, mark=*, mark size=1.5pt] coordinates {
(1,0) (2,0) (3,140) (4,1700) (5,5400) (6,6600) (7,4300) (8,2200) (9,1300) (10,900) (11,1100) (12,1500)
};
\end{axis}
\end{tikzpicture}
\end{center}
\legende{Évolution approximative du nombre de nouveaux cas de COVID-19 détectés par mois au Cameroun en 2020. On observe une première \emph{vague épidémique} : les cas augmentent fortement à partir de mars, atteignent un \emph{pic} vers mai-juin, puis diminuent progressivement.}
\end{popfigure}

\begin{experience}[Analyse du document : lecture d'une courbe épidémique]
\begin{enumerate}
\item \textbf{Observer} : à partir de quel mois les premiers cas apparaissent-ils au Cameroun ?
\item \textbf{Repérer} : en quels mois le nombre de nouveaux cas est-il le plus élevé ?
\item \textbf{Décrire} : comment évolue la courbe après ce pic ?
\item \textbf{Interpréter} : à ton avis, quelles mesures prises par les autorités ont pu faire baisser le nombre de cas ?
\end{enumerate}
\textbf{Réponses attendues} : les premiers cas apparaissent en \emph{mars 2020} ; le pic se situe vers \emph{mai-juin} ; ensuite la courbe \emph{diminue} progressivement. Cette baisse s'explique par les mesures de lutte : fermeture temporaire des écoles, port du masque, lavage des mains, limitation des rassemblements. L'ensemble forme une \cle{vague épidémique}.
\end{experience}

\courssub{Stratégie générale face à une épidémie virale}

\begin{methode}[Face à une épidémie virale : rompre les chaînes de transmission]
Une épidémie se propage parce que chaque malade contamine d'autres personnes : c'est la \cle{chaîne de transmission}. Pour arrêter une épidémie virale, il faut \emph{casser} cette chaîne. La stratégie de base comporte trois étapes :
\begin{enumerate}
\item \textbf{Identifier} les malades (reconnaître les signes de la maladie, dépistage si possible) ;
\item \textbf{Isoler} les malades et protéger les personnes saines (isolement, équipements de protection, gestes barrières adaptés au mode de transmission) ;
\item \textbf{Vacciner} la population lorsqu'un vaccin existe, afin de rendre les personnes insensibles au virus.
\end{enumerate}
Cette stratégie s'adapte à chaque maladie : contre Ebola, on insiste sur l'isolement, la protection des soignants et les enterrements sécurisés ; contre la COVID-19, sur le masque, la distanciation et le lavage des mains.
\end{methode}

\begin{aretenir}[Point commun d'Ebola et de la COVID-19 : ce sont des maladies virales]
La fièvre Ebola et la COVID-19 sont toutes les deux causées par des \cle{virus}. Or les \cle{antibiotiques sont inefficaces contre les virus} : ils ne tuent que les bactéries. Il est donc inutile — et dangereux — de prendre des antibiotiques contre Ebola ou la COVID-19. Face à ces maladies, on soigne les signes (réhydratation, oxygène si besoin) et surtout on \emph{prévient} la transmission : isolement, gestes barrières, vaccination.
\end{aretenir}

\begin{attention}[Pièges à éviter]
\begin{itemize}
\item Ne confonds pas les modes de transmission : Ebola = \emph{contact} avec les liquides biologiques ; COVID-19 = \emph{voie respiratoire} (gouttelettes et aérosols).
\item Ne dis pas qu'un antibiotique peut guérir la COVID-19 ou Ebola : c'est faux, ce sont des virus !
\item Une \emph{épidémie} est limitée à une région ou un pays ; une \emph{pandémie} s'étend à plusieurs pays ou continents. La COVID-19 est une pandémie ; les épidémies d'Ebola sont restées des épidémies régionales.
\end{itemize}
\end{attention}

\begin{exoresolu}[Choisir les bonnes mesures de lutte]
\textbf{Énoncé.} Dans un village, deux maladies se déclarent en même temps : un cas de fièvre Ebola et plusieurs cas de COVID-19. Pour chacune des mesures suivantes, indique si elle est efficace contre Ebola, contre la COVID-19, ou contre les deux : (a) port de gants et de combinaisons par les soignants ; (b) port du masque dans les lieux publics ; (c) isolement des malades ; (d) enterrements sécurisés ; (e) prise d'antibiotiques.
\tcblower
\textbf{Solution détaillée.}
\begin{itemize}
\item (a) \emph{Gants et combinaisons} : efficace surtout contre \textbf{Ebola}, car le virus se transmet par contact avec les liquides biologiques ; le personnel soignant doit se protéger la peau et les muqueuses.
\item (b) \emph{Port du masque} : efficace surtout contre la \textbf{COVID-19}, car le virus se transmet par les gouttelettes respiratoires ; le masque bloque ces gouttelettes. Contre Ebola, il ne suffit pas.
\item (c) \emph{Isolement des malades} : efficace contre \textbf{les deux} maladies, car il empêche le malade de contaminer d'autres personnes : c'est la rupture de la chaîne de transmission.
\item (d) \emph{Enterrements sécurisés} : efficace contre \textbf{Ebola}, car le corps d'un défunt reste contaminant (contact avec les liquides biologiques).
\item (e) \emph{Antibiotiques} : efficaces contre \textbf{aucune des deux}, car Ebola et la COVID-19 sont des maladies \emph{virales}, et les antibiotiques n'agissent que sur les bactéries.
\end{itemize}
\textbf{Conclusion} : les mesures de lutte doivent être adaptées au mode de transmission de chaque maladie, mais l'isolement des malades est une mesure commune à toutes les épidémies.
\end{exoresolu}

\begin{exercice}[Pour t'entraîner]
\begin{enumerate}
\item Définis les termes : épidémie, pandémie.
\item Cite le virus responsable : (a) de la fièvre Ebola ; (b) de la COVID-19.
\item Recopie et complète le tableau suivant :
\end{enumerate}
\begin{center}
\begin{tabularx}{\linewidth}{|l|X|X|}
\hline
\rowcolor{popblueL} & \textbf{Fièvre Ebola} & \textbf{COVID-19} \\
\hline
Agent causal & \ldots & \ldots \\
\hline
Mode de transmission & \ldots & \ldots \\
\hline
Deux manifestations & \ldots & \ldots \\
\hline
Deux moyens de lutte & \ldots & \ldots \\
\hline
\end{tabularx}
\end{center}
\begin{enumerate}
\setcounter{enumi}{3}
\item Explique pourquoi il est inutile de prescrire des antibiotiques à un malade atteint de COVID-19.
\end{enumerate}
\end{exercice}

\begin{corrige}[Exercice « Pour t'entraîner »]
\begin{enumerate}
\item Une \emph{épidémie} est l'apparition rapide d'un grand nombre de cas d'une maladie dans une région ou un pays donné. Une \emph{pandémie} est une épidémie qui s'étend à plusieurs pays ou continents.
\item (a) Le virus Ebola ; (b) le coronavirus SARS-CoV-2.
\item Fièvre Ebola : virus Ebola ; contact avec le sang, les liquides biologiques ou les objets souillés ; fièvre brutale, vomissements, diarrhées, hémorragies ; isolement, équipements de protection, enterrements sécurisés, vaccination. COVID-19 : SARS-CoV-2 ; gouttelettes respiratoires et aérosols ; fièvre, toux sèche, perte d'odorat et de goût ; gestes barrières (masque, lavage des mains, distanciation), dépistage, isolement, vaccination.
\item La COVID-19 est causée par un \emph{virus}. Or les antibiotiques n'agissent que sur les \emph{bactéries} : ils sont donc inefficaces contre les virus. En prescrire serait inutile et favoriserait la résistance des bactéries aux antibiotiques.
\end{enumerate}
\end{corrige}

\coursec{Le choléra : épidémie hydrique et synthèse}

Nous venons d'étudier deux épidémies virales : la fièvre Ebola et la COVID-19, qui se transmettent par contact direct avec les liquides biologiques d'un malade ou par les gouttelettes respiratoires. Toutes les épidémies ne se transmettent cependant pas ainsi. Certaines maladies explosent en épidémies lorsque l'eau de boisson est souillée : c'est le cas du \cle{choléra}, qui frappe régulièrement certaines régions du Cameroun (Extrême-Nord, Nord, Littoral) lors des saisons des pluies. Cette dernière section nous permettra aussi de dresser une synthèse comparative des quatre maladies étudiées dans cette leçon.

\courssub{Le choléra : une maladie bactérienne à transmission féco-orale}

\begin{definition}[Maladie hydrique ou à transmission féco-orale]
Une \cle{maladie hydrique} (ou à \cle{transmission féco-orale}) est une maladie qui se transmet par l'ingestion d'\textbf{eau} ou d'\textbf{aliments souillés} par les matières fécales (selles) d'une personne malade. Le microbe passe ainsi des selles d'un malade à la bouche d'une personne saine.
\end{definition}

Le choléra est une maladie causée par une \textbf{bactérie} : le \cle{vibrion cholérique}, dont le nom scientifique est \emph{Vibrio cholerae}. C'est une petite bactérie en forme de virgule, découverte en 1883 par le médecin allemand \textbf{Robert Koch} (le même savant qui découvrit le bacille de la tuberculose). Ce fait est un savoir admis : le vibrion cholérique se multiplie dans l'intestin de l'homme et y provoque une diarrhée très abondante.

\begin{savaistu}[John Snow, père de l'épidémiologie]
En 1854, une terrible épidémie de choléra frappe le quartier de Soho à Londres. Le médecin \textbf{John Snow} a une idée géniale : il reporte sur une carte de la ville l'adresse de chaque malade. Il observe alors que les cas se regroupent autour d'une \textbf{même pompe à eau}, située dans Broad Street. Il fait fermer la pompe : l'épidémie s'arrête. Cette étude a montré, pour la première fois, le rôle de l'eau souillée dans la transmission du choléra. John Snow est aujourd'hui considéré comme le \textbf{père de l'épidémiologie moderne}, la science qui étudie la répartition des maladies dans les populations.
\end{savaistu}

\begin{propriete}[Mode de transmission du choléra]
Le choléra se transmet par l'\textbf{eau contaminée} (et les aliments souillés), comme l'a établi historiquement l'étude de John Snow à Londres en 1854. Le vibrion cholérique est éliminé en grand nombre dans les selles des malades ; si ces selles souillent l'eau de boisson, les aliments ou les mains, la bactérie peut être avalée par une personne saine qui tombe malade à son tour. C'est pourquoi le choléra est appelé la \emph{maladie des mains sales}.
\end{propriete}

\begin{experience}[La démarche de John Snow : observation, hypothèse, vérification]
\begin{enumerate}
\item \textbf{Observation} : en quelques jours, des centaines de cas de choléra apparaissent dans le quartier de Soho.
\item \textbf{Hypothèse} : John Snow suppose que la maladie se transmet par l'eau de boisson (à l'époque, on croyait que le choléra se transmettait par l'air !).
\item \textbf{Vérification} : il cartographie les cas. Résultat : la quasi-totalité des malades boit l'eau de la pompe de Broad Street. De plus, les employés d'une brasserie voisine, qui boivent de la bière et non l'eau de la pompe, ne sont \textbf{pas} malades.
\item \textbf{Conclusion} : l'eau de la pompe est la source de la contamination. La fermeture de la pompe stoppe l'épidémie : l'hypothèse est validée.
\end{enumerate}
Cette démarche (observation $\rightarrow$ hypothèse $\rightarrow$ vérification $\rightarrow$ conclusion) est la démarche d'investigation scientifique que tu dois savoir rédiger au BEPC.
\end{experience}

\begin{popfigure}
\begin{center}
\begin{tikzpicture}[font=\small, >=Stealth, node distance=1.1cm]
% Blocs principaux
\node[draw=popink, fill=poppinkL, rounded corners, thick, minimum width=3.4cm, minimum height=1.1cm, align=center] (malade) {\textbf{Homme malade}\\ (vibrions dans l'intestin)};
\node[draw=poporange, fill=poporangeL, rounded corners, thick, minimum width=3.4cm, minimum height=1.1cm, align=center, right=2.6cm of malade] (selles) {\textbf{Selles souillées}\\ (millions de vibrions)};
\node[draw=popblue, fill=popblueL, rounded corners, thick, minimum width=3.6cm, minimum height=1.1cm, align=center, below=1.5cm of selles] (eau) {\textbf{Eau / aliments souillés}\\ (puits, rivière, légumes crus)};
\node[draw=popgreen, fill=popgreenL, rounded corners, thick, minimum width=3.4cm, minimum height=1.1cm, align=center, left=2.6cm of eau] (sain) {\textbf{Homme sain}\\ (ingère les vibrions)};
% Flèches du cycle
\draw[->, very thick, popink] (malade) -- (selles);
\draw[->, very thick, poporange] (selles) -- (eau);
\draw[->, very thick, popblue] (eau) -- node[above, sloped, font=\footnotesize]{ingestion} (sain);
\draw[->, very thick, popgreen] (sain) -- (malade);
% Barrières d'hygiène
\node[draw=popgold, fill=popgoldL, rounded corners, thick, align=center, font=\footnotesize, below=0.9cm of eau, minimum width=8.6cm] (barrieres) {\textbf{Barrières d'hygiène qui rompent le cycle :} lavage des mains au savon ;\\ ébullition ou traitement de l'eau (eau de Javel) ; hygiène alimentaire ; assainissement (latrines)};
\draw[->, very thick, popgold, dashed] (barrieres.north -| eau) -- (eau.south);
\end{tikzpicture}
\end{center}
\legende{Cycle de transmission féco-orale du choléra : les vibrions éliminés dans les selles du malade souillent l'eau ou les aliments, qui contaminent un homme sain. Les mesures d'hygiène (en jaune) brisent ce cycle.}
\end{popfigure}

\courssub{Manifestations et danger du choléra}

Après une courte incubation (quelques heures à \SI{5}{jours}), la maladie se déclare brutalement :

\begin{itemize}
\item des \textbf{diarrhées très abondantes} : les selles, liquides et blanchâtres, ressemblent à de l'eau de riz ; on dit qu'elles sont \cle{riziformes} ;
\item des \textbf{vomissements} répétés ;
\item une \textbf{déshydratation rapide} : le malade perd d'énormes quantités d'eau et de sels minéraux (parfois plus de \SI{10}{L} par jour !). Sa peau se ride, ses yeux se creusent, sa tension artérielle chute.
\end{itemize}

\begin{attention}[Le vrai danger : la déshydratation]
Erreur fréquente : croire que c'est la bactérie elle-même qui tue. En réalité, le \textbf{danger majeur du choléra est la déshydratation} provoquée par les diarrhées et les vomissements : la perte d'eau et de sels minéraux perturbe le fonctionnement du cœur et des reins. La priorité absolue est donc de \textbf{réhydrater} le malade, le plus vite possible.
\end{attention}

\begin{exemplebox}[Un exemple chiffré à retenir]
Sans traitement, le choléra peut tuer en \textbf{moins de \SI{24}{h}}, et la proportion de malades qui meurent (létalité) peut dépasser \SI{50}{\%}. En revanche, avec une \textbf{réhydratation bien conduite}, la létalité tombe \textbf{sous \SI{1}{\%}}. C'est la preuve que la réhydratation est la mesure qui sauve des vies.
\end{exemplebox}

\courssub{Les moyens de lutte contre le choléra}

\textbf{1. La lutte préventive (éviter la contamination)}

\begin{itemize}
\item \textbf{Hygiène de l'eau de boisson} : faire bouillir l'eau ou la traiter avec quelques gouttes d'eau de Javel ; la conserver dans des récipients propres et couverts ;
\item \textbf{Lavage des mains au savon} après être allé aux toilettes et avant de préparer ou de prendre les repas ;
\item \textbf{Hygiène alimentaire} : bien laver et cuire les aliments, éviter les crudités et les aliments vendus à même le sol en période d'épidémie ;
\item \textbf{Assainissement du milieu} : utiliser des latrines propres, ne pas déféquer près des points d'eau, évacuer les ordures ;
\item \textbf{Vaccination} des populations en zone à risque, organisée par les autorités sanitaires.
\end{itemize}

\textbf{2. La lutte curative (soigner le malade)}

\begin{itemize}
\item la \textbf{réhydratation orale} à l'aide d'une \cle{solution de réhydratation orale} (SRO), donnée au malade en petites gorgées fréquentes ;
\item la \textbf{réhydratation intraveineuse} (perfusion) dans les cas graves, en milieu hospitalier ;
\item les \textbf{antibiotiques}, uniquement dans les cas graves et sur prescription médicale ;
\item l'isolement et la prise en charge dans des \textbf{centres de traitement du choléra} (CTC), où les selles et les vomissures sont désinfectées pour éviter la propagation.
\end{itemize}

\begin{methode}[Préparer une solution de réhydratation orale (SRO) maison]
En attendant l'arrivée au centre de santé, on peut préparer une SRO ainsi :
\begin{enumerate}
\item Prendre \textbf{\SI{1}{L} d'eau potable} (bouillie puis refroidie, ou traitée).
\item Ajouter \textbf{6 cuillères à café rases de sucre}.
\item Ajouter \textbf{1/2 cuillère à café de sel}.
\item Bien mélanger jusqu'à dissolution complète.
\item Faire boire le malade en petites quantités, très souvent (après chaque selle liquide).
\end{enumerate}
Le sucre permet à l'intestin d'absorber l'eau et le sel : c'est ce mélange simple qui sauve des vies.
\end{methode}

\courssub{Synthèse comparative des quatre maladies étudiées}

\begin{center}
\renewcommand{\arraystretch}{1.35}
\begin{tabularx}{\linewidth}{|l|X|X|X|X|}
\hline
\rowcolor{popblueL} \textbf{Critère} & \textbf{Paludisme} & \textbf{Fièvre Ebola} & \textbf{COVID-19} & \textbf{Choléra} \\
\hline
\textbf{Agent causal} & Parasite : \emph{Plasmodium} & Virus Ebola & Virus SARS-CoV-2 & Bactérie : \emph{Vibrio cholerae} \\
\hline
\textbf{Mode de transmission} & Piqûre d'un moustique (anophèle femelle) & Contact avec les liquides biologiques d'un malade & Gouttelettes respiratoires (toux, parole) & Eau et aliments souillés (voie féco-orale) \\
\hline
\textbf{Type de maladie} & Endémie & Épidémie & Épidémie (pandémie) & Épidémie \\
\hline
\textbf{Principaux moyens de lutte} & Moustiquaires imprégnées, destruction des gîtes larvaires, médicaments antipaludiques & Isolement des malades, hygiène, équipements de protection, vaccination & Gestes barrières, lavage des mains, masque, vaccination & Eau potable, lavage des mains, assainissement, réhydratation (SRO), vaccination \\
\hline
\end{tabularx}
\end{center}

\begin{exoresolu}[Identifier une épidémie et proposer des mesures de lutte]
\textbf{Énoncé.} Dans un village de l'Extrême-Nord du Cameroun, en pleine saison des pluies, plusieurs personnes présentent en quelques jours des diarrhées abondantes et des vomissements. Deux malades, non soignés, meurent en moins de \SI{24}{h}. L'enquête montre que tous les malades puisent leur eau dans le même puits, situé près de latrines mal construites.
1. Quelle est la maladie la plus probable ? Justifie.
2. Quel est l'agent causal de cette maladie ?
3. Explique comment le puits a pu être contaminé.
4. Cite trois mesures pour stopper cette épidémie.
\tcblower
\textbf{Solution détaillée.}
1. Il s'agit très probablement du \textbf{choléra} : l'apparition brutale de nombreux cas groupés (épidémie), les diarrhées abondantes avec vomissements, la mort rapide par déshydratation et le lien avec une source d'eau commune sont des indices caractéristiques.
2. L'agent causal est une \textbf{bactérie}, le vibrion cholérique (\emph{Vibrio cholerae}).
3. Les latrines mal construites, situées près du puits, ont laissé les selles des malades (riches en vibrions) souiller l'eau du puits, surtout avec les pluies. En buvant cette eau contaminée, les villageois ont ingéré la bactérie : c'est la transmission féco-orale.
4. Trois mesures de lutte : \textbf{traiter l'eau} du puits (ébullition ou eau de Javel) ou interdire provisoirement son usage ; \textbf{réhydrater} rapidement les malades avec des SRO et les conduire dans un centre de traitement ; \textbf{laver les mains au savon} et améliorer l'assainissement (latrines éloignées des points d'eau). On peut ajouter la vaccination des habitants.
\end{exoresolu}

\begin{exercice}[Pour t'entraîner]
Recopie et complète le tableau suivant avec les mots : virus, bactérie, parasite, moustique, eau souillée, gouttelettes respiratoires.
\begin{center}
\begin{tabular}{|l|c|c|}
\hline
\rowcolor{popgreenL} \textbf{Maladie} & \textbf{Agent} & \textbf{Transmission} \\
\hline
Paludisme & \ldots & \ldots \\
\hline
COVID-19 & \ldots & \ldots \\
\hline
Choléra & \ldots & \ldots \\
\hline
\end{tabular}
\end{center}
\end{exercice}

\begin{corrige}[Exercice 1]
Paludisme : parasite (\emph{Plasmodium}) ; transmission par piqûre de moustique (anophèle femelle). COVID-19 : virus ; transmission par les gouttelettes respiratoires. Choléra : bactérie (vibrion cholérique) ; transmission par l'eau souillée (et les aliments contaminés).
\end{corrige}

\courssub{Conclusion générale de la leçon}

\begin{aretenir}[L'essentiel de la leçon]
\begin{itemize}
\item Une \cle{endémie} (ex. le paludisme) sévit en permanence dans une région ; une \cle{épidémie} (ex. Ebola, COVID-19, choléra) est l'apparition brutale de nombreux cas en peu de temps.
\item Le choléra est dû à une bactérie, le vibrion cholérique ; il se transmet par l'eau et les aliments souillés (voie féco-orale) et tue par déshydratation : la réhydratation (SRO) est le traitement prioritaire.
\item La lutte contre les endémies et les épidémies repose toujours sur les mêmes piliers : la \textbf{connaissance de l'agent causal et du mode de transmission}, l'\textbf{hygiène}, la \textbf{protection individuelle et collective} (moustiquaires, gestes barrières, eau potable) et la \textbf{vaccination} quand elle existe.
\end{itemize}
Retiens bien : connaître le mode de transmission d'une maladie, c'est détenir la clé pour l'empêcher de se propager. C'est cette démarche scientifique, inaugurée par John Snow, qui sauve chaque année des millions de vies, y compris ici au Cameroun.
\end{aretenir}

\newpage

\coursec{Activité d'intégration}

\courssub{Activité d'intégration : Enquête épidémiologique sur une épidémie de choléra (Cas de Nkol-Eton)}

\courssub{Objectifs de l'activité}

À l'issue de cette activité d'intégration, tu seras capable de :
\begin{itemize}
\item mobiliser les notions d'\cle{agent causal}, de \cle{mode de transmission} et de \cle{moyens de lutte} pour analyser une situation épidémique réelle ;
\item exploiter des documents variés (carte, tableau de données, interviews, texte historique) et en croiser les informations ;
\item tracer et interpréter une \cle{courbe épidémique} (nombre de cas en fonction du temps) ;
\item formuler une \cle{hypothèse} scientifique et la confronter aux faits ;
\item rédiger et justifier une démarche d'\cle{investigation épidémiologique} et proposer un \cle{plan d'action} argumenté ;
\item travailler en groupe, présenter et défendre une conclusion devant la classe.
\end{itemize}

\courssub{Situation : Une épidémie de choléra frappe le quartier Nkol-Eton (Yaoundé)}

Nous sommes en juillet, en pleine saison des pluies. Le Centre de Santé Intégré (CSI) du quartier Nkol-Eton, à Yaoundé, reçoit depuis huit semaines des malades présentant tous les mêmes symptômes : \textbf{diarrhées abondantes soudaines (selles liquides « en eau de riz »), vomissements, crampes abdominales et déshydratation rapide}. Le médecin du CSI suspecte une épidémie de \cle{choléra} et alerte la Délégation régionale de la Santé publique. Tu fais partie de l'équipe d'enquêteurs épidémiologues dépêchée sur place. Ta mission : \textbf{identifier la source de la contamination et proposer un plan d'action pour stopper l'épidémie.}

\textbf{Document 1 : Carte simplifiée du quartier Nkol-Eton.}
Le quartier compte quatre points d'eau utilisés par les habitants :
\begin{itemize}
\item un \textbf{puits P1} (au nord, près du marché) ;
\item un \textbf{forage F2} (au centre, rue des Écoles) ;
\item une \textbf{borne-fontaine B3} (à l'est, quartier administratif) ;
\item un \textbf{puits P4} (au sud, près du marécage où s'écoulent les eaux usées en saison des pluies).
\end{itemize}
La localisation des 47 cas recensés est la suivante :
\begin{itemize}
\item 38 cas habitent dans un rayon d'environ \SI{200}{m} autour du puits P4 ;
\item 5 cas autour du puits P1 ;
\item 3 cas autour du forage F2 ;
\item 1 cas près de la borne-fontaine B3.
\end{itemize}

\textbf{Document 2 : Nombre de nouveaux cas de choléra déclarés au CSI par semaine.}

\begin{center}
\begin{tabular}{|c|*{8}{c}|}\hline
\rowcolor{popblueL} \textbf{Semaine} & \textbf{S1} & \textbf{S2} & \textbf{S3} & \textbf{S4} & \textbf{S5} & \textbf{S6} & \textbf{S7} & \textbf{S8} \\ \hline
\textbf{Nouveaux cas} & 2 & 3 & 5 & 9 & 14 & 10 & 5 & 1 \\ \hline
\end{tabular}
\end{center}

\emph{Précision importante :} à la fin de la semaine S5, les autorités sanitaires ont \textbf{fermé le puits P4} et distribué de l'eau de Javel et du soluté de réhydratation orale (SRO) dans le quartier.

\textbf{Document 3 : Extraits d'interviews d'habitants.}
\begin{itemize}
\item \textbf{Mme Abena} (riveraine du puits P4) : « Nous puisons toute notre eau de boisson et de cuisine au puits P4. Quand il pleut beaucoup, l'eau du marécage monte et coule vers le puits. »
\item \textbf{M. Etoundi} (malade guéri) : « Je n'ai pas de latrines à la maison. En saison des pluies, tout est emporté par l'eau de ruissellement. »
\item \textbf{Mme Ngo Bell} (riveraine de la borne-fontaine B3) : « Nous, nous utilisons l'eau de la borne-fontaine, et à la maison nous faisons bouillir l'eau de boisson. Personne n'est tombé malade dans ma famille. »
\end{itemize}

\textbf{Document 4 : Rappel historique — John Snow et le choléra à Londres (1854).}
En 1854, à Londres, le médecin John Snow fait face à une épidémie de choléra. En reportant sur une carte l'adresse de chaque malade, il observe que les cas se concentrent autour d'une même pompe à eau publique (Broad Street). Il fait retirer la poignée de la pompe : l'épidémie s'arrête. Il venait de démontrer, sans connaître le microbe, que le choléra se transmet par l'\textbf{eau contaminée}.

\courssub{Consignes}

Travail en groupes de 4 à 5 élèves. Chaque groupe rédige un rapport d'enquête structuré qui sera présenté à la classe.

\begin{enumerate}
\item Rappelle la définition d'une \cle{épidémie} et explique pourquoi la situation de Nkol-Eton correspond à une épidémie et non à une \cle{endémie}.
\item Quel est l'\cle{agent causal} du choléra ? À quel groupe de microorganismes appartient-il ?
\item À partir du document 2, trace la \cle{courbe épidémique} : nombre de nouveaux cas (ordonnée) en fonction des semaines S1 à S8 (abscisse). Décris l'évolution de l'épidémie (montée, pic, descente) en citant des valeurs chiffrées.
\item En croisant le document 1 (carte) et le document 3 (interviews), identifie le \textbf{point d'eau suspect}. Justifie ta réponse par au moins deux arguments (chiffrés ou tirés des interviews).
\item Formule une \cle{hypothèse} sur le \cle{mode de transmission} du choléra dans ce quartier. Quel événement climatique a favorisé la contamination du point d'eau suspect ?
\item Explique pourquoi le nombre de cas diminue à partir de la semaine S6. Quel lien fais-tu avec l'expérience historique de John Snow (document 4) ?
\item Rédige un \textbf{plan de lutte} en deux volets : (a) les \textbf{mesures immédiates} pour stopper l'épidémie ; (b) les \textbf{mesures durables} pour empêcher son retour. Justifie chaque mesure par la notion scientifique qu'elle mobilise.
\item Présente ton rapport à la classe et défends ta démarche face aux questions des autres groupes.
\end{enumerate}

\begin{methode}[Tracer et interpréter une courbe épidémique]
\begin{enumerate}
\item Choisis les axes : abscisse = temps (semaines), ordonnée = nombre de nouveaux cas.
\item Repère les échelles : graduations régulières, unités indiquées.
\item Trace les points (semaine, effectif) et relie-les par des segments de droite.
\item Interprète : phase d'ascension (propagation), pic (maximum), phase de descente (contrôle), éventuel rebond.
\item Repère les événements marquants (mesures de lutte) par une ligne verticale pointillée annotée.
\end{enumerate}
\end{methode}

\courssub{Ressources à mobiliser}

\begin{aretenir}[Notions utiles de la leçon]
\begin{itemize}
\item Une \cle{épidémie} est l'apparition soudaine, localisée dans le temps, d'un nombre de cas d'une maladie supérieur à ce qui est habituellement attendu dans une population donnée. Une \cle{endémie} est la présence permanente et habituelle d'une maladie dans une région (ex. : le paludisme au Cameroun).
\item Le choléra est causé par une \cle{bactérie} (le vibrion cholérique, \emph{Vibrio cholerae}), transmise par l'\textbf{eau ou les aliments souillés} par les matières fécales de malades ou de porteurs sains : c'est une \textbf{épidémie hydrique} à transmission \textbf{féco-orale}.
\item Moyens de lutte contre le choléra : boire une eau saine (bouillie ou javellisée), se laver les mains au savon (surtout après les selles et avant les repas), utiliser des latrines hygiéniques, réhydrater les malades au \cle{SRO} (soluté de réhydratation orale), isoler et traiter les cas (antibiotiques si nécessaire).
\item Démarche d'investigation épidémiologique : \textbf{observer} (compter, décrire) $\rightarrow$ \textbf{localiser} (carte des cas) $\rightarrow$ \textbf{hypothétiser} (source, mode de transmission) $\rightarrow$ \textbf{vérifier} (croiser documents, analyses de laboratoire) $\rightarrow$ \textbf{agir} (couper la transmission, soigner) $\rightarrow$ \textbf{évaluer} (suivi de la courbe) $\rightarrow$ \textbf{prévenir} (mesures durables).
\end{itemize}
\end{aretenir}

\begin{savaistu}[John Snow, père de l'épidémiologie moderne]
En 1854, John Snow n'avait ni microscope ni connaissance de la bactérie. Par la seule puissance de la \textbf{cartographie des cas} et du \textbf{raisonnement déductif}, il a identifié la pompe de Broad Street comme source. Son action (retirer la poignée) a stoppé l'épidémie, prouvant que l'épidémiologie descriptive permet d'agir efficacement avant même de connaître l'agent causal. Aujourd'hui encore, au Cameroun (régions de l'Extrême-Nord, du Nord, du Littoral), les équipes du Ministère de la Santé publique utilisent cette même méthode « shoe-leather epidemiology » (épidémiologie de terrain) pour circonscrire les foyers de choléra.
\end{savaistu}

\courssub{Démarche et corrigé commenté}

\begin{exoresolu}[Rapport d'enquête épidémiologique de Nkol-Eton]
Reprends les huit consignes et rédige le rapport complet de l'équipe d'enquêteurs.
\tcblower
\textbf{1. Épidémie ou endémie ?}
Une épidémie se définit par l'apparition \textbf{brutale} et \textbf{localisée dans le temps} d'un nombre anormalement élevé de cas d'une maladie dans une population donnée. À Nkol-Eton, 47 cas de choléra sont déclarés en seulement 8 semaines dans un quartier unique, alors que le choléra n'y est pas habituel : c'est bien une \textbf{épidémie}. Ce n'est pas une endémie car le choléra n'est pas présent en permanence dans ce quartier (contrairement au paludisme, endémique au Cameroun, qui sévit toute l'année avec des pics saisonniers).

\textbf{2. Agent causal.}
Le choléra est causé par une \textbf{bactérie} (procaryote), le \emph{Vibrio cholerae} (vibrion cholérique), bacille Gram-négatif à mobilité dartante. C'est donc une maladie infectieuse bactérienne, à distinguer des épidémies virales comme Ebola (filovirus) ou la COVID-19 (coronavirus).

\textbf{3. Courbe épidémique.}

\begin{popfigure}
\begin{tikzpicture}
\begin{axis}[
width=\linewidth, height=7cm,
xlabel={Semaines}, ylabel={Nombre de nouveaux cas},
xmin=0.5, xmax=8.5, ymin=0, ymax=16,
xtick={1,2,3,4,5,6,7,8},
xticklabels={S1,S2,S3,S4,S5,S6,S7,S8},
ytick={0,2,4,6,8,10,12,14},
grid=major, grid style={popline!50},
tick label style={font=\small},
label style={font=\small},
legend style={at={(0.02,0.98)}, anchor=north west, font=\small, fill=white, fill opacity=0.8},
]
\addplot[popcurve, mark=*, mark options={fill=poporange, scale=0.8}] coordinates {(1,2) (2,3) (3,5) (4,9) (5,14) (6,10) (7,5) (8,1)};
\addlegendentry{Nouveaux cas}
\draw[dashed, popink, thick] (axis cs:5.5,0) -- (axis cs:5.5,14);
\node[popink, anchor=west, font=\small, align=left] at (axis cs:5.6,13) {Fin S5 : fermeture\\du puits P4};
\end{axis}
\end{tikzpicture}
\legende{Courbe épidémique hebdomadaire du choléra à Nkol-Eton. Le trait vertical pointillé rose marque la fermeture du puits P4 à la fin de la semaine S5.}
\end{popfigure}

\textbf{Interprétation :} La courbe montre une \textbf{phase d'ascension} régulière de S1 (2 cas) à S5 où elle atteint un \textbf{pic épidémique à 14 cas}. Elle entre ensuite en \textbf{phase de descente} nette : 10 cas en S6, 5 en S7, 1 seul cas en S8. L'épidémie est en voie d'extinction.

\textbf{4. Point d'eau suspect.}
C'est le \textbf{puits P4}. Deux arguments convergents :
\begin{itemize}
\item \textbf{Argument cartographique (quantitatif) :} 38 cas sur 47, soit \textbf{81 \%} des malades, habitent dans un rayon de \SI{200}{m} autour du puits P4. La répartition est très inégale : P1 (5 cas), F2 (3 cas), B3 (1 cas). Cette concentration spatiale forte oriente vers une source commune localisée.
\item \textbf{Argument des interviews (qualitatif) :} Mme Abena confirme l'usage exclusif du puits P4 pour l'eau de boisson et signale un risque de contamination par le marécage (eaux usées) en saison des pluies. À l'inverse, la famille de Mme Ngo Bell, desservie par la borne-fontaine B3 (eau traitée) et pratiquant l'ébullition, compte \textbf{zéro cas}. Cela renforce l'hypothèse d'une transmission par l'eau du puits P4.
\end{itemize}

\textbf{5. Hypothèse sur le mode de transmission.}
\textbf{Hypothèse :} Le choléra se transmet ici par la \textbf{consommation de l'eau du puits P4 contaminée par des matières fécales}.
\textbf{Mécanisme :} En saison des pluies (juillet), le ruissellement intense entraîne les matières fécales (M. Etoundi n'a pas de latrines, défécation à l'air libre) vers le marécage situé en amont du puits P4. L'eau du marécage, chargée en vibrions cholériques, s'infiltre ou déborde dans le puits P4. Les habitants boivent cette eau non traitée : la bactérie passe ainsi des selles des malades/porteurs à la bouche des nouveaux hôtes (\textbf{transmission féco-orale hydrique}).

\textbf{6. Explication de la baisse des cas (à partir de S6).}
La chute brutale des cas (14 $\rightarrow$ 10 $\rightarrow$ 5 $\rightarrow$ 1) coïncide exactement avec la \textbf{fermeture du puits P4 à la fin de la S5} et la distribution d'eau de Javel (désinfection de l'eau de boisson) et de SRO (prise en charge). En supprimant l'accès à la source contaminée, on \textbf{coupe la voie de transmission} : les nouveaux contacts avec le vibrion cessent, donc les nouveaux cas disparaissent après la période d'incubation (quelques heures à 5 jours).
\textbf{Lien avec John Snow :} C'est la réplique exacte de l'expérience de 1854. John Snow a stoppé l'épidémie de Londres en \textbf{retirant la poignée de la pompe de Broad Street} (source contaminée). À Nkol-Eton, fermer le puits P4 a eu le même effet : \textbf{couper la voie de transmission stoppe l'épidémie}, confirmant l'hypothèse hydrique.

\textbf{7. Plan de lutte.}

\begin{center}
\begin{tabularx}{\linewidth}{|l|X|X|}\hline
\rowcolor{popblueL} \textbf{Volet} & \textbf{Mesures} & \textbf{Justification scientifique (Notion mobilisée)} \\ \hline
\textbf{(a) Mesures immédiates}\\\textbf{(Urgence : couper la transmission, soigner)} &
\begin{itemize}
\item Fermer et condamner définitivement le puits P4 (remblaiement).
\item Distribuer de l'eau de Javel (hypochlorite de sodium) pour la désinfection de l'eau de boisson à domicile (1 goutte/L, attendre 30 min).
\item Distribuer massivement du \textbf{SRO} (Soluté de Réhydratation Orale) aux familles et au CSI.
\item Prise en charge médicale des cas au CSI (réhydratation IV si choc, antibiotiques : doxycycline/azithromycine pour réduire portage et durée).
\item Sensibilisation d'urgence : lavage des mains au savon (après selles, avant manger), ébullition ou javellisation de l'eau, cuisson complète des aliments.
\end{itemize} &
\begin{itemize}
\item \textbf{Supprimer la source} (réservoir + véhicule) : élimine le contact eau/bactérie/homme.
\item \textbf{Désinfection chimique} : l'hypochlorite tue \emph{V. cholerae} dans l'eau (action oxydante).
\item \textbf{Lutter contre la déshydratation} : cause majeure de décès (perte $> \SI{10}{L}$/jour d'eau + électrolytes). SRO = glucose + NaCl + KCl + citrate (cotransport Na+/glucose intact).
\item \textbf{Réduire le réservoir humain} : antibiotiques diminuent l'excrétion fécale du vibrion.
\item \textbf{Barrières hygiéniques} : savon élimine bactéries mains ; chaleur (ébullition) tue vibrion.
\end{itemize} \\ \hline
\textbf{(b) Mesures durables}\\\textbf{(Prévention : éviter le retour)} &
\begin{itemize}
\item Construire des latrines familiales améliorées (fosse ventilée, dalle lavable) \textbf{éloignées des points d'eau} ($> \SI{30}{m}$, en aval hydrogéologique).
\item Réaliser un réseau d'assainissement (caniveaux, fossés) pour évacuer les eaux usées et pluviales loin des habitations et puits.
\item Aménager/réhabiliter des points d'eau protégés : forages équipés de pompes à motricité humaine (PMH), bornes-fontaines raccordées au réseau traité (CAMWATER).
\item Éducation à l'hygiène pérenne (écoles, chefs de quartier, radios communautaires) : « Eau propre, Mains propres, Latrines propres ».
\item Surveillance épidémiologique active : déclaration obligatoire, veille syndromique au CSI.
\end{itemize} &
\begin{itemize}
\item \textbf{Barrière fécale} : empêche la contamination fécale de l'environnement (sol, eau de surface, nappe).
\item \textbf{Assainissement} : coupe la voie « ruissellement $\rightarrow$ marécage $\rightarrow$ puits ».
\item \textbf{Eau saine à la source} : élimine le besoin de traiter l'eau au quotidien (fiabilité accrue).
\item \textbf{Changement de comportement} : ancrage des gestes barrières (savoir-être).
\item \textbf{Détection précoce} : permet d'agir au premier cas (éviter amplification).
\end{itemize} \\ \hline
\end{tabularx}
\end{center}

\textbf{8. Présentation orale.}
Chaque groupe expose en 5 minutes : courbe épidémique commentée, carte annotée (foyer P4, flèches ruissellement), hypothèse argumentée, plan de lutte justifié. L'évaluation porte sur la rigueur du raisonnement (lien preuve $\rightarrow$ conclusion), la justesse du vocabulaire (agent causal, transmission féco-orale, SRO, endémie/épidémie) et la faisabilité des mesures proposées dans le contexte camerounais.

\end{exoresolu}

\courssub{Bilan de l'activité}

Cette enquête t'a permis de mettre en pratique la \cle{démarche d'investigation épidémiologique} complète, celle-là même que déploient les équipes du Ministère de la Santé publique (MINSANTE) et de l'OMS lors des épidémies de choléra qui touchent périodiquement le Cameroun (notamment dans les régions de l'Extrême-Nord, du Nord, du Littoral et du Centre) :

\begin{enumerate}
\item \textbf{Observer et quantifier} : tableau de déclaration hebdomadaire $\rightarrow$ courbe épidémique (temps, amplitude, pic).
\item \textbf{Localiser spatialement} : carte de répartition des cas $\rightarrow$ identification d'un foyer (cluster) autour du puits P4.
\item \textbf{Croiser les sources} : carte + interviews + données climatiques $\rightarrow$ formulation d'une hypothèse étayée (contamination fécale du puits par ruissellement pluvial).
\item \textbf{Agir sur la voie de transmission} : fermeture du puits, javellisation, SRO $\rightarrow$ vérification de l'efficacité par la chute de la courbe (évaluation).
\item \textbf{Prévenir la récidive} : mesures structurelles (assainissement, eau potable, latrines, éducation) $\rightarrow$ passage du curatif au préventif durable.
\end{enumerate}

\begin{aretenir}[L'essentiel à retenir pour le BEPC]
\begin{itemize}
\item Une \cle{épidémie} est une survenue brutale et temporaire ; une \cle{endémie} est une présence constante.
\item Le \cle{choléra} est une \textbf{épidémie hydrique} à transmission \textbf{féco-orale}, causée par la bactérie \emph{Vibrio cholerae}.
\item L'\cle{enquête épidémiologique} (courbe temporelle + carte spatiale + indices environnementaux) permet d'identifier la source et le mode de transmission \textbf{sans attendre l'analyse de laboratoire}.
\item La \cle{lutte} repose sur : (1) \textbf{Couper la transmission} (eau saine, hygiène, assainissement), (2) \textbf{Soigner} (SRO = priorité vitale, antibiotiques), (3) \textbf{Prévenir} (infrastructures durables, éducation).
\item \textbf{John Snow (1854)} a fondé l'épidémiologie moderne : une carte et du bon sens peuvent sauver des vies avant de connaître le microbe.
\end{itemize}
\end{aretenir}

\newpage

\coursec{Méthodes et exercices résolus}

\coursec{Quelques exemples d'endémies et d'épidémies}

\begin{definition}[Endémie, épidémie, pandémie]
On distingue trois niveaux de propagation d'une maladie transmissible :
\begin{itemize}
\item \cle{Endémie} : présence \textbf{permanente} et habituelle d'une maladie dans une région donnée (ex. : le paludisme en zone tropicale).
\item \cle{Épidémie} : apparition \textbf{brutale} et \textbf{rapide} d'un nombre de cas anormalement élevé dans une région et sur une période données (ex. : choléra, fièvre Ebola).
\item \cle{Pandémie} : épidémie qui s'étend à \textbf{plusieurs pays ou continents}, touchant une large population mondiale (ex. : COVID-19, grippe de 1918).
\end{itemize}
\end{definition}

\courssub{Méthode 1 : Distinguer une endémie d'une épidémie}

\begin{methode}[Reconnaître endémie, épidémie et pandémie]
Pour qualifier une situation sanitaire décrite dans un énoncé :
\begin{enumerate}
\item \textbf{Repérer la fréquence de la maladie} : est-elle présente en permanence dans la région, ou apparaît-elle brutalement ?
\item \textbf{Repérer l'étendue géographique} : la maladie touche-t-elle une région limitée, un pays entier, ou plusieurs continents ?
\item \textbf{Conclure avec le vocabulaire exact} :
\begin{itemize}
\item \cle{Endémie} : maladie présente de façon \textbf{permanente} (ou très fréquente) dans une région donnée (ex. : le paludisme en zone tropicale).
\item \cle{Épidémie} : apparition \textbf{brutale} et \textbf{rapide} d'un grand nombre de cas dans une région et un temps donnés (ex. : choléra, Ebola).
\item \cle{Pandémie} : épidémie qui s'étend à \textbf{plusieurs pays ou continents} (ex. : COVID-19).
\end{itemize}
\item \textbf{Justifier la réponse} en citant les indices de l'énoncé (durée, nombre de cas, zone touchée).
\end{enumerate}
\end{methode}

\begin{exoresolu}[Endémie ou épidémie ?]
\textbf{Énoncé.} Voici trois situations sanitaires observées au Cameroun et dans le monde :
\begin{itemize}
\item Situation A : dans la région du Centre, des cas de paludisme sont enregistrés chaque mois de l'année, pendant de nombreuses années.
\item Situation B : en 2022, plusieurs centaines de cas de choléra apparaissent en quelques semaines dans la région du Littoral.
\item Situation C : à partir de 2020, la COVID-19 touche presque tous les pays du monde en quelques mois.
\end{itemize}
Pour chaque situation, indique s'il s'agit d'une endémie, d'une épidémie ou d'une pandémie, en justifiant ta réponse.
\tcblower
\textbf{Solution détaillée.}
\begin{enumerate}
\item \textbf{Situation A} : le paludisme est présent \emph{en permanence} (chaque mois, chaque année) dans la région du Centre. Il s'agit donc d'une \textbf{endémie}. Le paludisme est d'ailleurs l'exemple type d'endémie en zone tropicale.
\item \textbf{Situation B} : les cas de choléra apparaissent \emph{brutalement} et en \emph{grand nombre} (plusieurs centaines) en \emph{quelques semaines}, dans une région limitée (le Littoral). Il s'agit d'une \textbf{épidémie}.
\item \textbf{Situation C} : la COVID-19 s'est propagée en quelques mois à \emph{presque tous les pays du monde}, donc à plusieurs continents. Une épidémie de cette ampleur est une \textbf{pandémie}.
\end{enumerate}
\textbf{Conclusion.} A est une endémie (présence permanente), B est une épidémie (apparition brutale et localisée), C est une pandémie (extension mondiale).
\end{exoresolu}

\courssub{Méthode 2 : Décrire l'agent causal et le cycle de transmission du paludisme}

\begin{methode}[Expliquer la transmission du paludisme]
Pour répondre à une question sur les causes et la transmission du paludisme :
\begin{enumerate}
\item \textbf{Nommer l'agent causal} : le paludisme est dû à un \cle{parasite} microscopique, le \emph{Plasmodium} (ce n'est ni un virus, ni une bactérie ; c'est un protiste).
\item \textbf{Identifier le vecteur} : la transmission se fait par la piqûre d'un moustique femelle, l'\cle{anophèle}, qui pique surtout la nuit (entre 22 h et 5 h).
\item \textbf{Décrire la chaîne de transmission} : anophèle infecté pique un homme sain $\rightarrow$ le parasite (sporozoïtes) pénètre dans le sang $\rightarrow$ multiplication dans le \textbf{foie} (phase hépatique) puis dans les \textbf{globules rouges} (phase érythrocytaire) $\rightarrow$ destruction des globules rouges (hémolyse) $\rightarrow$ fièvre et anémie $\rightarrow$ un autre anophèle aspire le sang infecté (gamétocytes) et peut contaminer une nouvelle personne.
\item \textbf{Préciser les conditions favorables} : eaux stagnantes (mares, flaques, récipients abandonnés, empreintes de pas) où les moustiques pondent (gîtes larvaires), climat chaud et humide.
\end{enumerate}
\textbf{Réflexe rédaction} : toujours distinguer l'\emph{agent} (Plasmodium), le \emph{vecteur} (anophèle femelle) et l'\emph{hôte} (l'homme).
\end{methode}

\begin{savaistu}[Le paludisme au Cameroun]
Au Cameroun, le paludisme est la première cause de consultation et d'hospitalisation. Le \emph{Programme National de Lutte contre le Paludisme (PNLP)} coordonne la distribution gratuite de moustiquaires imprégnées à longue durée d'action (MILD) lors de campagnes de masse (tous les 3 ans) et la prise en charge gratuite des enfants de moins de 5 ans et des femmes enceintes. Les espèces principales sont \emph{Plasmodium falciparum} (le plus dangereux) et \emph{P. malariae}, \emph{P. ovale}.
\end{savaistu}

\begin{exoresolu}[La chaîne de transmission du paludisme]
\textbf{Énoncé.} Dans un village de la région de l'Est, de nombreux enfants souffrent de paludisme pendant la saison des pluies. On observe autour du village de nombreuses flaques d'eau stagnante.
\begin{enumerate}
\item Quel est l'agent causal du paludisme ? Est-ce un virus ou une bactérie ?
\item Explique, en une chaîne logique, comment un enfant sain du village peut être contaminé.
\item Pourquoi les cas augmentent-ils pendant la saison des pluies ?
\end{enumerate}
\tcblower
\textbf{Solution détaillée.}
\begin{enumerate}
\item L'agent causal du paludisme est un \textbf{parasite unicellulaire} (protiste) appelé \emph{Plasmodium}. Ce n'est \textbf{ni un virus ni une bactérie} : c'est un microorganisme parasite qui vit dans le sang de l'homme.
\item Chaîne de contamination : un \textbf{anophèle femelle infecté} (moustique ayant préalablement piqué une personne malade) pique l'enfant sain, généralement \textbf{la nuit}. Lors de la piqûre, elle injecte le \emph{Plasmodium} (stade sporozoïte) dans le sang de l'enfant. Le parasite se multiplie d'abord dans le \textbf{foie} (phase silencieuse), puis envahit les \textbf{globules rouges} qu'il détruit (phase clinique) : l'enfant devient malade. Si un autre anophèle pique cet enfant malade, il aspire le parasite (stade gamétocyte) et pourra contaminer à son tour d'autres personnes.
\item Pendant la saison des pluies, les \textbf{eaux stagnantes} (flaques, mares, ornières) se multiplient : ce sont des \textbf{gîtes larvaires} où les anophèles pondent et où se développent leurs larves. Le nombre de moustiques vecteurs augmente donc fortement, ce qui multiplie les contacts homme-moustique et les cas de paludisme.
\end{enumerate}
\textbf{Conclusion.} Le paludisme est transmis par la piqûre de l'anophèle femelle infecté ; la saison des pluies favorise la reproduction des moustiques, d'où l'augmentation des cas.
\end{exoresolu}

\courssub{Méthode 3 : Citer les manifestations du paludisme et les moyens de lutte}

\begin{methode}[Manifestations et lutte contre le paludisme]
Pour construire une réponse complète :
\begin{enumerate}
\item \textbf{Manifestations (signes cliniques)} : \cle{fièvre élevée} (accès palustre souvent cyclique), frissons, maux de tête (céphalées), fatigue intense, vomissements, douleurs musculaires et articulaires ; dans les cas graves (paludisme sévère), anémie sévère (destruction massive des globules rouges), détresse respiratoire, troubles neurologiques (paludisme cérébral) pouvant entraîner la mort, surtout chez les jeunes enfants ($< 5$ ans) et les femmes enceintes.
\item \textbf{Moyens de lutte préventive} (éviter les piqûres et réduire la densité vectorielle) :
\begin{itemize}
\item dormir sous \textbf{moustiquaire imprégnée d'insecticide à longue durée d'action (MILD)}, bien bordée ;
\item détruire les \textbf{gîtes larvaires} : éliminer les eaux stagnantes autour des habitations (vider les récipients, curer les caniveaux, combler les flaques) ;
\item pulvérisation intra-domiciliaire d'insecticides à effet rémanent (PID) ; utiliser des répulsifs cutanés, porter des vêtements couvrants le soir.
\end{itemize}
\item \textbf{Moyens de lutte curative} : consulter \textbf{rapidement} un centre de santé dès les premiers symptômes ; confirmation du diagnostic par \textbf{test de diagnostic rapide (TDR)} ou \textbf{goutte épaisse} (microscopie) ; traitement par \textbf{thérapie combinée à base d'artémisinine (ACT)} prescrite par le personnel de santé (ex. : artéméther-luméfantrine, artesunate-amodiaquine).
\item \textbf{Structurer la réponse} en deux parties distinctes : prévention / traitement.
\end{enumerate}
\end{methode}

\begin{exoresolu}[Lutter contre le paludisme dans un quartier]
\textbf{Énoncé.} Un élève de 3ème présente une forte fièvre, des frissons et des maux de tête. Le médecin du centre de santé diagnostique un paludisme.
\begin{enumerate}
\item Cite deux autres manifestations possibles du paludisme.
\item Propose trois mesures que la famille de cet élève peut prendre pour éviter de nouveaux cas.
\item Pourquoi est-il dangereux de se soigner soi-même (automédication) ?
\end{enumerate}
\tcblower
\textbf{Solution détaillée.}
\begin{enumerate}
\item Deux autres manifestations : les \textbf{vomissements} et la \textbf{fatigue intense} (on peut aussi citer les douleurs musculaires, l'anémie dans les cas graves).
\item Trois mesures de prévention :
\begin{itemize}
\item faire dormir chaque membre de la famille sous une \textbf{moustiquaire imprégnée à longue durée d'action (MILD)}, bien bordée sous le matelas ;
\item \textbf{éliminer les eaux stagnantes} autour de la maison (vider les seaux, pneus usagés, curer les caniveaux) pour détruire les gîtes larvaires des anophèles ;
\item utiliser des \textbf{insecticides} domestiques (spirales, aérosols) ou des produits répulsifs et porter des vêtements à manches longues le soir.
\end{itemize}
\item L'automédication est dangereuse car une fièvre peut avoir d'autres causes que le paludisme (infections bactériennes, virales, typhoïde, etc.) ; seul un \textbf{diagnostic biologique au centre de santé} (TDR ou goutte épaisse) permet de confirmer la maladie et de prescrire le \textbf{bon médicament à la bonne dose} (ACT). Un traitement mal suivi (mauvais dosage, arrêt précoce, monothérapie) peut aggraver la maladie (évolution vers le paludisme grave, parfois mortel) et sélectionner des souches de \emph{Plasmodium} \textbf{résistantes} aux antipaludiques.
\end{enumerate}
\textbf{Conclusion.} La lutte contre le paludisme repose sur la prévention des piqûres (MILD, destruction des gîtes larvaires) et sur une prise en charge médicale précoce et correcte.
\end{exoresolu}

\courssub{Méthode 4 : TP N°9 — Utiliser correctement une moustiquaire imprégnée}

\begin{methode}[Conduite du TP : moustiquaire imprégnée]
Pour réussir le TP et rédiger le compte rendu :
\begin{enumerate}
\item \textbf{Définir le matériel} : une \cle{moustiquaire imprégnée à longue durée d'action (MILD)} est une moustiquaire dont les fibres du tissu intègrent un \textbf{insecticide} (pyréthrinoïde) qui tue ou repousse les moustiques pendant plusieurs années (norme OMS : 3 ans / 20 lavages).
\item \textbf{Protocole d'utilisation correcte} :
\begin{enumerate}
\item suspendre la moustiquaire au-dessus du lit avant le coucher (idéalement avant la nuit) ;
\item vérifier qu'elle n'est pas déchirée (ou réparer les trous par couture/nœuds) ;
\item bien \textbf{border} les bords inférieurs sous le matelas ou la natte, sans laisser aucune ouverture ;
\item dormir au centre du lit, en évitant que le corps ne touche le tissu (risque de piqûre à travers la maille).
\end{enumerate}
\item \textbf{Entretien} : laver la moustiquaire \textbf{rarement} (max 1 fois tous les 3 mois) et avec de l'eau et un \textbf{savon doux} (pas de détergent, pas de javel, pas de brosse) ; la faire sécher à l'ombre (le soleil dégrade l'insecticide) ; ne pas la réimprégner soi-même (technologie MILD) ; la remplacer après 3 ans ou si elle est trop abîmée.
\item \textbf{Interprétation} : la moustiquaire constitue une \textbf{barrière physique} (maillage fin $< 1,5$ mm) ; l'insecticide ajoute une \textbf{action chimique} (effet « knock-down » et mortalité) qui tue les moustiques au contact : double protection contre le paludisme.
\end{enumerate}
\end{methode}

\begin{exoresolu}[Compte rendu de TP : la moustiquaire imprégnée]
\textbf{Énoncé.} Au cours du TP N°9, deux groupes d'élèves comparent deux dispositifs de protection contre les piqûres de moustiques :
\begin{itemize}
\item Dispositif 1 : une moustiquaire ordinaire (non imprégnée), mal bordée, laissant une ouverture.
\item Dispositif 2 : une moustiquaire imprégnée d'insecticide (MILD), correctement suspendue et bien bordée.
\end{itemize}
\begin{enumerate}
\item Quel dispositif protège le mieux ? Justifie ta réponse.
\item Explique le double rôle de la moustiquaire imprégnée.
\item Pourquoi ne faut-il pas laver la moustiquaire imprégnée trop souvent ni avec des produits agressifs ?
\end{enumerate}
\tcblower
\textbf{Solution détaillée.}
\begin{enumerate}
\item Le \textbf{dispositif 2} protège le mieux. En effet, la moustiquaire ordinaire mal bordée laisse une \textbf{ouverture} par laquelle les moustiques peuvent entrer et piquer le dormeur ; de plus, elle ne tue pas les moustiques qui se posent dessus. La moustiquaire imprégnée (MILD) bien bordée, elle, ne laisse aucun passage physique et détruit les moustiques qui se posent sur le tissu ou tentent de traverser la maille.
\item La moustiquaire imprégnée a un \textbf{double rôle} :
\begin{itemize}
\item un rôle de \textbf{barrière physique} : le tissu à mailles fines empêche mécaniquement le moustique d'atteindre le dormeur ;
\item un rôle \textbf{chimique (insecticide)} : l'insecticide (pyréthrinoïde) incorporé aux fibres \textbf{tue (effet knock-down + mortalité) ou repousse} les moustiques qui entrent en contact avec la moustiquaire.
\end{itemize}
\item Les lavages fréquents ou l'utilisation de détergents/brosse \textbf{éliminent progressivement l'insecticide} fixé sur les fibres : la moustiquaire perd alors son action chimique et redevient une simple barrière physique moins efficace (les moustiques peuvent piquer à travers la maille). Il faut donc la laver rarement (1 fois / trimestre), avec un savon doux, à la main, et la sécher à l'ombre.
\end{enumerate}
\textbf{Conclusion.} Correctement installée (bordée) et entretenue (peu de lavages doux, séchage à l'ombre), la moustiquaire imprégnée à longue durée d'action (MILD) est le moyen de prévention individuelle le plus efficace et le plus rentable contre le paludisme.
\end{exoresolu}

\courssub{Méthode 5 : Présenter les épidémies virales : fièvre Ebola et COVID-19}

\begin{methode}[Décrire une épidémie virale (causes et lutte)]
Pour traiter une question sur Ebola ou la COVID-19 :
\begin{enumerate}
\item \textbf{Identifier la cause} : ces deux maladies sont dues à des \cle{virus} (virus Ebola, genre \emph{Ebolavirus}, famille Filoviridae ; coronavirus \emph{SARS-CoV-2} pour la COVID-19).
\item \textbf{Préciser le mode de transmission} :
\begin{itemize}
\item \textbf{Ebola} : contact \textbf{direct} avec le \textbf{sang ou les liquides biologiques} (selles, urines, salive, sperme, lait maternel, larmes) d'une personne malade symptomatique ou décédée, ou avec des animaux infectés (chauves-souris, primates) ; contamination possible par objets souillés (aiguilles, vêtements). Pas de transmission aérienne.
\item \textbf{COVID-19} : voie \textbf{respiratoire} (inhalation de \textbf{gouttelettes} et \textbf{aérosols} émis en toussant, parlant, chantant, éternuant) ; contact direct avec mains souillées puis visage (bouche, nez, yeux). Transmission possible par personnes asymptomatiques ou pré-symptomatiques.
\end{itemize}
\item \textbf{Citer les manifestations principales} :
\begin{itemize}
\item Ebola : fièvre brutale ($> 38,5\,^\circ\mathrm{C}$), fatigue intense, maux de tête, douleurs musculaires, vomissements, diarrhées, \textbf{hémorragies} (saignements nez, gencives, yeux, selles) en phase tardive ; létalité élevée ($25\,\%$ à $90\,\%$).
\item COVID-19 : fièvre, toux sèche, fatigue, \textbf{perte du goût (agueusie) ou de l'odorat (anosmie)} (signe assez spécifique), difficultés respiratoires (dyspnée) dans les formes graves ; formes asymptomatiques fréquentes.
\end{itemize}
\item \textbf{Citer les moyens de lutte} :
\begin{itemize}
\item Ebola : \textbf{isolement strict} des cas confirmés (centres de traitement), \textbf{recherche des contacts} et surveillance, \textbf{enterrements sécurisés et dignes} (équipes formées), \textbf{lavage des mains} au savon/eau ou solution hydroalcoolique, \textbf{équipements de protection individuelle (EPI)} pour soignants, \textbf{vaccination} (vaccin rVSV-ZEBOV, Ervebo\textregistered) en anneau autour des cas.
\item COVID-19 : \textbf{gestes barrières} (lavage mains fréquent, port du masque chirurgical/FFP2 en espace clos/affluence, distanciation physique $\ge 1$ m, étiquette respiratoire), \textbf{aération} des lieux clos (10 min / heure), \textbf{dépistage} (tests antigéniques/PCR) et isolement des cas positifs, \textbf{vaccination} (vaccins à ARNm, vecteur viral, protéine recombinante) pour prévenir les formes graves.
\end{itemize}
\item \textbf{Point de vigilance absolu} : ne jamais écrire que ces maladies se soignent avec des antibiotiques (les antibiotiques sont inefficaces contre les virus). Pas d'automédication.
\end{enumerate}
\end{methode}

\begin{savaistu}[Riposte au Cameroun]
Le Cameroun a géré plusieurs épidémies d'Ebola (importées depuis la RDC ou la Guinée équatoriale) grâce à la surveillance aux frontières et au laboratoire du Centre Pasteur du Cameroun. Pour la COVID-19, la riposte a été coordonnée par le \emph{Comité Opérationnel de Gestion des Urgences de Santé Publique (COGUSP)} : confinement partiel, port du masque obligatoire, campagne de vaccination (AstraZeneca, Sinopharm, Pfizer, Johnson\&Johnson) ciblant les personnes à risque et le personnel de santé.
\end{savaistu}

\begin{exoresolu}[Ebola et COVID-19 : deux épidémies virales]
\textbf{Énoncé.} Recopie et complète le tableau comparatif suivant, puis réponds à la question.

\begin{center}
\begin{tabularx}{\linewidth}{|l|X|X|}\hline
\rowcolor{popblueL} \textbf{Critère} & \textbf{Fièvre Ebola} & \textbf{COVID-19} \\ \hline
Agent causal & \ldots & \ldots \\ \hline
Mode de transmission principal & \ldots & \ldots \\ \hline
Deux moyens de lutte majeurs & \ldots & \ldots \\ \hline
\end{tabularx}
\end{center}

Question : pourquoi les antibiotiques sont-ils inefficaces contre ces deux maladies ?
\tcblower
\textbf{Solution détaillée.}

Tableau complété :
\begin{center}
\begin{tabularx}{\linewidth}{|l|X|X|}\hline
\rowcolor{popblueL} \textbf{Critère} & \textbf{Fièvre Ebola} & \textbf{COVID-19} \\ \hline
Agent causal & Virus Ebola (\emph{Ebolavirus}) & Coronavirus \emph{SARS-CoV-2} \\ \hline
Mode de transmission principal & Contact direct avec sang/liquides biologiques d'un malade/décédé ou animal infecté & Voie respiratoire : inhalation de gouttelettes/aérosols (toux, parole, éternuements) \\ \hline
Deux moyens de lutte majeurs & Isolement des cas + enterrements sécurisés ; Vaccination (anneau) + EPI soignants & Gestes barrières (masque, distanciation, hygiène mains) + Aération ; Vaccination + Dépistage/isolement \\ \hline
\end{tabularx}
\end{center}

\textbf{Réponse à la question.} Les antibiotiques sont des médicaments qui agissent \textbf{uniquement contre les bactéries} (ils ciblent la paroi cellulaire, la synthèse des protéines ou de l'ADN bactériens). Or Ebola et la COVID-19 sont causées par des \textbf{virus}, qui sont des agents infectieux acellulaires, beaucoup plus petits, dépourvus de métabolisme propre et se répliquant \emph{uniquement} à l'intérieur des cellules hôtes. Les cibles des antibiotiques n'existent pas chez les virus. Les antibiotiques sont donc \textbf{totalement inefficaces} contre ces maladies ; leur emploi abusif est même dangereux car il sélectionne des bactéries \textbf{résistantes} (antibiorésistance). La lutte repose sur la prévention (gestes barrières, isolement, hygiène) et la \textbf{vaccination} (immunité active).

\textbf{Conclusion.} Ebola et la COVID-19 sont deux épidémies virales dont la transmission et les moyens de lutte diffèrent (contact direct vs respiratoire), mais qui partagent un point commun : aucun antibiotique ne peut les soigner.
\end{exoresolu}

\courssub{Méthode 6 : Expliquer le choléra, épidémie hydrique, et synthétiser les moyens de lutte}

\begin{methode}[Traiter une question sur le choléra et comparer les maladies]
\begin{enumerate}
\item \textbf{Cause du choléra} : une \cle{bactérie} à Gram négatif, mobile, le \emph{Vibrio cholerae} (sérogroupe O1 ou O139), appelée \emph{vibrion cholérique}, présente en très grand nombre dans les selles des malades (aspect « eau de riz »).
\item \textbf{Mode de transmission} : voie \textbf{féco-orale}, principalement \textbf{hydrique} (eau souillée) et \textbf{alimentaire} — absorption d'\textbf{eau contaminée} par des selles (puits, rivières, réseau de distribution défaillant) ou d'aliments contaminés (mains sales lors de la préparation, mouches, légumes arrosés aux eaux usées). C'est la maladie type des « mains sales » et de l'insalubrité de l'eau.
\item \textbf{Manifestations} : après une incubation courte (quelques heures à 5 jours), début brutal de \textbf{diarrhées abondantes, indolores, en jets} (selles liquides « eau de riz ») et de \textbf{vomissements} en « jet » sans nausées préalables. Perte massive d'eau et d'électrolytes (Na$^+$, K$^+$, HCO$_3^-$, Cl$^-$) $\rightarrow$ \textbf{déshydratation rapide} (peau plissée, yeux enfoncés, soif intense, oligurie, choc hypovolémique) pouvant être mortelle en quelques heures sans prise en charge. Pas de fièvre habituellement.
\item \textbf{Moyens de lutte} :
\begin{itemize}
\item \textbf{Prévention (hygiène « Eau, Assainissement, Hygiène - EAH »)} : boire de l'\textbf{eau potable} (eau traitée, bouillie 1 min, filtrée, ou chlorée), se \textbf{laver les mains au savon} (cendres si pas de savon) avant les repas, après les toilettes, après avoir soigné un malade ; laver/éplucher/cuire les aliments ; utiliser des \textbf{latrines hygiéniques} (fosse septique, VIP) et gérer les déchets ; lutte anti-vectorielle (mouches).
\item \textbf{Traitement d'urgence} : \textbf{réhydratation} rapide et massive du malade : \textbf{Soluté de Réhydratation Orale (SRO)} (selon formule OMS : NaCl, KCl, Glucose, Citrate/Trisodium citrate) par voie orale (si conscient, pas de vomissements incoercibles) ou \textbf{perfusion intraveineuse} (Ringer Lactate / Saline Normale) en cas de choc/déshydratation sévère ; antibiotique (Doxycycline, Azithromycine, Ciprofloxacine) en appoint pour réduire la durée d'excrétion et le volume des selles (réservé aux cas modérés/sévères) ; supplémentation en Zinc (enfants).
\end{itemize}
\item \textbf{Pour une synthèse de l'unité} : construire un tableau comparatif récapitulatif des quatre maladies étudiées (agent causal, réservoir/vecteur, mode de transmission, manifestations clés, prévention, traitement).
\end{enumerate}
\end{methode}

\begin{exoresolu}[Épidémie de choléra et synthèse de l'unité]
\textbf{Énoncé.} Après de fortes pluies, une épidémie de choléra se déclare dans un quartier où l'eau de boisson provient d'un puits situé près de latrines mal entretenues.
\begin{enumerate}
\item Quel est l'agent causal du choléra ? Comment a-t-il pu contaminer l'eau du puits ?
\item Pourquoi le choléra est-il qualifié de maladie hydrique ?
\item Cite trois mesures à prendre d'urgence pour stopper l'épidémie.
\item Recopie et complète : « Le paludisme est une \ldots transmise par un \ldots, tandis que le choléra est une \ldots transmise par \ldots »
\end{enumerate}
\tcblower
\textbf{Solution détaillée.}
\begin{enumerate}
\item L'agent causal du choléra est une \textbf{bactérie}, le \emph{Vibrio cholerae} (vibrion cholérique). Rejetée en très grand nombre ($10^7$ à $10^9$ par mL) dans les \textbf{selles des malades}, elle a pu, à la faveur des fortes pluies (ruissellement, infiltration), être entraînée depuis les \textbf{latrines mal entretenues} (fosses non étanches, débordement) jusqu'au puits voisin, souillant ainsi l'eau de boisson.
\item Le choléra est une maladie \textbf{hydrique} car sa transmission se fait principalement par l'\textbf{eau souillée} par des matières fécales : on se contamine en buvant de l'eau (ou en mangeant des aliments lavés/préparés avec cette eau) contaminée par le vibrion cholérique.
\item Trois mesures d'urgence :
\begin{itemize}
\item \textbf{Interdire l'utilisation de l'eau du puits contaminé} et fournir de l'eau potable (eau en sachets/bouteilles, eau traitée au chlore, eau bouillie 1 min) à toute la population ;
\item \textbf{Sensibiliser au lavage des mains au savon} (ou cendres) aux moments critiques (avant manger, après toilettes, après soins) et à l'hygiène alimentaire ;
\item \textbf{Prise en charge immédiate des malades} : installation d'un point de réhydratation orale (SRO) au niveau communautaire + référence rapide des cas graves (déshydratation sévère, vomissements incoercibles) vers le centre de santé pour perfusion ; traitement antibiotique des cas modérés/sévères ; désinfection des déjections (chaux, chlore) et des lieux souillés.
\end{itemize}
\item Phrase complétée : « Le paludisme est une \textbf{endémie} transmise par un \textbf{moustique vecteur (anophèle femelle)}, tandis que le choléra est une \textbf{épidémie} transmise par \textbf{l'eau et les aliments souillés (voie féco-orale)} ».
\end{enumerate}
\textbf{Conclusion.} Le choléra, épidémie liée à l'insalubrité de l'eau et à l'absence d'assainissement, se prévient par l'hygiène (eau potable, lavage des mains, latrines) et se traite d'urgence par la réhydratation (SRO/perfusion) ; il illustre, avec le paludisme, Ebola et la COVID-19, la diversité des agents (parasite, virus, bactérie), des réservoirs, des modes de transmission et des stratégies de lutte des maladies transmissibles.
\end{exoresolu}

\begin{aretenir}[Tableau de synthèse : Quatre maladies transmissibles au programme]
\begin{center}
\small
\begin{tabularx}{\linewidth}{|l|>{\centering\arraybackslash}X|>{\centering\arraybackslash}X|>{\centering\arraybackslash}X|>{\centering\arraybackslash}X|}\hline
\rowcolor{popblueL} \textbf{Maladie} & \textbf{Agent causal} & \textbf{Transmission} & \textbf{Prévention clé} & \textbf{Traitement clé} \\ \hline
\textbf{Paludisme} (Endémie) & Parasite \emph{Plasmodium} & Piqûre anophèle ♀ (nuit) & MILD, destruction gîtes larvaires, chimioprévention & ACT (artémisinine) + diagnostic TDR/Goutte épaisse \\ \hline
\textbf{Fièvre Ebola} (Épidémie) & Virus \emph{Ebolavirus} & Contact direct fluides corporels (malade/décédé) & Isolement cas, enterrements sécurisés, EPI, Vaccin (anneau) & Soins de support (réhydratation, hémostase), Antiviraux (expérimentaux) \\ \hline
\textbf{COVID-19} (Pandémie) & Virus \emph{SARS-CoV-2} & Respiratoire (gouttelettes/aérosols) & Gestes barrières, masque, aération, Vaccination & Soins de support (O$_2$, anti-inflammatoires), Antiviraux (Paxlovid\textregistered\ldots) \\ \hline
\textbf{Choléra} (Épidémie) & Bactérie \emph{Vibrio cholerae} & Féco-orale : eau/aliments souillés (mains sales) & Eau potable, Lavage mains savon, Latrines, Cuisson & \textbf{Réhydratation (SRO / Perfusion)} + Antibiotique (appoint) \\ \hline
\end{tabularx}
\end{center}
\end{aretenir}

\begin{attention}[Les 5 erreurs qui coûtent des points au BEPC]
\begin{enumerate}
\item \textbf{Confondre endémie et épidémie.} Une endémie est \emph{permanente} dans une région (paludisme) ; une épidémie apparaît \emph{brutalement} (choléra, Ebola). Toujours justifier le terme choisi par un indice de l'énoncé (durée, nombre de cas, zone).
\item \textbf{Dire que le paludisme est dû à un virus ou à une bactérie.} Faux : l'agent est un \textbf{parasite} (protiste), le \emph{Plasmodium}, transmis par la piqûre de l'\textbf{anophèle femelle} (le moustique est le \emph{vecteur}, pas l'agent causal).
\item \textbf{Prescrire des antibiotiques contre Ebola ou la COVID-19.} Les antibiotiques n'agissent que sur les \textbf{bactéries} ; ils sont inefficaces contre les \textbf{virus}. La prévention repose sur les gestes barrières, l'isolement et la vaccination.
\item \textbf{Oublier le rôle de l'eau et des mains dans le choléra.} Le choléra se transmet par l'\textbf{eau et les aliments souillés} (et les mains sales), pas par l'air ni par les moustiques. Le danger majeur est la \textbf{déshydratation} : le premier soin vital est la \textbf{réhydratation} (SRO ou perfusion).
\item \textbf{Mal décrire l'utilisation de la moustiquaire (TP N°9).} Une moustiquaire (MILD) doit être \textbf{imprégnée d'insecticide}, \textbf{bien bordée} sous le matelas (sans ouverture) et \textbf{peu lavée} (lavages rares, savon doux, séchage à l'ombre). Une moustiquaire déchirée ou mal bordée ne protège pas.
\end{enumerate}
\end{attention}

\begin{exercice}[Entraînement BEPC — Unité : Endémies et Épidémies]
\begin{enumerate}
\item \textbf{QCM.} Une maladie présente en permanence dans une région est une : \\
A. Pandémie \quad B. Épidémie \quad C. Endémie \quad D. Zoonose
\item \textbf{Vrai / Faux / Justifier.}
\begin{itemize}
\item[a)] Le vecteur du paludisme est le \emph{Plasmodium}.
\item[b)] La COVID-19 se transmet par l'eau contaminée.
\item[c)] La réhydratation orale (SRO) est le traitement de première intention du choléra.
\item[d)] Les moustiquaires imprégnées se lavent chaque semaine avec de l'eau de Javel.
\end{itemize}
\item \textbf{Problème de synthèse.} Dans un camp de déplacés au Nord-Cameroun, on signale en une semaine : 50 cas de fièvre avec frissons (confirmés paludisme par TDR), 5 cas de diarrhée aqueuse abondante avec vomissements (choléra suspecté), et 2 cas de fièvre hémorragique (Ebola suspecté).
\begin{enumerate}
\item Pour chaque pathologie, cite l'agent causal (nature et nom) et le mode de transmission principal.
\item Propose trois mesures concrètes et adaptées au contexte de camp pour limiter la propagation de ces trois maladies simultanément.
\item Pourquoi la vaccination est-elle une stratégie majeure pour la COVID-19 et Ebola, mais pas encore disponible de façon routinière pour le paludisme au Cameroun (préciser le statut du vaccin RTS,S/AS01) ?
\end{enumerate}
\end{enumerate}
\end{exercice}

\begin{corrige}[Exercice Entraînement BEPC]
\begin{enumerate}
\item \textbf{C.} Endémie (définition : présence permanente/habituelle).
\item \textbf{Vrai / Faux / Justifier.}
\begin{itemize}
\item[a)] \textbf{Faux.} Le \emph{Plasmodium} est l'\textbf{agent causal} (parasite). Le \textbf{vecteur} est le moustique \textbf{anophèle femelle}.
\item[b)] \textbf{Faux.} La COVID-19 se transmet par voie \textbf{respiratoire} (gouttelettes/aérosols). Le choléra est la maladie hydrique type.
\item[c)] \textbf{Vrai.} La réhydratation orale (SRO) est le traitement de référence, sauveur de vie, pour le choléra (pertes hydro-électrolytiques massives).
\item[d)] \textbf{Faux.} Les MILD se lavent \textbf{rarement} (1 fois / 3 mois max), avec \textbf{savon doux}, séchage \textbf{à l'ombre}. L'eau de Javel et les lavages fréquents détruisent l'insecticide.
\end{itemize}
\item \textbf{Problème de synthèse.}
\begin{enumerate}
\item \textbf{Paludisme :} Agent = Parasite \emph{Plasmodium} (falciparum surtout). Transmission = Piqûre anophèle femelle infectée (nuit). \\
\textbf{Choléra :} Agent = Bactérie \emph{Vibrio cholerae}. Transmission = Féco-orale hydrique/alimentaire (eau/mains/aliments souillés par selles). \\
\textbf{Ebola :} Agent = Virus \emph{Ebolavirus}. Transmission = Contact direct sang/liquides biologiques malade/décédé (ou animal).
\item \textbf{Mesures communes / adaptées au camp :}
\begin{itemize}
\item \textbf{Eau / Hygiène / Assainissement (EAH) :} Fournir eau chlorée/ traitée en quantité suffisante (20 L/pers/jour) ; installer latrines d'urgence (fosses à simple fosse, VIP) éloignées des points d'eau ; organiser collecte/destruction déchets/excrétas ; distribution savon + promotion lavage mains (dispositifs lave-mains aux latrines, cuisines, points d'eau).
\item \textbf{Lutte vectorielle (Paludisme) :} Distribution massive MILD (1 pour 2 personnes) + installation correcte (bordage) + sensibilisation usage toutes les nuites ; destruction gîtes larvaires (assèchement flaques, larvicidage si possible) ; pulvérisation intra-domiciliaire (PID) si épi déco.
\item \textbf{Surveillance / Isolement / Riposte (Choléra / Ebola / COVID) :} Mise en place poste de santé / CTC (Centre de Traitement Choléra) avec stocks SRO/Perfusions/Antibiotiques ; zone d'isolement pour cas suspects Ebola/COVID (EPI pour soignants) ; alerte district/santé publique ; enterrements sécurisés (Ebola) ; vaccination réactive (Choléra : vaccin oral OCV ; COVID-19 : selon dispo ; Ebola : vaccination anneau si cas confirmé).
\end{itemize}
\item \textbf{Vaccination :} COVID-19 et Ebola : vaccins efficaces, sûrs, déployés en routine (COVID) ou en riposte (Ebola) induisant une immunité protectrice forte. Paludisme : le vaccin \textbf{RTS,S/AS01 (Mosquirix\textregistered)} a reçu une recommandation OMS (oct. 2021) pour les enfants en zone de transmission modérée/élevée. Au Cameroun, l'introduction dans le PEV (Programme Élargi de Vaccination) a débuté \textbf{en janvier 2024} (phase pilote dans 42 districts prioritaires). Il n'est pas encore « routinier » sur tout le territoire pour tous les enfants (déploiement progressif), son efficacité est partielle ($\sim 30-40\,\%$ contre paludisme sévère sur 4 ans, nécessite 4 doses), et il \textbf{ne remplace pas} les MILD et la chimioprévention (CPS). C'est un outil \textbf{complémentaire}.
\end{enumerate}
\end{enumerate}
\end{corrige}

\newpage

\coursec{Exercices}

\courssub{Je vérifie mes connaissances}

\begin{exercice}[QCM : une seule bonne réponse]
Pour chaque question, choisis la bonne réponse.
\begin{enumerate}
\item Une maladie qui sévit \textbf{en permanence} dans une région donnée, avec un nombre de cas relativement stable, est :
\begin{enumerate}
\item une épidémie ; \item une endémie ; \item une pandémie ; \item une infection.
\end{enumerate}
\item L'agent causal du paludisme est :
\begin{enumerate}
\item un virus ; \item une bactérie ; \item un protozoaire du genre \emph{Plasmodium} ; \item un champignon.
\end{enumerate}
\item Le paludisme est transmis à l'homme par :
\begin{enumerate}
\item l'eau souillée ; \item la piqûre d'une femelle d'anophèle infectée ; \item les gouttelettes de salive ; \item le contact avec le sang d'un malade.
\end{enumerate}
\item Le choléra se transmet principalement par :
\begin{enumerate}
\item les moustiques ; \item l'eau et les aliments souillés ; \item l'air ; \item les rapports sexuels.
\end{enumerate}
\end{enumerate}
\end{exercice}

\begin{exercice}[Vrai ou faux — justifier]
Réponds par vrai ou faux, puis justifie ta réponse en une ou deux phrases.
\begin{enumerate}
\item La fièvre Ebola est une maladie causée par une bactérie.
\item La COVID-19 se transmet notamment par les gouttelettes respiratoires émises par une personne infectée.
\item Dormir sous une moustiquaire imprégnée d'insecticide est un moyen de lutte contre le choléra.
\item Une épidémie qui s'étend à plusieurs pays ou continents est appelée pandémie.
\end{enumerate}
\end{exercice}

\begin{exercice}[Définitions]
Définis en deux phrases maximum chacun des termes suivants : \cle{endémie}, \cle{épidémie}, \cle{vecteur}, \cle{geste barrière}.
\end{exercice}

\begin{exercice}[Lecture de données]
Dans un centre de santé de la région du Centre, on a enregistré les nombres de cas de paludisme suivants : janvier : 120 cas ; février : 135 cas ; mars : 128 cas.
\begin{enumerate}
\item Calcule le nombre total de cas enregistrés au cours de ce trimestre.
\item Calcule le nombre moyen de cas par mois.
\item Ces valeurs relativement stables d'un mois à l'autre sont-elles plutôt en faveur d'une endémie ou d'une épidémie ? Justifie.
\end{enumerate}
\end{exercice}

\courssub{Je m'entraîne}

\begin{exercice}[Endémie ou épidémie ?]
Le tableau ci-dessous donne le nombre de cas de deux maladies A et B enregistrés dans un district de santé sur trois années.

\begin{center}
\begin{tabular}{|l|c|c|c|}
\hline
\rowcolor{popblueL} Année & 2022 & 2023 & 2024 \\
\hline
Maladie A (nombre de cas) & 480 & 510 & 495 \\
\hline
Maladie B (nombre de cas) & 5 & 12 & 340 \\
\hline
\end{tabular}
\end{center}

\begin{enumerate}
\item Laquelle de ces deux maladies présente un caractère endémique dans ce district ? Justifie ta réponse.
\item Laquelle présente un caractère épidémique en 2024 ? Justifie.
\item En t'appuyant sur tes connaissances, propose une maladie étudiée en classe qui pourrait correspondre à la maladie A, puis une qui pourrait correspondre à la maladie B.
\end{enumerate}
\end{exercice}

\begin{exercice}[Le cycle de transmission du paludisme]
Recopie et complète le texte suivant avec les mots de la liste : \emph{anophèle femelle, Plasmodium, piqûre, sang, foie, globules rouges}.

Le paludisme est causé par un parasite appelé \ldots\ldots\ldots\ldots\ldots. Il est transmis à l'homme lors de la \ldots\ldots\ldots\ldots\ldots d'un moustique, l'\ldots\ldots\ldots\ldots\ldots, préalablement infecté. Chez l'homme, le parasite se multiplie d'abord dans le \ldots\ldots\ldots\ldots\ldots, puis il envahit les \ldots\ldots\ldots\ldots\ldots du \ldots\ldots\ldots\ldots\ldots, qu'il détruit, provoquant les accès de fièvre.
\begin{enumerate}
\item Quel est le rôle de l'anophèle dans la transmission du paludisme ?
\item Pourquoi dit-on que l'anophèle est un \cle{vecteur} et non l'agent causal de la maladie ?
\end{enumerate}
\end{exercice}

\begin{exercice}[Étude de cas : la famille Ngo]
La famille Ngo habite une maison située à proximité d'un marigot, dans un quartier de Yaoundé. Les fenêtres n'ont pas de moustiquaires, la famille dort sans moustiquaire imprégnée, et des flaques d'eau stagnante entourent la maison. Deux enfants de la famille ont été traités pour paludisme au cours des six derniers mois.
\begin{enumerate}
\item Dégage de ce texte trois facteurs qui favorisent la transmission du paludisme dans cette famille.
\item Propose trois mesures de prévention adaptées à cette situation, en justifiant chacune d'elles.
\item Pourquoi les moustiquaires sont-elles dites « imprégnées » ? Quel est l'intérêt de l'insecticide qu'elles contiennent ?
\end{enumerate}
\end{exercice}

\begin{exercice}[Associer maladie, agent et lutte]
Recopie et complète le tableau suivant en associant à chaque maladie : la nature de son agent (virus, bactérie ou protozoaire), son mode de transmission principal et un moyen de lutte adapté.

\begin{center}
\begin{tabularx}{\linewidth}{|l|X|X|X|}
\hline
\rowcolor{popblueL} Maladie & Nature de l'agent & Mode de transmission & Moyen de lutte \\
\hline
Paludisme & & & \\
\hline
Fièvre Ebola & & & \\
\hline
COVID-19 & & & \\
\hline
Choléra & & & \\
\hline
\end{tabularx}
\end{center}
\end{exercice}

\courssub{Je me prépare au BEPC}

\begin{exercice}[Type BEPC : paludisme et choléra dans une région du Littoral]
Le tableau ci-dessous donne le nombre mensuel de cas de paludisme et de choléra enregistrés dans une région du Littoral pendant l'année 2024.

\begin{center}
\begin{tabular}{|l|c|c|c|c|c|c|c|c|c|c|c|c|}
\hline
\rowcolor{popblueL} Mois & J & F & M & A & M & J & J & A & S & O & N & D \\
\hline
Paludisme & 95 & 88 & 102 & 110 & 98 & 90 & 105 & 99 & 92 & 108 & 100 & 96 \\
\hline
Choléra & 0 & 0 & 0 & 0 & 0 & 0 & 0 & 12 & 85 & 140 & 30 & 2 \\
\hline
\end{tabular}
\end{center}

\begin{enumerate}
\item Calcule le nombre total de cas de paludisme et le nombre total de cas de choléra enregistrés dans l'année.
\item Représente sur un même graphique (axe des abscisses : mois ; axe des ordonnées : nombre de cas) l'évolution des deux maladies au cours de l'année.
\item En comparant les deux courbes, identifie la maladie endémique et la maladie épidémique dans cette région. Justifie ta réponse à partir des définitions.
\item Cite l'agent causal de chacune de ces deux maladies et précise son mode de transmission.
\item Les cas de choléra apparaissent en septembre-octobre, en pleine saison des pluies. Propose une explication à cette observation.
\item Dégage deux mesures de lutte adaptées à chacune de ces maladies.
\end{enumerate}
\end{exercice}

\begin{exercice}[Type BEPC : la COVID-19 au Cameroun en 2020]
Le tableau ci-dessous donne une estimation du nombre de nouveaux cas de COVID-19 confirmés par mois au Cameroun au cours de l'année 2020 (données simplifiées).

\begin{center}
\begin{tabular}{|l|c|c|c|c|c|c|c|c|c|}
\hline
\rowcolor{popblueL} Mois & Mars & Avril & Mai & Juin & Juillet & Août & Sept. & Oct. & Nov. \\
\hline
Nouveaux cas & 150 & \num{1500} & \num{3200} & \num{4100} & \num{2500} & \num{1200} & 600 & 400 & 900 \\
\hline
\end{tabular}
\end{center}

\begin{enumerate}
\item Trace la courbe du nombre de nouveaux cas en fonction des mois (échelle conseillée : 1 cm pour 1 mois en abscisse ; 1 cm pour 500 cas en ordonnée).
\item Repère sur la courbe la période correspondant au pic (ou « vague ») épidémique. Précise le mois et le nombre approximatif de cas.
\item La COVID-19 s'est étendue à presque tous les pays du monde en 2020. Comment appelle-t-on une épidémie d'une telle ampleur ? Donne sa définition.
\item Cite l'agent causal de la COVID-19 et son principal mode de transmission.
\item Cite deux gestes barrières recommandés pendant cette épidémie et explique, pour chacun, le mécanisme par lequel il réduit la transmission.
\item Compare le mode de transmission de la COVID-19 à celui du choléra, puis à celui de la fièvre Ebola.
\end{enumerate}
\end{exercice}

\begin{exercice}[Type BEPC : problème de synthèse — un cas de fièvre au retour du village]
Moussa, élève de 3ème, revient d'un séjour de deux semaines dans un village de la région de l'Est, où il dormait sans moustiquaire. Trois jours après son retour, il présente une fièvre élevée (39,5 °C), accompagnée de frissons intenses suivis de sueurs abondantes. Ces accès se répètent environ tous les deux jours. Sa mère l'amène au centre de santé intégré de l'arrondissement.
\begin{enumerate}
\item À partir des signes présentés par Moussa et des circonstances de son séjour, propose une hypothèse diagnostique. Justifie ta réponse par deux arguments tirés de l'énoncé.
\item Nomme l'agent causal de la maladie suspectée et précise comment Moussa a pu être contaminé.
\item Quel examen de laboratoire le personnel de santé peut-il réaliser pour confirmer le diagnostic ? Précise ce que l'on recherche et dans quel échantillon.
\item Le diagnostic est confirmé. Indique la conduite à tenir : le traitement de Moussa d'une part, et deux mesures de prévention à mettre en œuvre pour protéger toute la famille d'autre part.
\item Le caractère cyclique des accès de fièvre (tous les deux jours) est lié à la destruction répétée de certaines cellules du sang par le parasite. Quelles sont ces cellules ?
\item Le paludisme est-il une endémie ou une épidémie dans la plupart des régions du Cameroun ? Justifie ta réponse.
\end{enumerate}
\end{exercice}



\newpage

\coursec{Corrigés des exercices}

\begin{corrige}[Exercice 1 — QCM : une seule bonne réponse]
\begin{enumerate}
\item \textbf{Réponse b} : une maladie présente \textbf{en permanence} dans une région, avec un nombre de cas relativement stable, est une \cle{endémie}. Une épidémie correspond au contraire à une apparition brutale et limitée dans le temps.
\item \textbf{Réponse c} : l'agent causal du paludisme est un \cle{protozoaire} (organisme unicellulaire) du genre \emph{Plasmodium}. Ce n'est ni un virus, ni une bactérie, ni un champignon.
\item \textbf{Réponse b} : le paludisme est transmis par la \textbf{piqûre d'une femelle d'\emph{Anopheles} infectée}. Seule la femelle pique, car elle a besoin de sang pour la maturation de ses œufs ; le mâle se nourrit de sucs végétaux.
\item \textbf{Réponse b} : le choléra est une maladie \textbf{hydrique} : il se transmet par l'eau et les aliments souillés par les selles des malades. Il n'est transmis ni par les moustiques, ni par l'air, ni par les rapports sexuels.
\end{enumerate}
\end{corrige}

\begin{corrige}[Exercice 2 — Vrai ou faux — justifier]
\begin{enumerate}
\item \textbf{Faux.} La fièvre Ebola est causée par un \cle{virus} (le virus Ebola), et non par une bactérie. C'est pourquoi les antibiotiques, qui n'agissent que sur les bactéries, sont inefficaces contre cette maladie.
\item \textbf{Vrai.} Le virus de la COVID-19 (SARS-CoV-2) se transmet surtout par les \textbf{gouttelettes respiratoires} et les sécrétions projetées par une personne infectée lorsqu'elle tousse, éternue ou parle.
\item \textbf{Faux.} La moustiquaire imprégnée protège contre les piqûres d'\emph{Anopheles}, donc contre le \cle{paludisme}. Le choléra se transmet par l'eau et les aliments souillés : la lutte contre le choléra repose sur l'hygiène de l'eau et des aliments, pas sur les moustiquaires.
\item \textbf{Vrai.} Une \cle{pandémie} est une épidémie qui s'étend à une vaste échelle, touchant plusieurs pays ou continents. Exemple : la COVID-19, qui a touché presque tous les pays du monde en 2020.
\end{enumerate}
\end{corrige}

\begin{corrige}[Exercice 3 — Définitions]
\begin{itemize}
\item \textbf{Endémie} : maladie qui sévit en permanence dans une région donnée, avec un nombre de cas relativement stable au cours du temps. Exemple : le paludisme dans la plupart des régions du Cameroun.
\item \textbf{Épidémie} : apparition brutale et propagation rapide d'une maladie touchant un grand nombre de personnes dans une région, pendant une période limitée. Exemple : une épidémie de choléra.
\item \textbf{Vecteur} : organisme (souvent un insecte) qui transporte l'agent causal d'une maladie d'un individu malade vers un individu sain, sans être lui-même la cause de la maladie. Exemple : la femelle d'\emph{Anopheles}, vecteur du paludisme.
\item \textbf{Geste barrière} : mesure simple d'hygiène qui empêche ou réduit la transmission d'un agent pathogène. Exemples : lavage régulier des mains, port du masque, distanciation physique.
\end{itemize}
\end{corrige}

\begin{corrige}[Exercice 4 — Lecture de données]
\begin{enumerate}
\item Nombre total de cas sur le trimestre :
\[120 + 135 + 128 = 383 \text{ cas.}\]
\item Nombre moyen de cas par mois :
\[\frac{383}{3} \approx \num{127,7} \text{, soit environ } \num{128} \text{ cas par mois.}\]
\item Ces valeurs sont en faveur d'une \cle{endémie} : le nombre de cas reste élevé mais relativement stable d'un mois à l'autre (entre 120 et 135 cas, soit un écart maximal de 15 cas), sans augmentation brutale. Une épidémie se traduirait au contraire par une hausse rapide et importante du nombre de cas sur une courte période.
\end{enumerate}
\end{corrige}

\begin{corrige}[Exercice 5 — Endémie ou épidémie ?]
\begin{enumerate}
\item La maladie A présente un caractère \textbf{endémique} : le nombre de cas est élevé mais reste stable d'une année à l'autre (480, 510 puis 495 cas, soit un écart maximal de 30 cas autour de 500). Cela correspond à une présence permanente de la maladie dans le district.
\item La maladie B présente un caractère \textbf{épidémique} en 2024 : le nombre de cas passe brutalement de 12 en 2023 à 340 en 2024, soit une multiplication par environ \num{28,3} en un an ($340 \div 12 \approx 28,3$). Cette flambée soudaine et importante est la signature d'une épidémie.
\item La maladie A pourrait être le \textbf{paludisme}, endémique dans la plupart des régions du Cameroun. La maladie B pourrait être le \textbf{choléra}, qui sévit par épidémies, souvent pendant la saison des pluies.
\end{enumerate}
\textbf{Point de vigilance} : la justification doit toujours s'appuyer sur les \emph{données chiffrées} du tableau (stabilité ou variation brutale des nombres de cas), et non seulement sur la définition récitée.
\end{corrige}

\begin{corrige}[Exercice 6 — Le cycle de transmission du paludisme]
\textbf{Texte complété} : le paludisme est causé par un parasite appelé \emph{Plasmodium}. Il est transmis à l'homme lors de la \textbf{piqûre} d'un moustique, l'\textbf{\emph{Anopheles} femelle}, préalablement infecté. Chez l'homme, le parasite se multiplie d'abord dans le \textbf{foie}, puis il envahit les \textbf{globules rouges} du \textbf{sang}, qu'il détruit, provoquant les accès de fièvre.
\begin{enumerate}
\item L'\emph{Anopheles} femelle transporte le parasite : en piquant une personne infectée, elle prélève le \emph{Plasmodium} avec le sang ; en piquant ensuite une personne saine, elle le lui inocule avec sa salive. Elle assure ainsi la transmission du parasite d'un homme à l'autre.
\item L'\emph{Anopheles} est un \cle{vecteur} car il ne fait que \textbf{transporter} le parasite : ce n'est pas lui qui cause la maladie. L'agent causal est le \emph{Plasmodium}, qui se multiplie dans l'organisme humain et détruit les globules rouges.
\end{enumerate}
\end{corrige}

\begin{corrige}[Exercice 7 — Étude de cas : la famille Ngo]
\begin{enumerate}
\item Trois facteurs qui favorisent la transmission du paludisme :
\begin{itemize}
\item la proximité du marigot et les flaques d'eau stagnante, qui constituent des \textbf{gîtes larvaires} où les \emph{Anopheles} pondent et où leurs larves se développent ;
\item l'absence de moustiquaires aux fenêtres, qui laisse les moustiques entrer librement dans la maison ;
\item le fait de dormir sans moustiquaire imprégnée, ce qui expose les dormeurs aux piqûres nocturnes des \emph{Anopheles}.
\end{itemize}
\item Trois mesures de prévention adaptées :
\begin{itemize}
\item \textbf{assécher ou assainir les eaux stagnantes} autour de la maison (curage, remblaiement) : cela détruit les gîtes larvaires et réduit le nombre d'\emph{Anopheles} ;
\item \textbf{dormir sous des moustiquaires imprégnées} d'insecticide : elles empêchent les piqûres pendant le sommeil et tuent les moustiques qui s'y posent ;
\item \textbf{poser des moustiquaires aux fenêtres} : elles empêchent l'entrée des \emph{Anopheles} dans la maison.
\end{itemize}
\item Les moustiquaires sont dites « imprégnées » parce que leur tissu est \textbf{trempé dans un insecticide}. L'insecticide tue ou repousse les moustiques qui entrent en contact avec la moustiquaire : celle-ci ne se contente donc pas de constituer une barrière physique, elle réduit aussi le nombre d'\emph{Anopheles}, ce qui protège toute la famille, même si la moustiquaire présente un petit trou.
\end{enumerate}
\end{corrige}

\begin{corrige}[Exercice 8 — Associer maladie, agent et lutte]
\begin{center}
\begin{tabularx}{\linewidth}{|l|X|X|X|}
\hline
\rowcolor{popblueL} Maladie & Nature de l'agent & Mode de transmission & Moyen de lutte \\
\hline
Paludisme & Protozoaire (\emph{Plasmodium}) & Piqûre d'\emph{Anopheles} femelle infectée & Moustiquaires imprégnées, destruction des gîtes larvaires \\
\hline
Fièvre Ebola & Virus & Contact direct avec le sang et les liquides biologiques d'un malade & Isolement des malades, gants et tenues de protection, enterrements sécurisés \\
\hline
COVID-19 & Virus (coronavirus) & Gouttelettes respiratoires (toux, éternuements, parole) & Gestes barrières : masque, lavage des mains, distanciation \\
\hline
Choléra & Bactérie (vibrion cholérique) & Eau et aliments souillés par les selles des malades & Hygiène : eau potable, lavage des mains, aliments cuits et protégés \\
\hline
\end{tabularx}
\end{center}
\textbf{Point de vigilance} : ne pas confondre la \emph{nature de l'agent} (virus, bactérie, protozoaire) avec le \emph{mode de transmission} (piqûre, eau souillée, gouttelettes, contact) : ce sont deux colonnes différentes.
\end{corrige}

\begin{corrige}[Exercice 9 — Type BEPC : paludisme et choléra dans une région du Littoral]
\begin{enumerate}
\item \textbf{Paludisme} :
\[95+88+102+110+98+90+105+99+92+108+100+96 = \num{1183} \text{ cas.}\]
\textbf{Choléra} :
\[12+85+140+30+2 = 269 \text{ cas.}\]
\item \textbf{Graphique attendu} : en abscisse, les douze mois de l'année ; en ordonnée, le nombre de cas (échelle possible : 1 cm pour 20 cas). La courbe du paludisme reste presque horizontale, autour de 100 cas par mois ; la courbe du choléra est nulle de janvier à juillet, puis forme un pic net en septembre-octobre (140 cas en octobre) avant de redescendre en novembre-décembre.
\item Le \textbf{paludisme} est la maladie \textbf{endémique} : il sévit en permanence, tous les mois de l'année, avec un nombre de cas élevé et stable (entre 88 et 110 cas). Le \textbf{choléra} est la maladie \textbf{épidémique} : absent pendant sept mois, il apparaît brutalement et touche un grand nombre de personnes en peu de temps (269 cas en cinq mois, dont 140 en octobre seul).
\item \textbf{Paludisme} : agent causal = un protozoaire, le \emph{Plasmodium} ; transmission par la piqûre d'une femelle d'\emph{Anopheles} infectée. \textbf{Choléra} : agent causal = une bactérie, le vibrion cholérique ; transmission par l'eau et les aliments souillés par les selles des malades.
\item En saison des pluies, les eaux de ruissellement et les inondations transportent les matières fécales vers les puits, les marigots et les points d'eau utilisés par les populations. L'eau de boisson souillée par le vibrion cholérique devient alors un véhicule de la maladie, ce qui explique la flambée des cas en septembre-octobre.
\item \textbf{Contre le paludisme} : dormir sous moustiquaire imprégnée ; détruire les eaux stagnantes (gîtes larvaires) et assainir l'environnement. \textbf{Contre le choléra} : ne boire que de l'eau potable (bouillie ou traitée) ; se laver les mains à l'eau et au savon avant les repas et après être allé aux toilettes, et bien cuire les aliments.
\end{enumerate}
\textbf{Barème indicatif (sur 10)} : calculs des totaux (1 pt) ; graphique correct avec axes légendés et échelle (2 pts) ; identification endémie/épidémie justifiée par les données (2 pts) ; agents et modes de transmission (2 pts) ; explication du pic saisonnier (1,5 pt) ; mesures de lutte (1,5 pt).\\
\textbf{Points de vigilance} : la justification de la question 3 doit citer les \emph{définitions} et les \emph{valeurs du tableau} ; les courbes doivent être tracées avec des axes gradués et légendés (grandeur et unité en ordonnée).
\end{corrige}

\begin{corrige}[Exercice 10 — Type BEPC : la COVID-19 au Cameroun en 2020]
\begin{enumerate}
\item \textbf{Courbe attendue} : on place les points (mars ; 150), (avril ; \num{1500}), (mai ; \num{3200}), (juin ; \num{4100}), (juillet ; \num{2500}), (août ; \num{1200}), (septembre ; 600), (octobre ; 400), (novembre ; 900), puis on les relie par des segments. La courbe monte fortement de mars à juin, redescend jusqu'en octobre, puis remonte légèrement en novembre.
\item Le \textbf{pic épidémique} (ou « vague ») se situe en \textbf{juin 2020}, avec environ \num{4100} nouveaux cas dans le mois.
\item Une épidémie qui s'étend à presque tous les pays du monde est une \cle{pandémie} : c'est une épidémie de très grande ampleur, touchant simultanément plusieurs pays ou continents.
\item L'agent causal de la COVID-19 est un \textbf{virus} (un coronavirus, SARS-CoV-2). Il se transmet principalement par les \textbf{gouttelettes respiratoires} projetées par une personne infectée lorsqu'elle tousse, éternue ou parle, ainsi que par les mains portées au visage après contact avec des surfaces souillées.
\item Deux gestes barrières et leur mécanisme :
\begin{itemize}
\item le \textbf{port du masque} : il retient les gouttelettes émises par la personne infectée et protège les personnes proches ;
\item le \textbf{lavage régulier des mains} à l'eau et au savon : il élimine le virus déposé sur les mains avant qu'il ne soit porté à la bouche, au nez ou aux yeux.
\end{itemize}
(On peut aussi citer la distanciation physique d'au moins un mètre, qui réduit le risque de recevoir les gouttelettes d'une personne infectée.)
\item La COVID-19 se transmet par \textbf{voie aérienne} (gouttelettes respiratoires) ; le choléra se transmet par \textbf{voie digestive}, par l'eau et les aliments souillés ; la fièvre Ebola se transmet par \textbf{contact direct} avec le sang et les liquides biologiques d'un malade. Les trois maladies exigent donc des mesures de lutte différentes : masque et distanciation pour la COVID-19, hygiène de l'eau et des aliments pour le choléra, isolement des malades et équipements de protection pour Ebola.
\end{enumerate}
\textbf{Barème indicatif (sur 10)} : courbe correcte avec axes légendés et échelle respectée (2 pts) ; repérage du pic (1 pt) ; définition de la pandémie (1 pt) ; agent et mode de transmission (2 pts) ; deux gestes barrières avec mécanisme expliqué (2 pts) ; comparaison des trois modes de transmission (2 pts).\\
\textbf{Points de vigilance} : pour la question 5, il ne suffit pas de citer le geste barrière, il faut \emph{expliquer comment il réduit la transmission} ; pour la question 6, chaque maladie doit être associée à son mode de transmission propre, sans confusion.
\end{corrige}

\begin{corrige}[Exercice 11 — Type BEPC : problème de synthèse — un cas de fièvre au retour du village]
\begin{enumerate}
\item \textbf{Hypothèse diagnostique} : Moussa est atteint de \textbf{paludisme}. Deux arguments tirés de l'énoncé : (a) il a séjourné deux semaines dans une zone où sévit le paludisme en dormant \textbf{sans moustiquaire}, donc exposé aux piqûres d'\emph{Anopheles} ; (b) il présente les signes typiques de la maladie : \textbf{fièvre élevée (\SI{39,5}{\celsius}) avec frissons intenses puis sueurs abondantes}, en accès qui se répètent régulièrement.
\item L'agent causal est un \textbf{protozoaire}, le \emph{Plasmodium}. Moussa a pu être contaminé par la \textbf{piqûre d'une femelle d'\emph{Anopheles} infectée} pendant son sommeil au village : le moustique lui a inoculé le parasite avec sa salive.
\item Le personnel de santé peut réaliser une \textbf{goutte épaisse} (examen d'une goutte de sang au microscope) ou un test de diagnostic rapide (TDR) : on recherche le \emph{Plasmodium} dans un échantillon de \textbf{sang} prélevé au bout du doigt.
\item \textbf{Conduite à tenir} :
\begin{itemize}
\item traitement : administrer à Moussa un \textbf{médicament antipaludique} prescrit par le personnel de santé, en suivant le traitement jusqu'au bout ;
\item prévention pour la famille : (a) faire dormir tous les membres sous des \textbf{moustiquaires imprégnées} d'insecticide ; (b) éliminer les eaux stagnantes autour de la maison (gîtes larvaires) et, si possible, poser des moustiquaires aux fenêtres.
\end{itemize}
\item Ces cellules sont les \textbf{globules rouges} (hématies) : le parasite s'y multiplie puis les détruit en les faisant éclater ; la libération simultanée des parasites déclenche chaque accès de fièvre, d'où le caractère cyclique (tous les deux jours pour les formes tertiennes).
\item Le paludisme est une \textbf{endémie} dans la plupart des régions du Cameroun : il y sévit en permanence, toute l'année, avec un nombre de cas élevé mais relativement stable, contrairement à une épidémie qui apparaît brutalement sur une période limitée.
\end{enumerate}
\textbf{Barème indicatif (sur 10)} : hypothèse justifiée par deux arguments de l'énoncé (2 pts) ; agent causal et contamination (2 pts) ; examen de laboratoire (échantillon + élément recherché) (1,5 pt) ; traitement et deux mesures de prévention (2 pts) ; globules rouges (1 pt) ; caractère endémique justifié (1,5 pt).\\
\textbf{Points de vigilance} : à la question 1, les arguments doivent être \emph{tirés de l'énoncé} (séjour sans moustiquaire, accès de fièvre cycliques) et non généraux ; à la question 3, préciser à la fois l'échantillon (sang) et ce que l'on recherche (le \emph{Plasmodium}) ; à la question 6, la justification doit opposer endémie (permanence, stabilité) et épidémie (brutalité, durée limitée).
\end{corrige}

\newpage

\coursec{Fiche bilan}

\begin{popfigure}
\begin{tikzpicture}[mindmap, grow cyclic, every node/.style={concept, font=\small}, text width=2.6cm, root concept/.append style={concept color=popblue, font=\bfseries, text width=3cm}, level 1/.append style={level distance=4.2cm, sibling angle=60}, level 1 concept/.append style={font=\footnotesize}]
\node[root concept]{Endémies et épidémies}
  child[concept color=poporange]{ node {Définitions : endémie, épidémie, pandémie}
  }
  child[concept color=popgreen]{ node {Paludisme : \emph{Plasmodium}, moustique \emph{Anophèle}}
  }
  child[concept color=popgold]{ node {Paludisme : fièvres, anémie ; lutte (moustiquaires)}
  }
  child[concept color=poppurple]{ node {TP : moustiquaires imprégnées d'insecticide}
  }
  child[concept color=poppink]{ node {Fièvre Ebola et COVID-19 : virus, lutte}
  }
  child[concept color=popmuted]{ node {Choléra : bacille, eau souillée, lutte}
  };
\end{tikzpicture}
\legende{Carte mentale de la leçon : les six idées clés à retenir sur les endémies et les épidémies.}
\end{popfigure}

\begin{bilan}
\textbf{Définitions clés}
\begin{itemize}
  \item Une \cle{endémie} est une maladie \textbf{présente en permanence} dans une région donnée, où elle sévit à un niveau habituel (ex. : le paludisme au Cameroun).
  \item Une \cle{épidémie} est l'\textbf{apparition brutale} et l'extension rapide d'une maladie à un grand nombre de personnes dans une région (ex. : choléra, Ebola, COVID-19).
  \item Une \cle{pandémie} est une épidémie qui s'étend à \textbf{plusieurs pays ou continents} (ex. : COVID-19).
\end{itemize}

\textbf{Les maladies à connaître par c\oe ur}
\begin{center}
\begin{tabularx}{\linewidth}{|l|X|X|X|}
\hline
\rowcolor{popblueL} \textbf{Maladie} & \textbf{Agent causal} & \textbf{Mode de transmission} & \textbf{Moyens de lutte} \\
\hline
Paludisme (endémie) & \emph{Plasmodium} (protozoaire) & Piqûre du moustique femelle \emph{Anophèle} & Moustiquaires imprégnées, destruction des gîtes larvaires, médicaments antipaludiques \\
\hline
Fièvre Ebola (épidémie) & Virus Ebola & Contact avec les liquides biologiques d'un malade & Isolement des malades, hygiène, enterrements sécurisés, vaccination \\
\hline
COVID-19 (pandémie) & Virus SARS-CoV-2 & Gouttelettes respiratoires (toux, parole) & Lavage des mains, masque, distanciation, vaccination \\
\hline
Choléra (épidémie) & Bacille \emph{Vibrio cholerae} & Eau et aliments souillés par les selles & Eau potable, hygiène, latrines, réhydratation des malades \\
\hline
\end{tabularx}
\end{center}

\textbf{Manifestations du paludisme}
\begin{itemize}
  \item Accès palustre : \textbf{fièvre élevée}, frissons, maux de tête, vomissements, fatigue.
  \item En l'absence de soins : anémie, formes graves pouvant entraîner la mort (surtout chez les jeunes enfants et les femmes enceintes).
\end{itemize}

\textbf{Méthode express (TP N°9 : moustiquaire imprégnée)}
\begin{enumerate}
  \item La moustiquaire forme une \textbf{barrière physique} contre les moustiques.
  \item L'\textbf{insecticide} imprégné tue ou repousse les moustiques qui la touchent.
  \item Elle doit être bien installée (bordée sous le matelas), sans trou, et réimprégnée régulièrement.
  \item Dormir sous moustiquaire imprégnée chaque nuit protège surtout les enfants et les femmes enceintes.
\end{enumerate}

\textbf{Méthode express (rédiger une réponse)}
\begin{enumerate}
  \item Identifier la maladie et son agent causal.
  \item Expliquer le mode de transmission.
  \item Proposer des moyens de lutte adaptés à ce mode de transmission.
  \item Conclure en justifiant le choix des mesures.
\end{enumerate}
\end{bilan}

\begin{aretenir}[Checklist avant le Bac]
\begin{itemize}
  \item[$\square$] Je sais définir une endémie et une épidémie et les distinguer.
  \item[$\square$] Je sais que le paludisme est une endémie au Cameroun.
  \item[$\square$] Je connais l'agent causal du paludisme (\emph{Plasmodium}) et son vecteur (la femelle \emph{Anophèle}).
  \item[$\square$] Je sais citer les manifestations du paludisme (fièvre, frissons, anémie).
  \item[$\square$] Je connais les moyens de lutte contre le paludisme (moustiquaires imprégnées, assainissement, traitement).
  \item[$\square$] Je sais expliquer le rôle protecteur de la moustiquaire imprégnée (barrière + insecticide).
  \item[$\square$] Je sais que la fièvre Ebola et la COVID-19 sont des épidémies causées par des virus.
  \item[$\square$] Je connais les modes de transmission d'Ebola (liquides biologiques) et de la COVID-19 (gouttelettes respiratoires).
  \item[$\square$] Je sais que le choléra se transmet par l'eau et les aliments souillés.
  \item[$\square$] Je sais citer les moyens de lutte contre le choléra (eau potable, hygiène, réhydratation).
  \item[$\square$] Je sais rédiger et justifier une réponse complète (agent, transmission, lutte).
  \item[$\square$] Je sais appliquer ces notions à une situation concrète de la vie courante.
\end{itemize}
\end{aretenir}

\findecours

\end{document}
